\documentclass[11pt,a4paper]{article}
\newcommand{\corso}{Fisica dello Stato Solido}
\newcommand{\argomento}{Parte 4: bande reali}
\newcommand{\lezioni}{Lezioni 5d, 5e e 6a}
% ------------------------------------------------------------------
%  Preambolo delle dispense (skill dispense-latex, layout del 22/10/2026)
%  Modello: frontespizio + indice a pagina intera, etichette colorate
%  numerate per sezione, dimostrazioni a passi con barra laterale,
%  riquadro verde solo per i risultati essenziali.
%  Prima di % ------------------------------------------------------------------
%  Preambolo delle dispense (skill dispense-latex, layout del 22/10/2026)
%  Modello: frontespizio + indice a pagina intera, etichette colorate
%  numerate per sezione, dimostrazioni a passi con barra laterale,
%  riquadro verde solo per i risultati essenziali.
%  Prima di % ------------------------------------------------------------------
%  Preambolo delle dispense (skill dispense-latex, layout del 22/10/2026)
%  Modello: frontespizio + indice a pagina intera, etichette colorate
%  numerate per sezione, dimostrazioni a passi con barra laterale,
%  riquadro verde solo per i risultati essenziali.
%  Prima di \input{preambolo} definisci \corso, \argomento, \lezioni.
%  Compilare con pdflatex (scripts/setup_pdflatex_it.sh per la sillabazione).
% ------------------------------------------------------------------
\providecommand{\corso}{Nome del corso}
\providecommand{\argomento}{Argomento}
\providecommand{\lezioni}{Lezioni}
\usepackage[T1]{fontenc}
\usepackage{babel}
\babelprovide[import, main]{italian}
\usepackage{lmodern}            % Latin Modern (se manca: apt-get install -y lmodern)
\usepackage{microtype}
\usepackage[a4paper, left=2.4cm, right=2.4cm, top=2.5cm, bottom=2.7cm]{geometry}
\usepackage[skip=6pt plus 2pt, indent=0pt]{parskip}

\usepackage{amsmath, amssymb}
\usepackage{bm}
\numberwithin{equation}{section}
\usepackage{siunitx}
\sisetup{output-decimal-marker={,}, group-separator={\,}, exponent-product=\cdot}
\DeclareSIUnit{\angstrom}{\text{\AA}}
\usepackage{booktabs, array, tabularx, multirow, longtable}
\usepackage{enumitem}
\setlist{itemsep=3pt, topsep=4pt, parsep=0pt}
\usepackage[font=small, labelfont=bf, format=hang, skip=6pt]{caption}
\usepackage{float}
\usepackage{needspace}
% \spazio{n\baselineskip}: va a capo pagina solo se restano meno di n righe.
% Sostituisce \needspace, il cui \penalty-100 con \vskip elastico fa spezzare
% le pagine prima dei titoli anche quando c'è posto (pagine mezze vuote).
\makeatletter
\newcommand{\spazio}[1]{\par\begingroup\setlength{\dimen@}{#1}%
  \dimen@ii\pagegoal\advance\dimen@ii-\pagetotal
  \ifdim\dimen@>\dimen@ii\ifdim\dimen@ii>\z@\vfil\penalty-10000\fi\fi\endgroup}
\makeatother

% ---------------------------------------------------------- colori
\usepackage{xcolor}
\definecolor{csez}{HTML}{193C82}   % titoli di sezione
\definecolor{csub}{HTML}{3C5995}   % sottosezioni, link
\definecolor{cdef}{HTML}{1E5AAA}   % Definizione
\definecolor{cris}{HTML}{14823C}   % Risultato, riquadro verde
\definecolor{cbox}{HTML}{F1F7F3}   % fondo del riquadro verde
\definecolor{cform}{HTML}{E4F0E8}  % fondo di \formula
\definecolor{cformb}{HTML}{ACD3BA} % bordo di \formula
\definecolor{coss}{HTML}{BE1E1E}   % Osservazione (refusi, convenzioni)
\definecolor{cint}{HTML}{007D87}   % Interpretazione fisica, «Cosa dice»
\definecolor{cappr}{HTML}{783296}  % Approfondimento
\definecolor{crie}{HTML}{C86E00}   % Riepilogo
\definecolor{cdim}{HTML}{6E6E6E}   % Dimostrazione
\definecolor{cpasso}{HTML}{585858} % Passo n.
\definecolor{cesm}{HTML}{4A6FA5}   % Esempio
\definecolor{cenb}{HTML}{EEF3FA}   % fondo del riquadro blu dell'enunciato
% colori per le figure
\definecolor{accent}{HTML}{193C82}
\definecolor{vA}{HTML}{C0392B}
\definecolor{vB}{HTML}{2E6FB7}
\definecolor{vC}{HTML}{23884A}
\definecolor{atomA}{HTML}{E15759}
\definecolor{atomB}{HTML}{4E79A7}
\definecolor{atomC}{HTML}{F28E2B}
\definecolor{atomD}{HTML}{59A14F}
\definecolor{atomE}{HTML}{B07AA1}
\definecolor{cellP}{HTML}{3C9D4E}
\definecolor{cellC}{HTML}{1F4E8C}
\definecolor{cellS}{HTML}{7B3FA0}
\definecolor{bad}{HTML}{C0392B}
\definecolor{Na}{HTML}{F28E2B}
\definecolor{Cl}{HTML}{4E79A7}
% nomi del preambolo precedente (figure e snippet esistenti)
\colorlet{boxdef}{cdef}\colorlet{boxkey}{cris}\colorlet{boxwarn}{coss}
\colorlet{boxex}{black!65}\colorlet{boxdeep}{cappr}\colorlet{mapgray}{black!60}

% ---------------------------------------------------------- figure
\usepackage{tikz}
\usetikzlibrary{arrows.meta, calc, decorations.pathreplacing, patterns}
\usepackage{pgfplots}
\pgfplotsset{compat=1.18, every axis/.append style={font=\footnotesize}, trig format plots=rad}
\usepackage[europeanresistors, cuteinductors, europeanvoltages, straightvoltages]{circuitikz}
\ctikzset{bipoles/length=1.0cm}
\usepackage{tikz-3dplot}
\tikzset{
  vec/.style={-{Stealth[length=2.4mm, width=1.7mm]}, thick},
  bond/.style={black!22, line width=0.5pt},
  latline/.style={black!30, thin},
  edge/.style={black!75, line width=0.7pt},
  edgehid/.style={black!55, line width=0.6pt, densely dashed},
  edgelight/.style={black!50, line width=0.6pt},
  edgelighthid/.style={black!40, line width=0.5pt, densely dashed},
  cellprim/.style={fill=cellP!16, draw=cellP!75!black, thick},
  cellconv/.style={fill=cellC!10, draw=cellC, thick},
  cellbad/.style={fill=bad!14, draw=bad, thick},
  cellconvline/.style={draw=cellC, thick, densely dashed},
  cellshift/.style={draw=cellS, thick, dashed},
  tri/.style={fill=vB!14, draw=vB!55, thin},
  tri0/.style={fill=vB!45, draw=vB!80!black, thick},
}

% ---------------------------------------------------------- titoli e indice
\usepackage{titlesec}
\titleformat{\section}{\spazio{9\baselineskip}\Large\bfseries\color{csez}}{\thesection}{0.8em}{}%
  [\vspace{-0.6ex}{\color{csez}\titlerule[0.7pt]}]
\titleformat{\subsection}{\spazio{8\baselineskip}\large\bfseries\color{csub}}{\thesubsection}{0.8em}{}
\titleformat{\subsubsection}{\spazio{4\baselineskip}\normalsize\bfseries\color{csub}}{}{0pt}{}
\titlespacing*{\section}{0pt}{22pt plus 4pt minus 2pt}{10pt plus 2pt}
\titlespacing*{\subsection}{0pt}{16pt plus 3pt minus 2pt}{6pt plus 2pt}
\titlespacing*{\subsubsection}{0pt}{10pt plus 2pt}{3pt}
\setcounter{tocdepth}{2}
\usepackage[titles]{tocloft}          % titolo «Indice» con lo stile di \section*
\renewcommand{\cftsecfont}{\bfseries}
\renewcommand{\cftsecpagefont}{\bfseries}
\setlength{\cftbeforesecskip}{11pt}
\setlength{\cftbeforesubsecskip}{3pt}

\pagestyle{plain}
\usepackage[colorlinks, linkcolor=csub, citecolor=csub, urlcolor=csub]{hyperref}
\hypersetup{pdftitle={\corso{} - \argomento}, pdfauthor={}}

% ---------------------------------------------------------- etichette numerate
%  Definizione, Risultato, Osservazione (e i riquadri Teorema, Modello) condividono il
%  contatore: Definizione 1.2, Osservazione 1.3, ...  (numerazione per sezione)
%  NB: \if...\else...\fi solo DENTRO le macro (\tit, \sezvoce): un \else scritto
%  dentro \textbf{...} viene accoppiato con l'\ifmmode di \textbf («Extra \else»).
\newcounter{enun}[section]
\renewcommand{\theenun}{\thesection.\arabic{enun}}
\newcommand{\tit}[1]{\if\relax\detokenize{#1}\relax\else\ (#1)\fi}   % « (titolo)» solo se c'è
\newcommand{\etinum}[3]{% colore, nome, titolo facoltativo
  \refstepcounter{enun}\textbf{\color{#1}#2 \theenun\tit{#3}.}\ \ignorespaces}
\newcommand{\etisemp}[3]{% colore, nome, titolo facoltativo (senza numero)
  \textbf{\color{#1}#2\tit{#3}.}\ \ignorespaces}
\newenvironment{definizione}[1][]{\par\addvspace{4pt}\spazio{3\baselineskip}\noindent\etinum{cdef}{Definizione}{#1}}{\par\addvspace{4pt}}
\newenvironment{risultato}[1][]{\par\addvspace{4pt}\spazio{3\baselineskip}\noindent\etinum{cris}{Risultato}{#1}}{\par\addvspace{4pt}}
\newenvironment{osservazione}[1][]{\par\addvspace{4pt}\spazio{3\baselineskip}\noindent\etinum{coss}{Osservazione}{#1}}{\par\addvspace{4pt}}
\newenvironment{interpretazione}[1][]{\par\addvspace{4pt}\spazio{3\baselineskip}\noindent\etisemp{cint}{Interpretazione fisica}{#1}}{\par\addvspace{4pt}}
\newenvironment{riepilogo}[1][]{\par\addvspace{4pt}\spazio{3\baselineskip}\noindent\etisemp{crie}{Riepilogo}{#1}}{\par\addvspace{4pt}}
% introduzione: titolo e testo in nero, spaziatura stretta
\newenvironment{introduzione}{\begingroup\colorlet{csez}{black}\colorlet{crie}{black}%
  \section*{Introduzione}\addcontentsline{toc}{section}{Introduzione}%
  \setlength{\parskip}{3pt plus 1pt}\setlist{itemsep=1pt, topsep=2pt, parsep=0pt}%
  \setlength{\abovedisplayskip}{4pt plus 2pt}\setlength{\belowdisplayskip}{4pt plus 2pt}%
  \setlength{\abovedisplayshortskip}{2pt}\setlength{\belowdisplayshortskip}{3pt}}%
  {\par\endgroup}
% compatibilità con le dispense precedenti
\newenvironment{attenzione}[1][]{\begin{osservazione}[#1]}{\end{osservazione}}

% ---------------------------------------------------------- dimostrazioni a passi
\usepackage[skins, breakable]{tcolorbox}
\newcounter{passo}
\makeatletter
\newcommand{\passo}[1]{\par\addvspace{5pt}\stepcounter{passo}%
  \noindent{\bfseries\color{cpasso}Passo \thepasso.}\ \textit{#1}\@ifnextchar\[{\nopagebreak}{\@ifnextchar\begin{\nopagebreak}{\par\nopagebreak[4]}}}
\makeatother
\newcommand{\qed}{\ifmmode\qquad\blacksquare\else\unskip\nobreak\ \ensuremath{\blacksquare}\fi}
\newcommand{\iniziobarra}{\setlength{\parskip}{4pt plus 1pt}\setcounter{passo}{0}}
\tcbset{barra/.style={blanker, breakable, parbox=false, left=10pt, right=0pt, top=1pt, bottom=1pt,
  before skip=9pt plus 2pt, after skip=9pt plus 2pt}}
\newtcolorbox{dimostrazione}[1][]{barra, borderline west={1.8pt}{0pt}{cdim},
  before upper={\iniziobarra\textbf{\color{cdim}Dimostrazione\tit{#1}.}\par}}
\newtcolorbox{approfondimento}[1][]{barra, borderline west={1.8pt}{0pt}{cappr},
  before upper={\iniziobarra\textbf{\color{cappr}Approfondimento\tit{#1}.}\ \ignorespaces}}
\newtcolorbox{esempio}[1][]{barra, borderline west={1.8pt}{0pt}{cesm},
  before upper={\iniziobarra\textbf{\color{cesm}Esempio\tit{#1}.}\ \ignorespaces}}
% niente titolo di dimostrazione/approfondimento isolato in fondo alla pagina
\AddToHook{env/dimostrazione/before}{\par\spazio{6\baselineskip}}
\AddToHook{env/approfondimento/before}{\par\spazio{5\baselineskip}}
\AddToHook{env/esempio/before}{\par\spazio{4\baselineskip}}
% conseguenze introdotte dalla domanda a cui rispondono
\newcommand{\domanda}[1]{\par\addvspace{5pt}\spazio{5\baselineskip}\noindent{\color{cris}$\blacktriangleright$}\ \textbf{#1}\par\nopagebreak[4]}
\newcommand{\nsp}[1]{\spazio{#1\baselineskip}}

% ---------------------------------------------------------- riquadri
%  enunciato: riquadro blu con l'enunciato sopra e, dopo \tcblower, il significato
%  \begin{teorema}[di Bloch] enunciato \tcblower significato \end{teorema}
%  (anche: modello). Numerazione comune con Definizione, Osservazione, ...
\newtcolorbox{enunciato}[1]{enhanced, unbreakable, parbox=false,
  colback=cenb, colframe=csez, boxrule=0.6pt, arc=2pt,
  colbacktitle=csez, coltitle=white, fonttitle=\bfseries, title={#1},
  left=9pt, right=9pt, top=6pt, bottom=6pt, toptitle=2pt, bottomtitle=2pt,
  middle=5pt, segmentation style={csez!45, dashed, line width=0.5pt},
  before skip=12pt plus 2pt, after skip=12pt plus 2pt,
  before upper={\setlength{\parskip}{5pt}},
  before lower={\setlength{\parskip}{5pt}\textbf{\color{cint}Cosa dice.}\ \ignorespaces}}
\newenvironment{teorema}[1][]{\par\refstepcounter{enun}\begin{enunciato}{Teorema \theenun\tit{#1}}}{\end{enunciato}}
\newenvironment{modello}[1][]{\par\refstepcounter{enun}\begin{enunciato}{Modello \theenun\tit{#1}}}{\end{enunciato}}
%  chiave: SOLO i risultati essenziali (barra del titolo verde, non spezzabile)
\newtcolorbox{chiave}[1][Risultato]{enhanced, unbreakable, parbox=false,
  colback=cbox, colframe=cris, boxrule=0.6pt, arc=2pt,
  colbacktitle=cris, coltitle=white, fonttitle=\bfseries, title={#1},
  left=9pt, right=9pt, top=5pt, bottom=5pt, toptitle=2pt, bottomtitle=2pt,
  before skip=11pt plus 2pt, after skip=11pt plus 2pt,
  before upper={\setlength{\parskip}{5pt}}}
%  formula da evidenziare senza titolo
\newcommand{\formula}[1]{\par\addvspace{4pt}\begin{center}\begin{tcolorbox}[enhanced, hbox, colback=cform,
  colframe=cformb, boxrule=0.5pt, arc=2pt, left=10pt, right=10pt, top=2pt, bottom=2pt]
  $\displaystyle #1$\end{tcolorbox}\end{center}\par\addvspace{2pt}}
%  informazioni sul corso o curiosità sull'hardware (rarissima)
\newcommand{\titinfo}[1]{\if\relax\detokenize{#1}\relax\else\textbf{#1.}\ \fi}
\newtcolorbox{infobox}[1][]{enhanced, breakable, parbox=false, colback=black!4, frame hidden,
  boxrule=0pt, arc=2pt, left=9pt, right=9pt, top=5pt, bottom=5pt,
  before upper={\setlength{\parskip}{4pt}\titinfo{#1}\ignorespaces}}
%  formulario: riquadro verde spezzabile, una \voce per riga (si spezza tra le voci)
\newtcolorbox{formulario}[1][Formule della lezione]{enhanced, breakable, parbox=false,
  colback=cbox, colframe=cris, boxrule=0.6pt, arc=2pt,
  colbacktitle=cris, coltitle=white, fonttitle=\bfseries, title={#1},
  left=9pt, right=9pt, top=5pt, bottom=7pt, before upper={\setlength{\parskip}{0pt}}}
\newcommand{\gruppo}[1]{\par\addvspace{6pt}\noindent\textbf{#1}\par\nopagebreak\addvspace{4pt}}
\newcommand{\sezvoce}[1]{\if\relax\detokenize{#1}\relax\else\ {\footnotesize\color{black!55}(\S#1)}\fi}
\newcommand{\voce}[3][]{\par\noindent\begin{tabularx}{\linewidth}{@{}>{\raggedright\arraybackslash}p{3.9cm}@{\hspace{8pt}}X@{}}
  #2\sezvoce{#1} & $\displaystyle #3$\end{tabularx}\par\addvspace{4pt}}
%  domande di verifica: \begin{domande}\item \textbf{Domanda?} traccia...\end{domande}
\newlist{domande}{enumerate}{1}
\setlist[domande]{label=\textbf{Q\arabic*.}, leftmargin=2.3em, itemsep=6pt}

% ---------------------------------------------------------- macro generiche
\newcommand{\ingl}[1]{\textit{\color{black!55}(#1)}}   % termine inglese delle slide
\newcommand{\va}{\mathbf a}
\newcommand{\vb}{\mathbf b}
\newcommand{\vk}{\mathbf k}
\newcommand{\vr}{\mathbf r}
\newcommand{\vR}{\mathbf R}
\newcommand{\vG}{\mathbf G}
\newcommand{\vc}{\mathbf c}
\newcommand{\vT}{\mathbf T}
\newcommand{\xh}{\hat{\mathbf x}}
\newcommand{\yh}{\hat{\mathbf y}}
\newcommand{\zh}{\hat{\mathbf z}}
\newcommand{\coord}[1]{\textsf{#1}}
\newcommand{\h}{\tfrac12}
\newcommand{\jj}{\mathrm j}
\newcommand{\ee}{\mathrm e}
\newcommand{\dd}{\mathrm d}
\newcommand{\pd}[2]{\frac{\partial #1}{\partial #2}}
\newcommand{\figura}[1]{\input{figs/#1.tex}}

% ---------------------------------------------------------- float
\renewcommand{\topfraction}{0.9}
\renewcommand{\bottomfraction}{0.8}
\renewcommand{\textfraction}{0.1}
\renewcommand{\floatpagefraction}{0.8}
\setcounter{topnumber}{3}
\setcounter{bottomnumber}{2}
\setcounter{totalnumber}{4}

% ---------------------------------------------------------- aggiunte della Parte 3
\newcommand{\EF}{E_{\mathrm F}}
\newcommand{\kF}{k_{\mathrm F}}
\newcommand{\kB}{k_{\mathrm B}}
\newcommand{\Eg}{E_{\mathrm g}}
\newcommand{\vion}{v_{\mathrm{ion}}}
\newcommand{\Kt}{\tilde K}
\definecolor{occ}{HTML}{2E6FB7}
\definecolor{hole}{HTML}{C0392B}
\definecolor{exc}{HTML}{23884A}
\definecolor{bandfull}{HTML}{8C2A1F}
\definecolor{bandempty}{HTML}{F2E6A0}
\definecolor{bandpart}{HTML}{D9C84A}
%  le tre approssimazioni del §6: stesso riquadro blu di teoremi e modelli
\newenvironment{approssimazione}[1][]{\par\refstepcounter{enun}\begin{enunciato}{Approssimazione \theenun\tit{#1}}}{\end{enunciato}}
\makeatletter
\setlength{\@fptop}{0pt}
\makeatother

% ---------------------------------------------------------- aggiunte della Parte 4
\usetikzlibrary{decorations.pathmorphing}
\newcommand{\hp}{\hat{\mathbf p}}
\newcommand{\vq}{\mathbf q}
% rimando esplicito alle dispense di Quantum Physics (V. Penna, PoliTo): \QP{§4.1, Teorema~2, p.~56}
\newcommand{\QP}[1]{\textsf{QP\_2025v2}, #1}
\newcommand{\cer}[1]{\text{\textcircled{\scriptsize #1}}}   % termini 1-4 della dimostrazione a mano
\usepgfplotslibrary{groupplots}
% ---------------------------------------------------------- aggiunte della Parte 4 (6a seconda parte, 5/11/2026)
\newcommand{\vF}{v_{\mathrm F}}                    % velocità degli elettroni del grafene (slide: c_g)
\newcommand{\RH}{R_{\mathrm H}}                    % coefficiente di Hall
 definisci \corso, \argomento, \lezioni.
%  Compilare con pdflatex (scripts/setup_pdflatex_it.sh per la sillabazione).
% ------------------------------------------------------------------
\providecommand{\corso}{Nome del corso}
\providecommand{\argomento}{Argomento}
\providecommand{\lezioni}{Lezioni}
\usepackage[T1]{fontenc}
\usepackage{babel}
\babelprovide[import, main]{italian}
\usepackage{lmodern}            % Latin Modern (se manca: apt-get install -y lmodern)
\usepackage{microtype}
\usepackage[a4paper, left=2.4cm, right=2.4cm, top=2.5cm, bottom=2.7cm]{geometry}
\usepackage[skip=6pt plus 2pt, indent=0pt]{parskip}

\usepackage{amsmath, amssymb}
\usepackage{bm}
\numberwithin{equation}{section}
\usepackage{siunitx}
\sisetup{output-decimal-marker={,}, group-separator={\,}, exponent-product=\cdot}
\DeclareSIUnit{\angstrom}{\text{\AA}}
\usepackage{booktabs, array, tabularx, multirow, longtable}
\usepackage{enumitem}
\setlist{itemsep=3pt, topsep=4pt, parsep=0pt}
\usepackage[font=small, labelfont=bf, format=hang, skip=6pt]{caption}
\usepackage{float}
\usepackage{needspace}
% \spazio{n\baselineskip}: va a capo pagina solo se restano meno di n righe.
% Sostituisce \needspace, il cui \penalty-100 con \vskip elastico fa spezzare
% le pagine prima dei titoli anche quando c'è posto (pagine mezze vuote).
\makeatletter
\newcommand{\spazio}[1]{\par\begingroup\setlength{\dimen@}{#1}%
  \dimen@ii\pagegoal\advance\dimen@ii-\pagetotal
  \ifdim\dimen@>\dimen@ii\ifdim\dimen@ii>\z@\vfil\penalty-10000\fi\fi\endgroup}
\makeatother

% ---------------------------------------------------------- colori
\usepackage{xcolor}
\definecolor{csez}{HTML}{193C82}   % titoli di sezione
\definecolor{csub}{HTML}{3C5995}   % sottosezioni, link
\definecolor{cdef}{HTML}{1E5AAA}   % Definizione
\definecolor{cris}{HTML}{14823C}   % Risultato, riquadro verde
\definecolor{cbox}{HTML}{F1F7F3}   % fondo del riquadro verde
\definecolor{cform}{HTML}{E4F0E8}  % fondo di \formula
\definecolor{cformb}{HTML}{ACD3BA} % bordo di \formula
\definecolor{coss}{HTML}{BE1E1E}   % Osservazione (refusi, convenzioni)
\definecolor{cint}{HTML}{007D87}   % Interpretazione fisica, «Cosa dice»
\definecolor{cappr}{HTML}{783296}  % Approfondimento
\definecolor{crie}{HTML}{C86E00}   % Riepilogo
\definecolor{cdim}{HTML}{6E6E6E}   % Dimostrazione
\definecolor{cpasso}{HTML}{585858} % Passo n.
\definecolor{cesm}{HTML}{4A6FA5}   % Esempio
\definecolor{cenb}{HTML}{EEF3FA}   % fondo del riquadro blu dell'enunciato
% colori per le figure
\definecolor{accent}{HTML}{193C82}
\definecolor{vA}{HTML}{C0392B}
\definecolor{vB}{HTML}{2E6FB7}
\definecolor{vC}{HTML}{23884A}
\definecolor{atomA}{HTML}{E15759}
\definecolor{atomB}{HTML}{4E79A7}
\definecolor{atomC}{HTML}{F28E2B}
\definecolor{atomD}{HTML}{59A14F}
\definecolor{atomE}{HTML}{B07AA1}
\definecolor{cellP}{HTML}{3C9D4E}
\definecolor{cellC}{HTML}{1F4E8C}
\definecolor{cellS}{HTML}{7B3FA0}
\definecolor{bad}{HTML}{C0392B}
\definecolor{Na}{HTML}{F28E2B}
\definecolor{Cl}{HTML}{4E79A7}
% nomi del preambolo precedente (figure e snippet esistenti)
\colorlet{boxdef}{cdef}\colorlet{boxkey}{cris}\colorlet{boxwarn}{coss}
\colorlet{boxex}{black!65}\colorlet{boxdeep}{cappr}\colorlet{mapgray}{black!60}

% ---------------------------------------------------------- figure
\usepackage{tikz}
\usetikzlibrary{arrows.meta, calc, decorations.pathreplacing, patterns}
\usepackage{pgfplots}
\pgfplotsset{compat=1.18, every axis/.append style={font=\footnotesize}, trig format plots=rad}
\usepackage[europeanresistors, cuteinductors, europeanvoltages, straightvoltages]{circuitikz}
\ctikzset{bipoles/length=1.0cm}
\usepackage{tikz-3dplot}
\tikzset{
  vec/.style={-{Stealth[length=2.4mm, width=1.7mm]}, thick},
  bond/.style={black!22, line width=0.5pt},
  latline/.style={black!30, thin},
  edge/.style={black!75, line width=0.7pt},
  edgehid/.style={black!55, line width=0.6pt, densely dashed},
  edgelight/.style={black!50, line width=0.6pt},
  edgelighthid/.style={black!40, line width=0.5pt, densely dashed},
  cellprim/.style={fill=cellP!16, draw=cellP!75!black, thick},
  cellconv/.style={fill=cellC!10, draw=cellC, thick},
  cellbad/.style={fill=bad!14, draw=bad, thick},
  cellconvline/.style={draw=cellC, thick, densely dashed},
  cellshift/.style={draw=cellS, thick, dashed},
  tri/.style={fill=vB!14, draw=vB!55, thin},
  tri0/.style={fill=vB!45, draw=vB!80!black, thick},
}

% ---------------------------------------------------------- titoli e indice
\usepackage{titlesec}
\titleformat{\section}{\spazio{9\baselineskip}\Large\bfseries\color{csez}}{\thesection}{0.8em}{}%
  [\vspace{-0.6ex}{\color{csez}\titlerule[0.7pt]}]
\titleformat{\subsection}{\spazio{8\baselineskip}\large\bfseries\color{csub}}{\thesubsection}{0.8em}{}
\titleformat{\subsubsection}{\spazio{4\baselineskip}\normalsize\bfseries\color{csub}}{}{0pt}{}
\titlespacing*{\section}{0pt}{22pt plus 4pt minus 2pt}{10pt plus 2pt}
\titlespacing*{\subsection}{0pt}{16pt plus 3pt minus 2pt}{6pt plus 2pt}
\titlespacing*{\subsubsection}{0pt}{10pt plus 2pt}{3pt}
\setcounter{tocdepth}{2}
\usepackage[titles]{tocloft}          % titolo «Indice» con lo stile di \section*
\renewcommand{\cftsecfont}{\bfseries}
\renewcommand{\cftsecpagefont}{\bfseries}
\setlength{\cftbeforesecskip}{11pt}
\setlength{\cftbeforesubsecskip}{3pt}

\pagestyle{plain}
\usepackage[colorlinks, linkcolor=csub, citecolor=csub, urlcolor=csub]{hyperref}
\hypersetup{pdftitle={\corso{} - \argomento}, pdfauthor={}}

% ---------------------------------------------------------- etichette numerate
%  Definizione, Risultato, Osservazione (e i riquadri Teorema, Modello) condividono il
%  contatore: Definizione 1.2, Osservazione 1.3, ...  (numerazione per sezione)
%  NB: \if...\else...\fi solo DENTRO le macro (\tit, \sezvoce): un \else scritto
%  dentro \textbf{...} viene accoppiato con l'\ifmmode di \textbf («Extra \else»).
\newcounter{enun}[section]
\renewcommand{\theenun}{\thesection.\arabic{enun}}
\newcommand{\tit}[1]{\if\relax\detokenize{#1}\relax\else\ (#1)\fi}   % « (titolo)» solo se c'è
\newcommand{\etinum}[3]{% colore, nome, titolo facoltativo
  \refstepcounter{enun}\textbf{\color{#1}#2 \theenun\tit{#3}.}\ \ignorespaces}
\newcommand{\etisemp}[3]{% colore, nome, titolo facoltativo (senza numero)
  \textbf{\color{#1}#2\tit{#3}.}\ \ignorespaces}
\newenvironment{definizione}[1][]{\par\addvspace{4pt}\spazio{3\baselineskip}\noindent\etinum{cdef}{Definizione}{#1}}{\par\addvspace{4pt}}
\newenvironment{risultato}[1][]{\par\addvspace{4pt}\spazio{3\baselineskip}\noindent\etinum{cris}{Risultato}{#1}}{\par\addvspace{4pt}}
\newenvironment{osservazione}[1][]{\par\addvspace{4pt}\spazio{3\baselineskip}\noindent\etinum{coss}{Osservazione}{#1}}{\par\addvspace{4pt}}
\newenvironment{interpretazione}[1][]{\par\addvspace{4pt}\spazio{3\baselineskip}\noindent\etisemp{cint}{Interpretazione fisica}{#1}}{\par\addvspace{4pt}}
\newenvironment{riepilogo}[1][]{\par\addvspace{4pt}\spazio{3\baselineskip}\noindent\etisemp{crie}{Riepilogo}{#1}}{\par\addvspace{4pt}}
% introduzione: titolo e testo in nero, spaziatura stretta
\newenvironment{introduzione}{\begingroup\colorlet{csez}{black}\colorlet{crie}{black}%
  \section*{Introduzione}\addcontentsline{toc}{section}{Introduzione}%
  \setlength{\parskip}{3pt plus 1pt}\setlist{itemsep=1pt, topsep=2pt, parsep=0pt}%
  \setlength{\abovedisplayskip}{4pt plus 2pt}\setlength{\belowdisplayskip}{4pt plus 2pt}%
  \setlength{\abovedisplayshortskip}{2pt}\setlength{\belowdisplayshortskip}{3pt}}%
  {\par\endgroup}
% compatibilità con le dispense precedenti
\newenvironment{attenzione}[1][]{\begin{osservazione}[#1]}{\end{osservazione}}

% ---------------------------------------------------------- dimostrazioni a passi
\usepackage[skins, breakable]{tcolorbox}
\newcounter{passo}
\makeatletter
\newcommand{\passo}[1]{\par\addvspace{5pt}\stepcounter{passo}%
  \noindent{\bfseries\color{cpasso}Passo \thepasso.}\ \textit{#1}\@ifnextchar\[{\nopagebreak}{\@ifnextchar\begin{\nopagebreak}{\par\nopagebreak[4]}}}
\makeatother
\newcommand{\qed}{\ifmmode\qquad\blacksquare\else\unskip\nobreak\ \ensuremath{\blacksquare}\fi}
\newcommand{\iniziobarra}{\setlength{\parskip}{4pt plus 1pt}\setcounter{passo}{0}}
\tcbset{barra/.style={blanker, breakable, parbox=false, left=10pt, right=0pt, top=1pt, bottom=1pt,
  before skip=9pt plus 2pt, after skip=9pt plus 2pt}}
\newtcolorbox{dimostrazione}[1][]{barra, borderline west={1.8pt}{0pt}{cdim},
  before upper={\iniziobarra\textbf{\color{cdim}Dimostrazione\tit{#1}.}\par}}
\newtcolorbox{approfondimento}[1][]{barra, borderline west={1.8pt}{0pt}{cappr},
  before upper={\iniziobarra\textbf{\color{cappr}Approfondimento\tit{#1}.}\ \ignorespaces}}
\newtcolorbox{esempio}[1][]{barra, borderline west={1.8pt}{0pt}{cesm},
  before upper={\iniziobarra\textbf{\color{cesm}Esempio\tit{#1}.}\ \ignorespaces}}
% niente titolo di dimostrazione/approfondimento isolato in fondo alla pagina
\AddToHook{env/dimostrazione/before}{\par\spazio{6\baselineskip}}
\AddToHook{env/approfondimento/before}{\par\spazio{5\baselineskip}}
\AddToHook{env/esempio/before}{\par\spazio{4\baselineskip}}
% conseguenze introdotte dalla domanda a cui rispondono
\newcommand{\domanda}[1]{\par\addvspace{5pt}\spazio{5\baselineskip}\noindent{\color{cris}$\blacktriangleright$}\ \textbf{#1}\par\nopagebreak[4]}
\newcommand{\nsp}[1]{\spazio{#1\baselineskip}}

% ---------------------------------------------------------- riquadri
%  enunciato: riquadro blu con l'enunciato sopra e, dopo \tcblower, il significato
%  \begin{teorema}[di Bloch] enunciato \tcblower significato \end{teorema}
%  (anche: modello). Numerazione comune con Definizione, Osservazione, ...
\newtcolorbox{enunciato}[1]{enhanced, unbreakable, parbox=false,
  colback=cenb, colframe=csez, boxrule=0.6pt, arc=2pt,
  colbacktitle=csez, coltitle=white, fonttitle=\bfseries, title={#1},
  left=9pt, right=9pt, top=6pt, bottom=6pt, toptitle=2pt, bottomtitle=2pt,
  middle=5pt, segmentation style={csez!45, dashed, line width=0.5pt},
  before skip=12pt plus 2pt, after skip=12pt plus 2pt,
  before upper={\setlength{\parskip}{5pt}},
  before lower={\setlength{\parskip}{5pt}\textbf{\color{cint}Cosa dice.}\ \ignorespaces}}
\newenvironment{teorema}[1][]{\par\refstepcounter{enun}\begin{enunciato}{Teorema \theenun\tit{#1}}}{\end{enunciato}}
\newenvironment{modello}[1][]{\par\refstepcounter{enun}\begin{enunciato}{Modello \theenun\tit{#1}}}{\end{enunciato}}
%  chiave: SOLO i risultati essenziali (barra del titolo verde, non spezzabile)
\newtcolorbox{chiave}[1][Risultato]{enhanced, unbreakable, parbox=false,
  colback=cbox, colframe=cris, boxrule=0.6pt, arc=2pt,
  colbacktitle=cris, coltitle=white, fonttitle=\bfseries, title={#1},
  left=9pt, right=9pt, top=5pt, bottom=5pt, toptitle=2pt, bottomtitle=2pt,
  before skip=11pt plus 2pt, after skip=11pt plus 2pt,
  before upper={\setlength{\parskip}{5pt}}}
%  formula da evidenziare senza titolo
\newcommand{\formula}[1]{\par\addvspace{4pt}\begin{center}\begin{tcolorbox}[enhanced, hbox, colback=cform,
  colframe=cformb, boxrule=0.5pt, arc=2pt, left=10pt, right=10pt, top=2pt, bottom=2pt]
  $\displaystyle #1$\end{tcolorbox}\end{center}\par\addvspace{2pt}}
%  informazioni sul corso o curiosità sull'hardware (rarissima)
\newcommand{\titinfo}[1]{\if\relax\detokenize{#1}\relax\else\textbf{#1.}\ \fi}
\newtcolorbox{infobox}[1][]{enhanced, breakable, parbox=false, colback=black!4, frame hidden,
  boxrule=0pt, arc=2pt, left=9pt, right=9pt, top=5pt, bottom=5pt,
  before upper={\setlength{\parskip}{4pt}\titinfo{#1}\ignorespaces}}
%  formulario: riquadro verde spezzabile, una \voce per riga (si spezza tra le voci)
\newtcolorbox{formulario}[1][Formule della lezione]{enhanced, breakable, parbox=false,
  colback=cbox, colframe=cris, boxrule=0.6pt, arc=2pt,
  colbacktitle=cris, coltitle=white, fonttitle=\bfseries, title={#1},
  left=9pt, right=9pt, top=5pt, bottom=7pt, before upper={\setlength{\parskip}{0pt}}}
\newcommand{\gruppo}[1]{\par\addvspace{6pt}\noindent\textbf{#1}\par\nopagebreak\addvspace{4pt}}
\newcommand{\sezvoce}[1]{\if\relax\detokenize{#1}\relax\else\ {\footnotesize\color{black!55}(\S#1)}\fi}
\newcommand{\voce}[3][]{\par\noindent\begin{tabularx}{\linewidth}{@{}>{\raggedright\arraybackslash}p{3.9cm}@{\hspace{8pt}}X@{}}
  #2\sezvoce{#1} & $\displaystyle #3$\end{tabularx}\par\addvspace{4pt}}
%  domande di verifica: \begin{domande}\item \textbf{Domanda?} traccia...\end{domande}
\newlist{domande}{enumerate}{1}
\setlist[domande]{label=\textbf{Q\arabic*.}, leftmargin=2.3em, itemsep=6pt}

% ---------------------------------------------------------- macro generiche
\newcommand{\ingl}[1]{\textit{\color{black!55}(#1)}}   % termine inglese delle slide
\newcommand{\va}{\mathbf a}
\newcommand{\vb}{\mathbf b}
\newcommand{\vk}{\mathbf k}
\newcommand{\vr}{\mathbf r}
\newcommand{\vR}{\mathbf R}
\newcommand{\vG}{\mathbf G}
\newcommand{\vc}{\mathbf c}
\newcommand{\vT}{\mathbf T}
\newcommand{\xh}{\hat{\mathbf x}}
\newcommand{\yh}{\hat{\mathbf y}}
\newcommand{\zh}{\hat{\mathbf z}}
\newcommand{\coord}[1]{\textsf{#1}}
\newcommand{\h}{\tfrac12}
\newcommand{\jj}{\mathrm j}
\newcommand{\ee}{\mathrm e}
\newcommand{\dd}{\mathrm d}
\newcommand{\pd}[2]{\frac{\partial #1}{\partial #2}}
\newcommand{\figura}[1]{\input{figs/#1.tex}}

% ---------------------------------------------------------- float
\renewcommand{\topfraction}{0.9}
\renewcommand{\bottomfraction}{0.8}
\renewcommand{\textfraction}{0.1}
\renewcommand{\floatpagefraction}{0.8}
\setcounter{topnumber}{3}
\setcounter{bottomnumber}{2}
\setcounter{totalnumber}{4}

% ---------------------------------------------------------- aggiunte della Parte 3
\newcommand{\EF}{E_{\mathrm F}}
\newcommand{\kF}{k_{\mathrm F}}
\newcommand{\kB}{k_{\mathrm B}}
\newcommand{\Eg}{E_{\mathrm g}}
\newcommand{\vion}{v_{\mathrm{ion}}}
\newcommand{\Kt}{\tilde K}
\definecolor{occ}{HTML}{2E6FB7}
\definecolor{hole}{HTML}{C0392B}
\definecolor{exc}{HTML}{23884A}
\definecolor{bandfull}{HTML}{8C2A1F}
\definecolor{bandempty}{HTML}{F2E6A0}
\definecolor{bandpart}{HTML}{D9C84A}
%  le tre approssimazioni del §6: stesso riquadro blu di teoremi e modelli
\newenvironment{approssimazione}[1][]{\par\refstepcounter{enun}\begin{enunciato}{Approssimazione \theenun\tit{#1}}}{\end{enunciato}}
\makeatletter
\setlength{\@fptop}{0pt}
\makeatother

% ---------------------------------------------------------- aggiunte della Parte 4
\usetikzlibrary{decorations.pathmorphing}
\newcommand{\hp}{\hat{\mathbf p}}
\newcommand{\vq}{\mathbf q}
% rimando esplicito alle dispense di Quantum Physics (V. Penna, PoliTo): \QP{§4.1, Teorema~2, p.~56}
\newcommand{\QP}[1]{\textsf{QP\_2025v2}, #1}
\newcommand{\cer}[1]{\text{\textcircled{\scriptsize #1}}}   % termini 1-4 della dimostrazione a mano
\usepgfplotslibrary{groupplots}
% ---------------------------------------------------------- aggiunte della Parte 4 (6a seconda parte, 5/11/2026)
\newcommand{\vF}{v_{\mathrm F}}                    % velocità degli elettroni del grafene (slide: c_g)
\newcommand{\RH}{R_{\mathrm H}}                    % coefficiente di Hall
 definisci \corso, \argomento, \lezioni.
%  Compilare con pdflatex (scripts/setup_pdflatex_it.sh per la sillabazione).
% ------------------------------------------------------------------
\providecommand{\corso}{Nome del corso}
\providecommand{\argomento}{Argomento}
\providecommand{\lezioni}{Lezioni}
\usepackage[T1]{fontenc}
\usepackage{babel}
\babelprovide[import, main]{italian}
\usepackage{lmodern}            % Latin Modern (se manca: apt-get install -y lmodern)
\usepackage{microtype}
\usepackage[a4paper, left=2.4cm, right=2.4cm, top=2.5cm, bottom=2.7cm]{geometry}
\usepackage[skip=6pt plus 2pt, indent=0pt]{parskip}

\usepackage{amsmath, amssymb}
\usepackage{bm}
\numberwithin{equation}{section}
\usepackage{siunitx}
\sisetup{output-decimal-marker={,}, group-separator={\,}, exponent-product=\cdot}
\DeclareSIUnit{\angstrom}{\text{\AA}}
\usepackage{booktabs, array, tabularx, multirow, longtable}
\usepackage{enumitem}
\setlist{itemsep=3pt, topsep=4pt, parsep=0pt}
\usepackage[font=small, labelfont=bf, format=hang, skip=6pt]{caption}
\usepackage{float}
\usepackage{needspace}
% \spazio{n\baselineskip}: va a capo pagina solo se restano meno di n righe.
% Sostituisce \needspace, il cui \penalty-100 con \vskip elastico fa spezzare
% le pagine prima dei titoli anche quando c'è posto (pagine mezze vuote).
\makeatletter
\newcommand{\spazio}[1]{\par\begingroup\setlength{\dimen@}{#1}%
  \dimen@ii\pagegoal\advance\dimen@ii-\pagetotal
  \ifdim\dimen@>\dimen@ii\ifdim\dimen@ii>\z@\vfil\penalty-10000\fi\fi\endgroup}
\makeatother

% ---------------------------------------------------------- colori
\usepackage{xcolor}
\definecolor{csez}{HTML}{193C82}   % titoli di sezione
\definecolor{csub}{HTML}{3C5995}   % sottosezioni, link
\definecolor{cdef}{HTML}{1E5AAA}   % Definizione
\definecolor{cris}{HTML}{14823C}   % Risultato, riquadro verde
\definecolor{cbox}{HTML}{F1F7F3}   % fondo del riquadro verde
\definecolor{cform}{HTML}{E4F0E8}  % fondo di \formula
\definecolor{cformb}{HTML}{ACD3BA} % bordo di \formula
\definecolor{coss}{HTML}{BE1E1E}   % Osservazione (refusi, convenzioni)
\definecolor{cint}{HTML}{007D87}   % Interpretazione fisica, «Cosa dice»
\definecolor{cappr}{HTML}{783296}  % Approfondimento
\definecolor{crie}{HTML}{C86E00}   % Riepilogo
\definecolor{cdim}{HTML}{6E6E6E}   % Dimostrazione
\definecolor{cpasso}{HTML}{585858} % Passo n.
\definecolor{cesm}{HTML}{4A6FA5}   % Esempio
\definecolor{cenb}{HTML}{EEF3FA}   % fondo del riquadro blu dell'enunciato
% colori per le figure
\definecolor{accent}{HTML}{193C82}
\definecolor{vA}{HTML}{C0392B}
\definecolor{vB}{HTML}{2E6FB7}
\definecolor{vC}{HTML}{23884A}
\definecolor{atomA}{HTML}{E15759}
\definecolor{atomB}{HTML}{4E79A7}
\definecolor{atomC}{HTML}{F28E2B}
\definecolor{atomD}{HTML}{59A14F}
\definecolor{atomE}{HTML}{B07AA1}
\definecolor{cellP}{HTML}{3C9D4E}
\definecolor{cellC}{HTML}{1F4E8C}
\definecolor{cellS}{HTML}{7B3FA0}
\definecolor{bad}{HTML}{C0392B}
\definecolor{Na}{HTML}{F28E2B}
\definecolor{Cl}{HTML}{4E79A7}
% nomi del preambolo precedente (figure e snippet esistenti)
\colorlet{boxdef}{cdef}\colorlet{boxkey}{cris}\colorlet{boxwarn}{coss}
\colorlet{boxex}{black!65}\colorlet{boxdeep}{cappr}\colorlet{mapgray}{black!60}

% ---------------------------------------------------------- figure
\usepackage{tikz}
\usetikzlibrary{arrows.meta, calc, decorations.pathreplacing, patterns}
\usepackage{pgfplots}
\pgfplotsset{compat=1.18, every axis/.append style={font=\footnotesize}, trig format plots=rad}
\usepackage[europeanresistors, cuteinductors, europeanvoltages, straightvoltages]{circuitikz}
\ctikzset{bipoles/length=1.0cm}
\usepackage{tikz-3dplot}
\tikzset{
  vec/.style={-{Stealth[length=2.4mm, width=1.7mm]}, thick},
  bond/.style={black!22, line width=0.5pt},
  latline/.style={black!30, thin},
  edge/.style={black!75, line width=0.7pt},
  edgehid/.style={black!55, line width=0.6pt, densely dashed},
  edgelight/.style={black!50, line width=0.6pt},
  edgelighthid/.style={black!40, line width=0.5pt, densely dashed},
  cellprim/.style={fill=cellP!16, draw=cellP!75!black, thick},
  cellconv/.style={fill=cellC!10, draw=cellC, thick},
  cellbad/.style={fill=bad!14, draw=bad, thick},
  cellconvline/.style={draw=cellC, thick, densely dashed},
  cellshift/.style={draw=cellS, thick, dashed},
  tri/.style={fill=vB!14, draw=vB!55, thin},
  tri0/.style={fill=vB!45, draw=vB!80!black, thick},
}

% ---------------------------------------------------------- titoli e indice
\usepackage{titlesec}
\titleformat{\section}{\spazio{9\baselineskip}\Large\bfseries\color{csez}}{\thesection}{0.8em}{}%
  [\vspace{-0.6ex}{\color{csez}\titlerule[0.7pt]}]
\titleformat{\subsection}{\spazio{8\baselineskip}\large\bfseries\color{csub}}{\thesubsection}{0.8em}{}
\titleformat{\subsubsection}{\spazio{4\baselineskip}\normalsize\bfseries\color{csub}}{}{0pt}{}
\titlespacing*{\section}{0pt}{22pt plus 4pt minus 2pt}{10pt plus 2pt}
\titlespacing*{\subsection}{0pt}{16pt plus 3pt minus 2pt}{6pt plus 2pt}
\titlespacing*{\subsubsection}{0pt}{10pt plus 2pt}{3pt}
\setcounter{tocdepth}{2}
\usepackage[titles]{tocloft}          % titolo «Indice» con lo stile di \section*
\renewcommand{\cftsecfont}{\bfseries}
\renewcommand{\cftsecpagefont}{\bfseries}
\setlength{\cftbeforesecskip}{11pt}
\setlength{\cftbeforesubsecskip}{3pt}

\pagestyle{plain}
\usepackage[colorlinks, linkcolor=csub, citecolor=csub, urlcolor=csub]{hyperref}
\hypersetup{pdftitle={\corso{} - \argomento}, pdfauthor={}}

% ---------------------------------------------------------- etichette numerate
%  Definizione, Risultato, Osservazione (e i riquadri Teorema, Modello) condividono il
%  contatore: Definizione 1.2, Osservazione 1.3, ...  (numerazione per sezione)
%  NB: \if...\else...\fi solo DENTRO le macro (\tit, \sezvoce): un \else scritto
%  dentro \textbf{...} viene accoppiato con l'\ifmmode di \textbf («Extra \else»).
\newcounter{enun}[section]
\renewcommand{\theenun}{\thesection.\arabic{enun}}
\newcommand{\tit}[1]{\if\relax\detokenize{#1}\relax\else\ (#1)\fi}   % « (titolo)» solo se c'è
\newcommand{\etinum}[3]{% colore, nome, titolo facoltativo
  \refstepcounter{enun}\textbf{\color{#1}#2 \theenun\tit{#3}.}\ \ignorespaces}
\newcommand{\etisemp}[3]{% colore, nome, titolo facoltativo (senza numero)
  \textbf{\color{#1}#2\tit{#3}.}\ \ignorespaces}
\newenvironment{definizione}[1][]{\par\addvspace{4pt}\spazio{3\baselineskip}\noindent\etinum{cdef}{Definizione}{#1}}{\par\addvspace{4pt}}
\newenvironment{risultato}[1][]{\par\addvspace{4pt}\spazio{3\baselineskip}\noindent\etinum{cris}{Risultato}{#1}}{\par\addvspace{4pt}}
\newenvironment{osservazione}[1][]{\par\addvspace{4pt}\spazio{3\baselineskip}\noindent\etinum{coss}{Osservazione}{#1}}{\par\addvspace{4pt}}
\newenvironment{interpretazione}[1][]{\par\addvspace{4pt}\spazio{3\baselineskip}\noindent\etisemp{cint}{Interpretazione fisica}{#1}}{\par\addvspace{4pt}}
\newenvironment{riepilogo}[1][]{\par\addvspace{4pt}\spazio{3\baselineskip}\noindent\etisemp{crie}{Riepilogo}{#1}}{\par\addvspace{4pt}}
% introduzione: titolo e testo in nero, spaziatura stretta
\newenvironment{introduzione}{\begingroup\colorlet{csez}{black}\colorlet{crie}{black}%
  \section*{Introduzione}\addcontentsline{toc}{section}{Introduzione}%
  \setlength{\parskip}{3pt plus 1pt}\setlist{itemsep=1pt, topsep=2pt, parsep=0pt}%
  \setlength{\abovedisplayskip}{4pt plus 2pt}\setlength{\belowdisplayskip}{4pt plus 2pt}%
  \setlength{\abovedisplayshortskip}{2pt}\setlength{\belowdisplayshortskip}{3pt}}%
  {\par\endgroup}
% compatibilità con le dispense precedenti
\newenvironment{attenzione}[1][]{\begin{osservazione}[#1]}{\end{osservazione}}

% ---------------------------------------------------------- dimostrazioni a passi
\usepackage[skins, breakable]{tcolorbox}
\newcounter{passo}
\makeatletter
\newcommand{\passo}[1]{\par\addvspace{5pt}\stepcounter{passo}%
  \noindent{\bfseries\color{cpasso}Passo \thepasso.}\ \textit{#1}\@ifnextchar\[{\nopagebreak}{\@ifnextchar\begin{\nopagebreak}{\par\nopagebreak[4]}}}
\makeatother
\newcommand{\qed}{\ifmmode\qquad\blacksquare\else\unskip\nobreak\ \ensuremath{\blacksquare}\fi}
\newcommand{\iniziobarra}{\setlength{\parskip}{4pt plus 1pt}\setcounter{passo}{0}}
\tcbset{barra/.style={blanker, breakable, parbox=false, left=10pt, right=0pt, top=1pt, bottom=1pt,
  before skip=9pt plus 2pt, after skip=9pt plus 2pt}}
\newtcolorbox{dimostrazione}[1][]{barra, borderline west={1.8pt}{0pt}{cdim},
  before upper={\iniziobarra\textbf{\color{cdim}Dimostrazione\tit{#1}.}\par}}
\newtcolorbox{approfondimento}[1][]{barra, borderline west={1.8pt}{0pt}{cappr},
  before upper={\iniziobarra\textbf{\color{cappr}Approfondimento\tit{#1}.}\ \ignorespaces}}
\newtcolorbox{esempio}[1][]{barra, borderline west={1.8pt}{0pt}{cesm},
  before upper={\iniziobarra\textbf{\color{cesm}Esempio\tit{#1}.}\ \ignorespaces}}
% niente titolo di dimostrazione/approfondimento isolato in fondo alla pagina
\AddToHook{env/dimostrazione/before}{\par\spazio{6\baselineskip}}
\AddToHook{env/approfondimento/before}{\par\spazio{5\baselineskip}}
\AddToHook{env/esempio/before}{\par\spazio{4\baselineskip}}
% conseguenze introdotte dalla domanda a cui rispondono
\newcommand{\domanda}[1]{\par\addvspace{5pt}\spazio{5\baselineskip}\noindent{\color{cris}$\blacktriangleright$}\ \textbf{#1}\par\nopagebreak[4]}
\newcommand{\nsp}[1]{\spazio{#1\baselineskip}}

% ---------------------------------------------------------- riquadri
%  enunciato: riquadro blu con l'enunciato sopra e, dopo \tcblower, il significato
%  \begin{teorema}[di Bloch] enunciato \tcblower significato \end{teorema}
%  (anche: modello). Numerazione comune con Definizione, Osservazione, ...
\newtcolorbox{enunciato}[1]{enhanced, unbreakable, parbox=false,
  colback=cenb, colframe=csez, boxrule=0.6pt, arc=2pt,
  colbacktitle=csez, coltitle=white, fonttitle=\bfseries, title={#1},
  left=9pt, right=9pt, top=6pt, bottom=6pt, toptitle=2pt, bottomtitle=2pt,
  middle=5pt, segmentation style={csez!45, dashed, line width=0.5pt},
  before skip=12pt plus 2pt, after skip=12pt plus 2pt,
  before upper={\setlength{\parskip}{5pt}},
  before lower={\setlength{\parskip}{5pt}\textbf{\color{cint}Cosa dice.}\ \ignorespaces}}
\newenvironment{teorema}[1][]{\par\refstepcounter{enun}\begin{enunciato}{Teorema \theenun\tit{#1}}}{\end{enunciato}}
\newenvironment{modello}[1][]{\par\refstepcounter{enun}\begin{enunciato}{Modello \theenun\tit{#1}}}{\end{enunciato}}
%  chiave: SOLO i risultati essenziali (barra del titolo verde, non spezzabile)
\newtcolorbox{chiave}[1][Risultato]{enhanced, unbreakable, parbox=false,
  colback=cbox, colframe=cris, boxrule=0.6pt, arc=2pt,
  colbacktitle=cris, coltitle=white, fonttitle=\bfseries, title={#1},
  left=9pt, right=9pt, top=5pt, bottom=5pt, toptitle=2pt, bottomtitle=2pt,
  before skip=11pt plus 2pt, after skip=11pt plus 2pt,
  before upper={\setlength{\parskip}{5pt}}}
%  formula da evidenziare senza titolo
\newcommand{\formula}[1]{\par\addvspace{4pt}\begin{center}\begin{tcolorbox}[enhanced, hbox, colback=cform,
  colframe=cformb, boxrule=0.5pt, arc=2pt, left=10pt, right=10pt, top=2pt, bottom=2pt]
  $\displaystyle #1$\end{tcolorbox}\end{center}\par\addvspace{2pt}}
%  informazioni sul corso o curiosità sull'hardware (rarissima)
\newcommand{\titinfo}[1]{\if\relax\detokenize{#1}\relax\else\textbf{#1.}\ \fi}
\newtcolorbox{infobox}[1][]{enhanced, breakable, parbox=false, colback=black!4, frame hidden,
  boxrule=0pt, arc=2pt, left=9pt, right=9pt, top=5pt, bottom=5pt,
  before upper={\setlength{\parskip}{4pt}\titinfo{#1}\ignorespaces}}
%  formulario: riquadro verde spezzabile, una \voce per riga (si spezza tra le voci)
\newtcolorbox{formulario}[1][Formule della lezione]{enhanced, breakable, parbox=false,
  colback=cbox, colframe=cris, boxrule=0.6pt, arc=2pt,
  colbacktitle=cris, coltitle=white, fonttitle=\bfseries, title={#1},
  left=9pt, right=9pt, top=5pt, bottom=7pt, before upper={\setlength{\parskip}{0pt}}}
\newcommand{\gruppo}[1]{\par\addvspace{6pt}\noindent\textbf{#1}\par\nopagebreak\addvspace{4pt}}
\newcommand{\sezvoce}[1]{\if\relax\detokenize{#1}\relax\else\ {\footnotesize\color{black!55}(\S#1)}\fi}
\newcommand{\voce}[3][]{\par\noindent\begin{tabularx}{\linewidth}{@{}>{\raggedright\arraybackslash}p{3.9cm}@{\hspace{8pt}}X@{}}
  #2\sezvoce{#1} & $\displaystyle #3$\end{tabularx}\par\addvspace{4pt}}
%  domande di verifica: \begin{domande}\item \textbf{Domanda?} traccia...\end{domande}
\newlist{domande}{enumerate}{1}
\setlist[domande]{label=\textbf{Q\arabic*.}, leftmargin=2.3em, itemsep=6pt}

% ---------------------------------------------------------- macro generiche
\newcommand{\ingl}[1]{\textit{\color{black!55}(#1)}}   % termine inglese delle slide
\newcommand{\va}{\mathbf a}
\newcommand{\vb}{\mathbf b}
\newcommand{\vk}{\mathbf k}
\newcommand{\vr}{\mathbf r}
\newcommand{\vR}{\mathbf R}
\newcommand{\vG}{\mathbf G}
\newcommand{\vc}{\mathbf c}
\newcommand{\vT}{\mathbf T}
\newcommand{\xh}{\hat{\mathbf x}}
\newcommand{\yh}{\hat{\mathbf y}}
\newcommand{\zh}{\hat{\mathbf z}}
\newcommand{\coord}[1]{\textsf{#1}}
\newcommand{\h}{\tfrac12}
\newcommand{\jj}{\mathrm j}
\newcommand{\ee}{\mathrm e}
\newcommand{\dd}{\mathrm d}
\newcommand{\pd}[2]{\frac{\partial #1}{\partial #2}}
\newcommand{\figura}[1]{\input{figs/#1.tex}}

% ---------------------------------------------------------- float
\renewcommand{\topfraction}{0.9}
\renewcommand{\bottomfraction}{0.8}
\renewcommand{\textfraction}{0.1}
\renewcommand{\floatpagefraction}{0.8}
\setcounter{topnumber}{3}
\setcounter{bottomnumber}{2}
\setcounter{totalnumber}{4}

% ---------------------------------------------------------- aggiunte della Parte 3
\newcommand{\EF}{E_{\mathrm F}}
\newcommand{\kF}{k_{\mathrm F}}
\newcommand{\kB}{k_{\mathrm B}}
\newcommand{\Eg}{E_{\mathrm g}}
\newcommand{\vion}{v_{\mathrm{ion}}}
\newcommand{\Kt}{\tilde K}
\definecolor{occ}{HTML}{2E6FB7}
\definecolor{hole}{HTML}{C0392B}
\definecolor{exc}{HTML}{23884A}
\definecolor{bandfull}{HTML}{8C2A1F}
\definecolor{bandempty}{HTML}{F2E6A0}
\definecolor{bandpart}{HTML}{D9C84A}
%  le tre approssimazioni del §6: stesso riquadro blu di teoremi e modelli
\newenvironment{approssimazione}[1][]{\par\refstepcounter{enun}\begin{enunciato}{Approssimazione \theenun\tit{#1}}}{\end{enunciato}}
\makeatletter
\setlength{\@fptop}{0pt}
\makeatother

% ---------------------------------------------------------- aggiunte della Parte 4
\usetikzlibrary{decorations.pathmorphing}
\newcommand{\hp}{\hat{\mathbf p}}
\newcommand{\vq}{\mathbf q}
% rimando esplicito alle dispense di Quantum Physics (V. Penna, PoliTo): \QP{§4.1, Teorema~2, p.~56}
\newcommand{\QP}[1]{\textsf{QP\_2025v2}, #1}
\newcommand{\cer}[1]{\text{\textcircled{\scriptsize #1}}}   % termini 1-4 della dimostrazione a mano
\usepgfplotslibrary{groupplots}
% ---------------------------------------------------------- aggiunte della Parte 4 (6a seconda parte, 5/11/2026)
\newcommand{\vF}{v_{\mathrm F}}                    % velocità degli elettroni del grafene (slide: c_g)
\newcommand{\RH}{R_{\mathrm H}}                    % coefficiente di Hall


\begin{document}

% ============================================================ frontespizio
\hypersetup{pageanchor=false}
\begin{titlepage}
\centering
\vspace*{5.2cm}
{\fontsize{24}{30}\selectfont\bfseries \corso\par}
\vspace{1.1cm}
{\Large \argomento\par}
\vspace{0.9cm}
{\large Legami forti con più orbitali e catena biatomica, metodo $\vk\cdot\mathbf p$,\\ lacune pesanti e leggere, pseudopotenziale,\\ bande di Al, Si e GaAs, riempimento delle bande,\\ semimetalli: bismuto e grafene, ibridazione sp$^2$,\\ pacchetti d'onde e quasi-impulso\par}
\vspace{0.35cm}
{\color{black!60} \lezioni{} \textperiodcentered{} A.A. 2026/2027\par}
\vspace{0.2cm}
{\small\color{black!60} Slide del corso, Ashcroft--Mermin capp.~8, 10 e 11, Kittel capp.~7--9, Hofmann capp.~5--6 e note integrative, Yu--Cardona cap.~2\par}
\vspace{2.6cm}
\begin{minipage}{0.86\linewidth}\small\setlength{\parskip}{0pt}
\textbf{Struttura del testo.} Gli enunciati sono numerati per sezione e contrassegnati da
un'etichetta colorata: \textbf{\color{cdef}Definizione}, \textbf{\color{cris}Risultato},
\textbf{\color{coss}Osservazione} (imprecisioni delle slide,
convenzioni di segno), \textbf{\color{cint}Interpretazione fisica}. Le dimostrazioni e gli
\textbf{\color{cappr}approfondimenti} (Kittel, Ashcroft--Mermin) sono segnalati da una barra
laterale; le dimostrazioni sono divise in passi numerati, ciascuno dei quali svolge un solo
passaggio. Gli enunciati di teoremi e modelli, con il loro significato, sono in un
\textbf{\color{csez}riquadro blu}; i \textbf{\color{cris}risultati essenziali} in un riquadro verde. I richiami di meccanica
quantistica rimandano alle dispense di Quantum Physics di V.~Penna (PoliTo, versione del 14 aprile 2025) con la sigla \QP{sezione e pagina}.
\end{minipage}
\vfill
\end{titlepage}
\hypersetup{pageanchor=true}

% ============================================================ indice
{\small\setlength{\cftbeforesecskip}{4pt}\setlength{\cftbeforesubsecskip}{0.5pt}\tableofcontents}
\newpage

% ============================================================ introduzione
\begin{introduzione}
Il filo della parte: \emph{dai modelli alle bande dei materiali veri}. Con più orbitali e più atomi per cella i legami forti danno bande di valenza e
di conduzione separate da un gap; vicino a $\Gamma$ il metodo $\vk\cdot\mathbf p$ ricava le masse efficaci e, con lo spin--orbita, le tre bande di
lacune dei semiconduttori; lo pseudopotenziale spiega perché gli elettroni quasi liberi danno le bande dei solidi veri; contando gli elettroni per cella
si capisce perché l'alluminio è un metallo e Si e GaAs sono semiconduttori, e perché bismuto e grafene stanno a metà strada. Si dà per nota la Parte~3 (teorema di Bloch, elettroni quasi liberi, legami forti, modello semiclassico); i rimandi sono
del tipo «Parte~3, §12.5».
\begin{enumerate}[label=\textbf{\Alph*.}, leftmargin=2.2em]
\item \textbf{Legami forti in tre dimensioni} (§\ref{sec:tb3d}). Con $M$ orbitali per atomo le bande vengono da un sistema $M\times M$, l'equazione
(I); la banda $s$ di una catena e del cubico semplice.
\item \textbf{La catena biatomica} (§\ref{sec:biat}). Due atomi diversi per cella: banda di valenza e banda di conduzione, con il gap $E_B-E_A$ al
bordo di zona. È un semiconduttore in una dimensione.
\item \textbf{Il metodo $\vk\cdot\mathbf p$} (§\ref{sec:kp}). Bande vicino a $\Gamma$ a partire dalle energie in $\Gamma$; la massa efficace è decisa
dalle bande vicine; lo spin--orbita separa lacune pesanti, leggere e split-off.
\item \textbf{Lo pseudopotenziale} (§\ref{sec:pseudo}). Le onde piane ortogonalizzate al core portano a un'equazione di Schrödinger per la sola
parte liscia della funzione di valenza, con un potenziale efficace debole: si tornano a usare gli elettroni quasi liberi, con $V_{\mathrm{ps}}$ al posto
di $V$; bande e masse di Si, Ge, GaAs.
\item \textbf{Bande dei materiali veri} (§\ref{sec:reali}). Schemi di zona; l'alluminio, metallo quasi libero; la fotoemissione che misura
$E(\vk)$; il gap assoluto di Si e GaAs; metallo se $g(\mu)\ne0$; una banda contiene 2 elettroni per cella, quindi un numero dispari di elettroni dà
un metallo.
\item \textbf{Semimetalli} (§\ref{sec:semimetalli}). Con un numero pari di elettroni una banda piena e una vuota possono sovrapporsi (bismuto: pochi
elettroni e lacune, effetto Hall enorme) o toccarsi in un punto (grafene: ibridi sp$^2$, cono di Dirac, $g(\EF)=0$ senza gap, massa efficace nulla).
\item \textbf{Trasporto} (§\ref{sec:trasporto}). Pacchetto d'onde di Bloch; il reticolo perfetto non fa resistenza; il quasi-impulso si conserva negli
urti a meno di $\hbar\mathbf G$.
\item \textbf{Esercizi delle slide} (§\ref{sec:eserslide}). Elettrone in una scatola; elettroni quasi liberi su un reticolo quadrato.
\end{enumerate}
\begin{riepilogo}[filo logico]
\begin{gather*}
\underbrace{(E-\bar E_m)\,b_m+\sum_n\sum_{\vR\ne0}e^{i\vk\cdot\vR}\,t_{mn}(\vR)\,b_n=0}_{\text{legami forti, (I)}}
\ \Longrightarrow\ \underbrace{E_\pm(k),\ \ \Eg=E_B-E_A}_{\text{catena }AB}\\
\Longrightarrow\ \underbrace{\frac1{m^*}=\frac1{m_0}+\frac2{m_0^2k^2}\sum_{n'\ne n}\frac{|\langle u_{n0}|\vk\cdot\hp|u_{n'0}\rangle|^2}{E_n(0)-E_{n'}(0)}}_{\vk\cdot\mathbf p:\ \text{massa dalle bande vicine}}
\\
\Longrightarrow\ \underbrace{\Bigl[\frac{\hp^2}{2m_0}+V+\sum_c(E-E_c)|\psi_{c\vk}\rangle\langle\psi_{c\vk}|\Bigr]\psi^{\mathrm f}_{\vk}=E\,\psi^{\mathrm f}_{\vk}}_{\text{pseudopotenziale debole}}\\
\Longrightarrow\ \underbrace{1\ \text{banda}=2N\ \text{elettroni};\ \ g(\mu)\ne0\ \text{o gap assoluto}}_{\text{metallo o semiconduttore}}
\ \Longrightarrow\ \underbrace{E_\pm=\pm\hbar\vF|\vq|}_{\text{grafene}}\\
\Longrightarrow\ \underbrace{\hbar\dot k=-e\mathcal E,\ \ \vk=\vk_f+\vq+\mathbf G}_{\text{trasporto}}
\end{gather*}
\end{riepilogo}
\end{introduzione}

% =================================================================
\section{Il modello a legami forti in tre dimensioni}\label{sec:tb3d}
% =================================================================
Il modello a legami forti della Parte~3 (§12) con tutti gli orbitali atomici insieme: la funzione di Bloch combina $M$ orbitali per atomo e le bande
si ottengono da un sistema lineare $M\times M$, l'equazione~(I) (slide 1--15, da Ashcroft--Mermin, cap.~10).

\begin{osservazione}[notazione]
Le slide scrivono $\psi_n(\vr-\vR)$ per gli orbitali atomici, $\Psi_{\vk}$ per la funzione di Bloch, $E_{\vk}$ o $\varepsilon(\vk)$ per la banda,
$\Delta U(\vr)$ (Ashcroft--Mermin) o $\sum_{\vR'}U(\vR,\vR')$ per il potenziale degli altri atomi, $\gamma_{mn}(\vR)$ per l'integrale di
sovrapposizione (nelle slide 12--15 $\alpha$, nella catena biatomica $\mu$, nell'esercizio delle slide 66--68 $\gamma$), $N$ per il numero di orbitali,
$\varepsilon$ per il livello atomico. Qui, come nella Parte~3: $\varphi_n$, $\psi_{\vk}$, $E(\vk)$, $\Delta V(\vr)$,
$t_{mn}(\vR)=-\int\varphi_m^*\,\Delta V\,\varphi_n(\vr-\vR)\,\dd\vr$ (quindi $\gamma_{mn}=\alpha=t$), $M$ orbitali per atomo, $n_{\mathrm{at}}$ atomi per
cella, $N$ celle, $E_n$ livello atomico e $\bar E_n$ livello spostato dagli altri atomi.
\end{osservazione}

\subsection{Il modello: atomi più interazioni tra primi vicini}
Il solido monoatomico in 3D ha l'Hamiltoniana vera $H=p^2/2m_0+V(\vr)$, con $V$ potenziale cristallino dotato di tutta la simmetria e la periodicità
del reticolo (slide 4). Il modello la approssima con le Hamiltoniane atomiche più le interazioni tra atomi:
\[
H\simeq\sum_{\vR}H_{\mathrm{at}}(\vR)+\sum_{\vR,\vR'}U(\vR,\vR'),
\]
con $H_{\mathrm{at}}(\vR)$ Hamiltoniana dell'atomo in $\vR$ e $U(\vR,\vR')$ interazione tra gli elettroni e gli atomi in $\vR$ e $\vR'$; più avanti
$U(\vR,\vR')=0$ se $\vR$ e $\vR'$ non sono primi vicini. $H$ è la somma delle Hamiltoniane a un elettrone, $H=\sum_eH_{1e}(\vR)$, con
\[
H_{1e}(\vR)=H_{\mathrm{at}}(\vR)+\sum_{\vR'}U(\vR,\vR');
\]
per l'elettrone dell'atomo nell'origine $H_{1e}=H_{\mathrm{at}}+\Delta V$, con $\Delta V(\vr)=\sum_{\vR\ne0}U(0,\vR)$ potenziale degli altri atomi.

\begin{modello}[a legami forti con $M$ orbitali per atomo]\label{mod:tb}
Ipotesi (slide 5--6):
\begin{enumerate}[leftmargin=*, label=\arabic*.]
\item l'Hamiltoniana a un elettrone è $H_{1e}=H_{\mathrm{at}}+\Delta V$;
\item il problema atomico è risolto: $H_{\mathrm{at}}\varphi_n=E_n\varphi_n$, con $n=s,p,d,\dots$ etichetta dell'orbitale ed $E_n$ energia
dell'elettrone atomico;
\item ogni $\varphi_n(\vr-\vR)$ è molto localizzato attorno a $\vR$, e gli orbitali dello stesso atomo sono ortonormali,
$\int\varphi_m^*\varphi_n\,\dd\vr=\delta_{mn}$;
\item (\textbf{ipotesi LCAO}) la funzione di Bloch è una combinazione lineare di orbitali atomici, con coefficienti $b_n$ da determinare:
\[
\psi_{\vk}(\vr)=\sum_{\vR}e^{i\vk\cdot\vR}f(\vr-\vR),\qquad f(\vr-\vR)=\sum_{n=1}^Mb_n\,\varphi_n(\vr-\vR);
\]
\item (usata in seguito) solo interazioni tra primi vicini.
\end{enumerate}
Allora i $b_n$ e la banda $E(\vk)$ risolvono, per $m=1,\dots,M$,
\begin{equation*}
\bigl[E(\vk)-\bar E_m\bigr]b_m+\sum_n\sum_{\vR\ne0}e^{i\vk\cdot\vR}\,t_{mn}(\vR)\,b_n=0,
\tag{I}
\end{equation*}
con l'integrale di sovrapposizione $t_{mn}(\vR)=-\int\varphi_m^*(\vr)\,\Delta V(\vr)\,\varphi_n(\vr-\vR)\,\dd\vr$.
\tcblower
Le funzioni di Bloch sono combinazioni di orbitali atomici \ingl{tight binding, LCAO}: ogni livello atomico dà una banda, centrata vicino al livello
e allargata dalle sovrapposizioni $t_{mn}$ con gli atomi vicini, che mescolano anche orbitali diversi. Il problema del cristallo si riduce, per ogni
$\vk$, a un sistema lineare $M\times M$.
\end{modello}

\subsection{Dimostrazione principale: l'equazione (I)}\label{sec:dimI}
Si proietta l'equazione di Schrödinger del cristallo su un orbitale dell'atomo nell'origine (slide 6--9; Ashcroft--Mermin, eq.~10.8--10.12). Nella
Parte~3 (§12.5) la (I) era ricavata in breve dalle funzioni di Bloch dei singoli orbitali; qui la strada delle slide.
\begin{dimostrazione}[l'equazione (I) dei legami forti]
\passo{Equazione di Schrödinger per un elettrone del cristallo (eq.~10.8).}
\[
H_{1e}\psi_{\vk}(\vr)=\bigl[H_{\mathrm{at}}+\Delta V(\vr)\bigr]\psi_{\vk}(\vr)=E(\vk)\,\psi_{\vk}(\vr).
\]
\passo{Proiettiamo sull'orbitale $\varphi_m$ dell'atomo nell'origine.}
Moltiplichiamo per $\varphi_m^*(\vr)$ e integriamo (in notazione di Dirac, $\langle\varphi_m|H_{1e}|\psi_{\vk}\rangle=E\langle\varphi_m|\psi_{\vk}\rangle$, slide 6;
\QP{§10.2, pp.~128--130}):
\[
\int\varphi_m^*\,H_{\mathrm{at}}\,\psi_{\vk}\,\dd\vr+\int\varphi_m^*\,\Delta V\,\psi_{\vk}\,\dd\vr=E(\vk)\int\varphi_m^*\,\psi_{\vk}\,\dd\vr .
\]
\passo{$H_{\mathrm{at}}$ è hermitiano: agisce a sinistra su $\varphi_m^*$ e dà $E_m$ (eq.~10.9).}
Definizione di operatore hermitiano, $(\psi,H\phi)=(H\psi,\phi)$: \QP{§2.4, p.~24}; $E_m$ è reale (\QP{§4.1, Teorema~1, p.~55}).
\[
\int\varphi_m^*(\vr)\,H_{\mathrm{at}}\psi_{\vk}(\vr)\,\dd\vr=\int\bigl(H_{\mathrm{at}}\varphi_m\bigr)^*\psi_{\vk}\,\dd\vr=E_m\int\varphi_m^*(\vr)\,\psi_{\vk}(\vr)\,\dd\vr .
\]
\passo{Portiamo $E_m$ a sinistra (eq.~10.10).}
\[
\bigl[E(\vk)-E_m\bigr]\int\varphi_m^*\,\psi_{\vk}\,\dd\vr=\int\varphi_m^*(\vr)\,\Delta V(\vr)\,\psi_{\vk}(\vr)\,\dd\vr .
\]
\passo{Sostituiamo la LCAO e separiamo $\vR=0$ da $\vR\ne0$ (eq.~10.12).}
Con $\int\varphi_m^*\varphi_n\,\dd\vr=\delta_{mn}$:
\begin{multline*}
\bigl[E-E_m\bigr]\Bigl[b_m+\sum_n\sum_{\vR\ne0}e^{i\vk\cdot\vR}\Bigl(\int\varphi_m^*(\vr)\,\varphi_n(\vr-\vR)\,\dd\vr\Bigr)b_n\Bigr]\\
=\sum_n\Bigl(\int\varphi_m^*\,\Delta V\,\varphi_n\,\dd\vr\Bigr)b_n+\sum_n\sum_{\vR\ne0}e^{i\vk\cdot\vR}\Bigl(\int\varphi_m^*(\vr)\,\Delta V(\vr)\,\varphi_n(\vr-\vR)\,\dd\vr\Bigr)b_n .
\end{multline*}
\passo{Trascuriamo le sovrapposizioni tra orbitali di atomi diversi.}
Gli orbitali sono localizzati: $\int\varphi_m^*(\vr)\,\varphi_n(\vr-\vR)\,\dd\vr\ll1$ per $\vR\ne0$, e a sinistra resta $[E-E_m]\,b_m$.
\passo{Il termine con $\vR=0$ sposta soltanto il livello atomico.}
$\Delta V$, dovuto agli altri atomi, varia molto più lentamente di $\varphi_n$ dove gli orbitali dell'origine sono apprezzabili: esce dall'integrale con
il valore $\Delta V(\vr_0)$ e, per l'ortogonalità, $\int\varphi_m^*\,\Delta V\,\varphi_n\,\dd\vr\simeq\Delta V(\vr_0)\,\delta_{mn}$. È una costante per
$b_m$ (slide 8: $\beta_mb_m$), che si assorbe nel livello, $E_m\to\bar E_m$:
\[
\bigl[E(\vk)-\bar E_m\bigr]b_m=\sum_n\sum_{\vR\ne0}e^{i\vk\cdot\vR}\Bigl(\int\varphi_m^*(\vr)\,\Delta V(\vr)\,\varphi_n(\vr-\vR)\,\dd\vr\Bigr)b_n .
\]
\passo{Chiamiamo $-t_{mn}(\vR)$ l'integrale rimasto (slide 9).}
È l'\textbf{integrale di sovrapposizione} \ingl{overlap energy integral}, l'analogo del $t$ della catena (Parte~3, §12.4):
\[
t_{mn}(\vR)=-\int\varphi_m^*(\vr)\,\Delta V(\vr)\,\varphi_n(\vr-\vR)\,\dd\vr .
\]
\passo{Conclusione.}
\[
\bigl[E(\vk)-\bar E_m\bigr]b_m+\sum_n\sum_{\vR\ne0}e^{i\vk\cdot\vR}\,t_{mn}(\vR)\,b_n=0 .\qed
\]
Proiettando su ciascun $\varphi_m$ si ottengono $M$ equazioni per gli $M$ coefficienti $b_n$.
\end{dimostrazione}

\begin{chiave}[Legami forti: il sistema (I)]
\[
\boxed{\ \bigl[E(\vk)-\bar E_m\bigr]b_m+\sum_{n=1}^M\sum_{\vR\ne0}e^{i\vk\cdot\vR}\,t_{mn}(\vR)\,b_n=0,\qquad m=1,\dots,M\ }
\]
Sistema lineare omogeneo nei $b_n$: ha soluzioni non nulle solo se il determinante dei coefficienti si annulla. Per ogni $\vk$ si trovano $M$ radici
$E(\vk)$, cioè $M$ bande.
\end{chiave}

\begin{osservazione}[le slide 1, 8 e 9]
Nella slide 8 «i primi due termini a destra sono più piccoli del terzo e si trascurano»: il primo (le sovrapposizioni, passo~6) lo è davvero; il
secondo (passo~7) non è piccolo in generale, ma è una costante che si assorbe nel livello, come dice la stessa slide. La slide 9 definisce
$-\gamma_{mn}(\vR)\equiv\langle\psi_m|U|\psi_n\rangle$, con il segno meno: $\gamma_{mn}=t_{mn}$. Le slide 1, 3 e 11 lo omettono, come la terza parte
della 5c (Parte~3, §12.5).
\end{osservazione}

\subsection{Conseguenze: quante bande, che cosa sono i parametri}
\domanda{Quante bande dà la (I)? (slide 1--2, 9--10)}
$M$ equazioni lineari omogenee: si annulla un determinante $M\times M$ e si ottengono \textbf{$M$ bande}, ciascuna con le radici in un orbitale
atomico. Si è assunto implicitamente un atomo per cella: con $n_{\mathrm{at}}$ atomi per cella la somma su $n$ corre su tutti gli orbitali della
cella, e equazioni e bande diventano $n_{\mathrm{at}}M$. Casi particolari («artificiali»):
\begin{itemize}[leftmargin=*]
\item \#1, un livello $s$ per atomo: 1 banda, di tipo $s$;
\item \#2, tre livelli $p$ per atomo: 3 bande, di tipo $p$;
\item \#3, un livello $s$ e tre $p$ con legami $sp^3$: 4 bande (nel Si, con 2 atomi per cella, 8 bande; Parte~3, §12.5).
\end{itemize}
Un atomo per cella con un orbitale dà una sola equazione; due atomi per cella con un orbitale ciascuno danno un sistema $2\times2$ (§\ref{sec:biat};
slide 9).

\domanda{Come si calcolano i $t_{mn}$?}
Con orbitali atomici veri gli integrali sono proibitivi: in pratica i $t_{mn}$ (e i livelli $\bar E_m$) sono \textbf{parametri} da adattare ai dati
sperimentali.

\domanda{Quando una banda resta vicina al livello atomico?}
Se i $t_{mn}$ sono «piccoli», ogni banda è un livello atomico più correzioni che dipendono da $\vk$: $E(\vk)\simeq E_n+{}$correzioni, cioè la
dispersione.

\subsection{La banda $s$: catena di atomi e cubico semplice}\label{sec:bandas}
Caso \#1 con le sole interazioni tra primi vicini (slide 3 e 11): un solo orbitale, quindi $b_m=1$ e la somma su $n$ ha il solo termine $n=m=s$:
\[
E(\vk)=\bar E_s-\sum_{\vR\ \text{p.v.}}t(\vR)\,e^{i\vk\cdot\vR}=\bar E_s-t\sum_{\vR\ \text{p.v.}}e^{i\vk\cdot\vR},
\]
perché $t(\vR)=t$ è lo stesso per tutti i primi vicini (p.v.). Nel cubico semplice di passo $a$ i 6 vicini sono in $\pm a\xh$, $\pm a\yh$, $\pm a\zh$:
\begin{equation}\label{eq:sc}
E(\vk)=\bar E_s-2t\bigl[\cos k_xa+\cos k_ya+\cos k_za\bigr],
\end{equation}
un coseno per ogni direzione (l'«effetto 3D» della slide 11).

\domanda{E in una dimensione? La catena monoatomica (slide 12--15)}
Le slide rifanno il caso della catena con l'Hamiltoniana scritta per esteso: catena infinita di passo $a$, un orbitale $s$ per atomo di energia nota
($\varepsilon\to\bar E_s$), solo primi vicini ($U(n,n')=0$ se $n'\ne n\pm1$):
\[
H=\frac{p^2}{2m_0}+V(x),\ \ V(x+a)=V(x)\qquad\longrightarrow\qquad H\simeq\sum_n\bigl[H_{\mathrm{at}}(n)+U(n,n-1)+U(n,n+1)\bigr].
\]
In notazione di Dirac (slide 15; l'elemento di matrice è un integrale nello spazio)
\[
E(k)=\langle\psi_k|H|\psi_k\rangle=\Bigl\langle\psi_k\Bigl|\sum_nH_{\mathrm{at}}(n)\Bigr|\psi_k\Bigr\rangle+\Bigl\langle\psi_k\Bigl|\sum_n\bigl[U(n,n-1)+U(n,n+1)\bigr]\Bigr|\psi_k\Bigr\rangle;
\]
con $H_{\mathrm{at}}(n)|\varphi_n\rangle=\bar E_s|\varphi_n\rangle$, orbitali ortogonali, $\langle\varphi_n|\varphi_{n'}\rangle=\delta_{nn'}$, e la sola sovrapposizione tra
primi vicini,
\[
\langle\varphi_{n''}|U(n,n\pm1)|\varphi_{n'}\rangle=-t\quad(n''=n,\ n'=n\pm1),\qquad 0\ \ \text{altrimenti},
\]
la (I) con $\vR=\pm a$ dà
\[
\boxed{\ E(k)=\bar E_s-2t\cos ka=\bar E_s-2t+4t\sin^2\frac{ka}2\ }\qquad\Bigl(\cos ka=1-2\sin^2\frac{ka}2\Bigr).
\]
Larghezza $4t$ e minimo $\bar E_s-2t$ in $k=0$; $E_T\equiv E(k)-\bar E_s+2t=4t\sin^2(ka/2)$ è l'energia misurata dal fondo della banda (figura
della slide 14, disegnata per $2t<\bar E_s$). Per orbitali $s$ ($\varphi>0$) e $\Delta V<0$ l'elemento di matrice è negativo, cioè $t>0$ (slide 15: «per $\alpha>0$ deve essere
negativo»). È la forma $E=A-B\cos ka$ del Kronig--Penney nell'approssimazione lineare (Parte~3, §10.6); come lì, $\bar E_s$ e $t$ si prendono come
parametri adattati alla banda sperimentale.

\begin{esempio}[esercizio delle slide 66--68: banda $s$ nel cubico semplice]
Con la notazione del Kittel, $E(\vk)=E_0-\alpha-\gamma\sum_je^{-i\vk\cdot\boldsymbol\rho_j}$, con $E_0$, $\alpha$, $\gamma$ energie positive e
$\boldsymbol\rho_j$ vettori verso i primi vicini: qui $E_0-\alpha=\bar E_s$ e $\gamma=t$. I sei vicini $\boldsymbol\rho=(\pm a,0,0)$, $(0,\pm a,0)$,
$(0,0,\pm a)$ ridanno la (\ref{eq:sc}); lungo le direzioni principali, con $\vk=\frac\pi a(k_x,k_y,k_z)$ ai bordi (figura~\ref{fig:tbsc}):
\begin{itemize}[leftmargin=*]
\item $\Gamma$, $\vk=0$: $E(\Gamma)=\bar E_s-6t$;
\item $[100]$, $\Gamma X$, $\vk=(k,0,0)$: $E=\bar E_s-4t-2t\cos ka$; $X=\frac\pi a(1,0,0)$ è sul bordo, $\cos ka=-1$: $E(X)=\bar E_s-2t$;
\item $[110]$, $\Gamma M$, $k_x=k_y=k/\sqrt2$: $E=\bar E_s-2t-4t\cos(ka/\sqrt2)$; in $M=\frac\pi a(1,1,0)$, $k^2=(\pi/a)^2(1+1+0)$: $E(M)=\bar E_s+2t$;
\item $[111]$, $\Gamma R$, $k_x=k_y=k_z=k/\sqrt3$: $E=\bar E_s-6t\cos(ka/\sqrt3)$; in $R=\frac\pi a(1,1,1)$, $k^2=3(\pi/a)^2$: $E(R)=\bar E_s+6t$.
\end{itemize}
Larghezza $12t$. Lungo $\Gamma X$ e $\Gamma M$ le curve sono quelle del reticolo quadrato, $\bar E_s-2t(\cos k_xa+\cos k_ya)$, spostate di $-2t$.
\end{esempio}

\begin{figure}[tbp]
\centering
\figura{tbsc}
\caption{Banda $s$ a legami forti nel cubico semplice lungo $R$--$\Gamma$--$X$--$R$ (figura~15 della slide 68): dal fondo $\bar E_s-6t$ in $\Gamma$ alla
cima $\bar E_s+6t$ nel vertice $R$ della zona; in $X$, centro di una faccia, $E=\bar E_s-2t$. In ascissa le lunghezze dei tratti.}\label{fig:tbsc}
\end{figure}

\begin{osservazione}[slide 66--68]
Le slide chiamano $W$ il vertice $\frac\pi a(1,1,1)$ della zona cubica: nel cubico semplice si chiama $R$ ($W$ è un punto della zona del reticolo FCC).
La slide 67 nota che, a meno di una traslazione in energia, lungo $[100]$ e $[110]$ si ottengono gli stessi risultati del reticolo quadrato; lì però
$E(\Gamma)=E_0-\alpha-4\gamma$ (la slide scrive $-6\gamma$, il valore del cubico).
\end{osservazione}

% =================================================================
\section{La catena biatomica: un semiconduttore in una dimensione}\label{sec:biat}
% =================================================================
Due atomi diversi per cella, un orbitale $s$ ciascuno: la (I) diventa un sistema $2\times2$ e dà due bande, una di valenza e una di conduzione,
separate da un gap. È il modello più semplice di semiconduttore (slide 16--32; dimostrazione a mano nelle slide 21--27).

\begin{osservazione}[notazione]
Slide: $\varepsilon_A,\varepsilon_B\to E_A,E_B$ (livelli atomici, già spostati; $E_A<E_B$ come nella slide 30); $C_A,C_B\to b_A,b_B$ (i $b_n$ della
(I)); $\mu\equiv\langle\psi_{An}|U(n,n-1)|\psi_{B,n-1}\rangle=-\gamma_{AB}\to-t$, con $t>0$ per orbitali $s$ come nel §\ref{sec:tb3d};
$\psi_{\gamma n}\to\varphi_{\gamma n}$, orbitale dell'atomo $\gamma=A,B$ della cella $n$; $E_{\vk}\to E(k)$; $H$ (larghezza delle bande, slide 30--31)
$\to W$, per non confonderla con l'Hamiltoniana.
\end{osservazione}

\subsection{Il modello}
Catena infinita di due specie atomiche alternate, $A$ e $B$ (cristallo con due atomi per cella), un elettrone di valenza $s$ per atomo, distanza di
ripetizione della cella $a$; solo interazioni tra primi vicini: quelle tra atomi più lontani di circa $a/2$ sono trascurabili (slide 16). È la
\textbf{catena biatomica} \ingl{diatomic chain} (figura~\ref{fig:catena}).

\begin{modello}[catena biatomica]\label{mod:biat}
Catena di passo $a$ con gli atomi $A$ in $x_{A_n}=na$ e $B$ in $x_{B_n}=(n+\frac12)a$; un orbitale $s$ per atomo,
$H_{\mathrm{at}}(\gamma,n)\varphi_{\gamma n}=E_\gamma\varphi_{\gamma n}$ ($\gamma=A,B$, $E_A<E_B$), orbitali localizzati e ortonormali; sovrapposizione
solo tra primi vicini $A$--$B$, $\langle\varphi_{A}|U|\varphi_{B}\rangle=-t$. Funzione di Bloch LCAO (slide 19), con i coefficienti $b_A$, $b_B$
incogniti:
\[
\psi_k(x)=\sum_ne^{ikna}\bigl[b_A\,\varphi_{An}(x)+b_B\,\varphi_{Bn}(x)\bigr],\qquad \varphi_{\gamma n}(x)=\varphi_\gamma(x-x_{\gamma_n}).
\]
Le bande sono
\[
E_\pm(k)=\frac{E_A+E_B}2\pm\sqrt{\Bigl(\frac{E_B-E_A}2\Bigr)^2+4t^2\cos^2\frac{ka}2}\ .
\]
\tcblower
Due atomi per cella danno due bande: quella bassa (valenza) viene soprattutto dal livello $E_A$, quella alta (conduzione) da $E_B$; la sovrapposizione
$t$ le mescola e le allontana. Con un elettrone per atomo la valenza è piena e la conduzione vuota: un «semiconduttore» in una dimensione, con gap
$E_B-E_A$ al bordo di zona.
\end{modello}
Le slide 17--19 partono dall'Hamiltoniana: quella vera, $H=p^2/2m_0+V(x)$ con $V(x+a)=V(x)$, si approssima con le Hamiltoniane atomiche $H_{\mathrm{at}}(\gamma,n)$
dell'atomo $\gamma$ nella cella $n$ e le interazioni $U(\gamma,n;\gamma',n')$ tra elettrone e atomo $\gamma$ della cella $n$ ed elettrone e atomo
$\gamma'$ della cella $n'$:
\[
H\simeq\sum_{\gamma,n}H_{\mathrm{at}}(\gamma,n)+\sum_{\gamma n,\gamma'n'}U(\gamma,n;\gamma',n')\ \ \longrightarrow\ \
H\simeq\sum_{\gamma,n}H_{\mathrm{at}}(\gamma,n)+\sum_n\bigl[U(n,n-1)+U(n,n+1)\bigr]:
\]
con i soli primi vicini restano le interazioni $A$--$B$, scritte $U(n,n\pm1)$ (con il vicino a sinistra e con quello a destra).

\begin{figure}[tbp]
\centering
\figura{catena}
\caption{Catena biatomica di passo $a$: l'atomo $A_n$ ha come primi vicini $B_{n-1}$ e $B_n$, a distanza $a/2$; l'atomo $B_n$ ha $A_n$ e $A_{n+1}$.
Le frecce sono gli integrali di sovrapposizione $-t$ tra primi vicini.}\label{fig:catena}
\end{figure}


\subsection{Dimostrazione principale: le equazioni per $b_A$ e $b_B$}
È la dimostrazione a mano delle slide 21--27 (pagine 1--7): si proietta l'equazione di Schrödinger sull'orbitale di un atomo $A$, con i vicini
dell'osservazione~\ref{oss:vicini}.
\begin{dimostrazione}[equazioni (1a) e (2a) delle slide]
\passo{Hamiltoniana a un elettrone (pag.~1).}
Solo primi vicini: $\sum_{n'}U(n,n')=U(n,n-1)+U(n,n+1)$. Con due specie $H_{\mathrm{at}}(n)=H_{\mathrm{at}}(A,n)+H_{\mathrm{at}}(B,n)$, e ciascuna agisce
solo sugli orbitali della propria specie:
\[
H_{1e}(n)=H_{\mathrm{at}}(A,n)+H_{\mathrm{at}}(B,n)+U(n,n-1)+U(n,n+1).
\]
\passo{Funzione di Bloch e proiezione su $\langle\varphi_{An}|$ (pag.~2).}
Con $|\psi_k\rangle=\sum_me^{ikma}\bigl(b_A|\varphi_{Am}\rangle+b_B|\varphi_{Bm}\rangle\bigr)$ e $H_{1e}|\psi_k\rangle=E|\psi_k\rangle$, l'uguaglianza
proiettata si scrive come quattro contributi:
\[
\underbrace{\langle\varphi_{An}|H_{\mathrm{at}}(A,n)|\psi_k\rangle}_{\cer1}+\underbrace{\langle\varphi_{An}|H_{\mathrm{at}}(B,n)|\psi_k\rangle}_{\cer2}
+\underbrace{\langle\varphi_{An}|U(n,n-1)+U(n,n+1)|\psi_k\rangle}_{\cer3}=\underbrace{E\,\langle\varphi_{An}|\psi_k\rangle}_{\cer4}.
\]
\passo{Termine \cer1 (pag.~3).}
$H_{\mathrm{at}}(A,n)$ agisce sugli orbitali $A$ e dà $E_A$ (hermiticità, come nel §\ref{sec:dimI}); poi $\langle\varphi_{An}|\varphi_{Am}\rangle=\delta_{nm}$
e $\langle\varphi_{An}|\varphi_{Bm}\rangle=0$:
\[
\cer1=\sum_me^{ikma}\langle\varphi_{An}|\varphi_{Am}\rangle\,E_A\,b_A=e^{ikna}E_A\,b_A .
\]
\passo{Termine \cer2 (pag.~4): è nullo.}
$H_{\mathrm{at}}(B,n)$ agisce sugli orbitali $B$, e $|\varphi_{An}\rangle\perp|\varphi_{Bm}\rangle$ sempre ($A\ne B$):
\[
\cer2=\sum_me^{ikma}\langle\varphi_{An}|\varphi_{Bm}\rangle\,E_B\,b_B=0 .
\]
\passo{Termine \cer4 (pag.~4).}
\[
\cer4=E\sum_me^{ikma}\bigl[b_A\langle\varphi_{An}|\varphi_{Am}\rangle+b_B\langle\varphi_{An}|\varphi_{Bm}\rangle\bigr]=E\,e^{ikna}\,b_A .
\]
\passo{Termine \cer3 (pag.~5--6): restano i due atomi $B$ vicini di $A_n$.}
L'interazione collega $\varphi_{An}$ solo agli orbitali dei primi vicini, $B_{n-1}$ (a sinistra, nella cella $n-1$) e $B_n$ (a destra, nella stessa
cella), con lo stesso integrale $-t$:
\begin{multline*}
\cer3=\langle\varphi_{An}|U(n,n-1)|\varphi_{B,n-1}\rangle\,e^{ik(n-1)a}\,b_B+\langle\varphi_{An}|U(n,n+1)|\varphi_{Bn}\rangle\,e^{ikna}\,b_B\\
=-t\bigl(e^{ik(n-1)a}+e^{ikna}\bigr)b_B .
\end{multline*}
\passo{Sommiamo $\cer1+\cer2+\cer3=\cer4$ e dividiamo per il fattore comune $e^{ikna}$ (pag.~7).}
\[
E_A\,b_A-t\bigl(1+e^{-ika}\bigr)b_B=E\,b_A .\qquad\text{(1a)}
\]
\passo{Stesso calcolo proiettando su $\langle\varphi_{Bn}|$: i vicini di $B_n$ sono $A_n$ e $A_{n+1}$.}
\[
E_B\,b_B-t\bigl(1+e^{ika}\bigr)b_A=E\,b_B .\qquad\text{(2a)}\qed
\]
\end{dimostrazione}

\begin{osservazione}[i vicini di $A_n$ nelle slide 20 e 25--30]\label{oss:vicini}
Le slide (e la pagina 5 della dimostrazione a mano, «catena BABABAB») prendono come vicini di $A_n$ gli atomi $B_{n-1}$ e $B_{n+1}$, con fasi
$e^{ik(n\mp1)a}$: ottengono $\mu\bigl[e^{ik(n-1)a}+e^{ik(n+1)a}\bigr]$ e poi $2\mu\cos ka$. Con il passo $a$ della slide 16 ($A$ e $B$ a distanza
$a/2$) i vicini di $A_n$ sono $B_{n-1}$ e $B_n$, e il fattore è $1+e^{-ika}$, di modulo $2|\cos(ka/2)|$. Le formule delle slide valgono se $a$ è la
distanza $A$--$B$ (cella di lato $2a$, bordo di zona in $\pi/2a$); con il passo $a$ si legge $\cos(ka/2)$ al posto di $\cos ka$, e il gap cade al bordo
$k=\pi/a$, come nella figura della slide 31 (verificato diagonalizzando una catena di 60 celle).
\end{osservazione}

\subsection{Le due bande}
$E_A$, $E_B$ e $t$ sono di solito parametri della teoria, non calcolati dagli orbitali (slide 29). Le (1a)--(2a) sono lineari e omogenee in $b_A$ e
$b_B$.
\begin{dimostrazione}[bande della catena biatomica (slide 29--30)]
\passo{Sistema omogeneo (eq.~(3)--(4) delle slide).}
\[
\begin{cases}(E_A-E)\,b_A-t\bigl(1+e^{-ika}\bigr)\,b_B=0 & \text{(atomo $A$)}\\[3pt] -t\bigl(1+e^{ika}\bigr)\,b_A+(E_B-E)\,b_B=0 & \text{(atomo $B$)}\end{cases}
\]
\passo{Determinante $2\times2$ nullo.}
Con $\bigl|1+e^{ika}\bigr|^2=2+2\cos ka=4\cos^2(ka/2)$ si ottiene un'equazione di secondo grado in $E$:
\[
(E_A-E)(E_B-E)-4t^2\cos^2\frac{ka}2=0 .
\]
\passo{Risolviamo.}
$E^2-(E_A+E_B)E+E_AE_B-4t^2\cos^2(ka/2)=0$, e $\frac14(E_A+E_B)^2-E_AE_B=\frac14(E_B-E_A)^2$:
\[
E_\pm(k)=\frac{E_A+E_B}2\pm\sqrt{\Bigl(\frac{E_B-E_A}2\Bigr)^2+4t^2\cos^2\frac{ka}2}\ .\qed
\]
\end{dimostrazione}
\begin{chiave}[Bande della catena biatomica]
\[
\boxed{\begin{gathered}E_\pm(k)=\frac{E_A+E_B}2\pm\sqrt{\Bigl(\frac{E_B-E_A}2\Bigr)^2+4t^2\cos^2\frac{ka}2}\\[2pt]
\Eg=E_+\bigl(\tfrac\pi a\bigr)-E_-\bigl(\tfrac\pi a\bigr)=E_B-E_A\end{gathered}}
\]
Due soluzioni: la \textbf{banda di valenza} $E_-$ e la \textbf{banda di conduzione} $E_+$, separate da un gap diretto al bordo di zona.
\end{chiave}
Con le fasi sulle posizioni degli atomi ($b_B\to b_Be^{ika/2}$) il sistema diventa reale e simmetrico,
$(E_A-E)b_A-2t\cos(ka/2)\,b_B=0$ e $-2t\cos(ka/2)\,b_A+(E_B-E)b_B=0$: sono le (3)--(4) delle slide con $\mu\to-t$ e $\cos ka\to\cos(ka/2)$.

\subsection{Conseguenze: gap, larghezze, carattere degli stati}
\domanda{Dov'è il gap e quanto vale? (slide 30--31)}
$\cos(ka/2)$ si annulla al bordo $k=\pm\pi/a$: lì $E_+=E_B$ ed $E_-=E_A$, e la distanza tra le bande è minima, $\Eg=E_B-E_A$. In $k=0$ le bande sono
invece alla massima distanza (la slide 30 la chiama «gap in $k=0$»):
\[
E_+(0)-E_-(0)=2\Bigl[\tfrac14(E_B-E_A)^2+4t^2\Bigr]^{1/2}.
\]
Il gap esiste solo perché $A\ne B$: lo apre la differenza tra i livelli atomici dei due atomi.

\domanda{Quanto sono larghe le bande? (slide 30--31)}
Se $4t^2\ll\frac14(E_B-E_A)^2$ si sviluppa la radice al primo ordine, $\sqrt{D^2+x}\simeq D+x/2D$ con $D=\frac12(E_B-E_A)$:
\[
E_+(k)\simeq E_B+W\cos^2\frac{ka}2,\qquad E_-(k)\simeq E_A-W\cos^2\frac{ka}2,\qquad W\equiv\frac{4t^2}{E_B-E_A}.
\]
Ogni banda è larga circa $W$; la conduzione sta \emph{sopra} $E_B$ (tra $E_B$ e $E_B+W$), la valenza \emph{sotto} $E_A$ (tra $E_A-W$ ed $E_A$): la
sovrapposizione respinge i due livelli, come legante e antilegante nella molecola $AB$ (Parte~3, §12.1). In $k=0$ la distanza è $\Eg+2W$
(figura~\ref{fig:biat}a); bande strette se $t\ll E_B-E_A$.

\begin{figure}[tbp]
\centering
\figura{biatomica}
\caption{Bande della catena biatomica. (a) $E_A<E_B$ (qui $t=0{,}3\,(E_B-E_A)$): gap $\Eg=E_B-E_A$ al bordo di zona (in giallo), bande larghe $W$
fuori dall'intervallo $[E_A,E_B]$, distanza massima $\Eg+2W$ in $k=0$. (b) $E_A=E_B$: le bande si toccano in $k=\pm\pi/a$ e il gap scompare.}\label{fig:biat}
\end{figure}

\begin{osservazione}[slide 30--31: quadrati e segni]
La slide 30 scrive gli sviluppi $\varepsilon_B+H[\cos ka]$ e $\varepsilon_A-H[\cos ka]$, senza il quadrato, e il gap $E_G(\pi/a)\simeq\varepsilon_A
-\varepsilon_B+2H$; la slide 31 disegna le bande \emph{dentro} l'intervallo $[\varepsilon_A,\varepsilon_B]$ e scrive $E_G\simeq\varepsilon_A-
\varepsilon_B-2H$. Dalla formula esatta $E_+\ge E_B$ ed $E_-\le E_A$ sempre, perché la radice vale almeno $\frac12(E_B-E_A)$: le bande stanno fuori
dall'intervallo, il gap al bordo è esattamente $E_B-E_A>0$, ed $E_B-E_A+2W$ è la distanza in $k=0$.
\end{osservazione}

\domanda{Su quali atomi stanno gli elettroni? Valenza e conduzione}
In $k=\pm\pi/a$ il termine misto si annulla: lo stato di valenza è puramente $A$ ($b_B=0$), quello di conduzione puramente $B$. Per $t\ll E_B-E_A$ la
valenza ha carattere $A$ e la conduzione carattere $B$ in tutta la zona (con $E_B-E_A=4t$, in $k=0$ la valenza ha $|b_A|^2=0{,}85$). Con un elettrone
per atomo ci sono 2 elettroni per cella: riempiono i 2 posti per cella della valenza (con lo spin), e la conduzione resta vuota. Sopra la banda piena
c'è il gap $E_B-E_A$: un isolante o, se il gap è piccolo, un semiconduttore, con gli elettroni di valenza sull'atomo di livello più basso.

\domanda{E se i due atomi sono uguali?}
Per $E_A=E_B$ il gap si chiude: $E_\pm=E_A\pm2t|\cos(ka/2)|$ si toccano in $k=\pm\pi/a$ (figura~\ref{fig:biat}b). È la banda della catena monoatomica
di passo $a/2$, $E_A-2t\cos(ka/2)$, ripiegata nella zona $|k|\le\pi/a$ della cella doppia (Parte~3, §9.3); con un elettrone per atomo è piena a metà:
un metallo.

\begin{risultato}[masse efficaci al gap]
Vicino al bordo, $k=\pi/a+q$, $\cos^2(ka/2)\simeq a^2q^2/4$ e lo sviluppo dà $E_\pm\simeq E_{B,A}\pm t^2a^2q^2/(E_B-E_A)$ (non nelle slide):
\[
m^*_c=-m^*_v=\frac{\hbar^2(E_B-E_A)}{2t^2a^2}.
\]
Gap piccolo, masse piccole: la stessa tendenza $m^*\propto\Eg$ dei quasi liberi (Parte~3, §13.3) e del metodo $\vk\cdot\mathbf p$ (§\ref{sec:kp}).
\end{risultato}

\domanda{Come si ritrova dalla (I)? (slide 32)}
Si pone $b_m=b_A$, $b_n=b_B$ e la somma su $\vR$ corre sui due vicini $B$ di $A_n$ (solo primi vicini). Proiettando su $\varphi_A$:
\[
(E-E_A)\,b_A+t_{AB}\bigl(1+e^{-ika}\bigr)b_B=0,\qquad t_{AB}=-\langle\varphi_A|\Delta V|\varphi_B\rangle=t ,
\]
cioè la (1a); proiettando su $\varphi_B$ si ritrova la (2a). La slide 32 scrive $\sum_n\sum_{\vR\ne0}e^{i\vk\cdot\vR}\gamma_{mn}(\vR)b_n=\frac12\gamma_{AB}
\sum_{n=\pm1}\sum_{l=\pm1}e^{ikal}b_n=\gamma_{AB}(e^{-ika}+e^{ika})C_B$ e $(E_k-E_A)C_A+\gamma_{AB}(e^{-ika}+e^{ika})C_B=0$ con $\gamma_{AB}=\mu$: con il
segno della (I) è $\gamma_{AB}=-\mu=t$ (come nella slide 20, $\mu=-\gamma_{AB}$), e per il fattore vale l'osservazione~\ref{oss:vicini}.

% =================================================================
\section{Il metodo $\vk\cdot\mathbf p$}\label{sec:kp}
% =================================================================
Un metodo semi-empirico per le bande vicino a un punto di alta simmetria, di solito $\Gamma$ \ingl{k dot p method}: note energie e funzioni d'onda in
$\vk=0$, le bande per $\vk$ piccolo si ottengono con la teoria delle perturbazioni. Dà soprattutto le masse efficaci (e le proprietà ottiche) e, con lo
spin--orbita, le tre bande di lacune dei semiconduttori (slide 33--48, «solo un breve riassunto»).

\begin{osservazione}[notazione]
Slide: $E_{n\vk}\to E_n(\vk)$, $E_{n0}\to E_n(0)$, $u_{n0}\equiv u_{n,\vk=0}$, $m$ o $m_o\to m_0$, $E_G\to\Eg$, $\mathcal H'\to H'_{\vk}$,
$\mathcal H_{\mathrm{so}}\to H_{\mathrm{so}}$, indici $i,j$ delle bande (slide 40) $\to n,n'$. $P$ è il parametro di Kane della slide 41 (un'energia per
una lunghezza), da non confondere con il $P$ del Kronig--Penney della Parte~3; $\Delta$ è l'energia di separazione spin--orbita. Lo zero delle energie
è la cima della banda di valenza.
\end{osservazione}

\subsection{Enunciato e significato}
\begin{modello}[metodo $\vk\cdot\mathbf p$]\label{mod:kp}
La parte periodica $u_{n\vk}$ della funzione di Bloch $\psi_{n\vk}=e^{i\vk\cdot\vr}u_{n\vk}$ risolve l'equazione di Schrödinger efficace
\[
\Bigl[\frac{\hp^2}{2m_0}+V(\vr)+\frac{\hbar\vk\cdot\hp}{m_0}+\frac{\hbar^2k^2}{2m_0}\Bigr]u_{n\vk}=E_n(\vk)\,u_{n\vk},\qquad
u_{n\vk}(\vr+\vR)=u_{n\vk}(\vr),\qquad \hp=-i\hbar\nabla .
\]
Con $H_0=\hp^2/2m_0+V$, esatta in $\Gamma$, e $H'_{\vk}=\hbar\vk\cdot\hp/m_0+\hbar^2k^2/2m_0$, piccola per $\vk$ piccolo, le bande vicino a $\Gamma$
si ottengono da $E_n(0)$ e $u_{n0}$ con la teoria delle perturbazioni.
\tcblower
Un metodo molto empirico: non serve $V(\vr)$ (potenziale vero o pseudopotenziale), bastano pochi parametri misurati in $\Gamma$ e adattati alle bande:
il gap, gli elementi di matrice di $\hp$ tra le bande («ottici»), lo spin--orbita. Vicino ai punti di alta simmetria le bande sono quasi parabole, e la
loro curvatura, la massa efficace, è fissata dall'accoppiamento $\vk\cdot\hp$ con le altre bande: conta di più la banda più vicina.
\end{modello}


\subsection{Dimostrazione principale (I): l'equazione per $u_{n\vk}$}\label{sec:kpu}
Slide 34--37; i passi 3--5 sono già nella Parte~3 (§8.3) per l'equazione di $u_{\vk}$.
\begin{dimostrazione}[equazione di Schrödinger efficace per $u_{n\vk}$]
\passo{Equazione di Schrödinger a un elettrone (eq.~(1)).}
$n$ è l'indice di banda:
\[
\Bigl[-\frac{\hbar^2\nabla^2}{2m_0}+V(\vr)\Bigr]\psi_{n\vk}(\vr)=E_n(\vk)\,\psi_{n\vk}(\vr).
\]
\passo{Forma di Bloch (eq.~(2)).}
\[
\psi_{n\vk}(\vr)=e^{i\vk\cdot\vr}u_{n\vk}(\vr),\qquad u_{n\vk}(\vr+\vR)=u_{n\vk}(\vr).
\]
\passo{Applichiamo l'impulso $\hp=-i\hbar\nabla$ (slide 35).}
Regola del prodotto:
\[
\hp\,\psi_{n\vk}=\hbar\vk\,\psi_{n\vk}+e^{i\vk\cdot\vr}\,\frac\hbar i\nabla u_{n\vk}(\vr)=e^{i\vk\cdot\vr}\,\hbar\Bigl(\vk+\frac1i\nabla\Bigr)u_{n\vk}(\vr).
\]
\passo{Applichiamo di nuovo: l'energia cinetica.}
Il fattore $e^{i\vk\cdot\vr}$ esce allo stesso modo:
\[
\frac{\hp^2}{2m_0}\,\psi_{n\vk}=e^{i\vk\cdot\vr}\,\frac{\hbar^2}{2m_0}\Bigl(\frac1i\nabla+\vk\Bigr)^2u_{n\vk}.
\]
\passo{Semplifichiamo $e^{i\vk\cdot\vr}$ (slide 35).}
\[
H_{\vk}u_{n\vk}\equiv\Bigl[\frac{\hbar^2}{2m_0}\Bigl(\frac1i\nabla+\vk\Bigr)^2+V(\vr)\Bigr]u_{n\vk}(\vr)=E_n(\vk)\,u_{n\vk}(\vr).
\]
\passo{Sviluppiamo il quadrato (slide 36).}
$\vk$ è un parametro costante, quindi $\vk\cdot\hp=\hp\cdot\vk$:
\[
\frac{\hbar^2}{2m_0}\Bigl(\frac1i\nabla+\vk\Bigr)^2=\frac{\hp^2}{2m_0}+\frac{\hbar^2k^2}{2m_0}+\frac{\hbar}{2m_0}\,2\,\vk\cdot\hp,\qquad
\vk\cdot\hp=k_x\Bigl(-i\hbar\frac{\partial}{\partial x}\Bigr)+k_y\Bigl(-i\hbar\frac{\partial}{\partial y}\Bigr)+k_z\Bigl(-i\hbar\frac{\partial}{\partial z}\Bigr).
\]
\passo{Conclusione (slide 34 e 37).}
\[
\Bigl[\frac{\hp^2}{2m_0}+\frac{\hbar\vk\cdot\hp}{m_0}+\frac{\hbar^2k^2}{2m_0}+V(\vr)\Bigr]u_{n\vk}(\vr)=E_n(\vk)\,u_{n\vk}(\vr).\qed
\]
\end{dimostrazione}
In forma esplicita (slide 38): $\bigl[-\frac{\hbar^2\nabla^2}{2m_0}+V(\vr)+\frac{\hbar^2k^2}{2m_0}-i\frac{\hbar^2}{m_0}\vk\cdot\nabla\bigr]u_{n\vk}(\vr)
=E_n(\vk)\,u_{n\vk}(\vr)$.

\begin{interpretazione}[non sono elettroni liberi (slide 37)]
$\hp$ è un operatore che agisce su $u_{n\vk}$, $\vk$ un numero: $\hp\ne\hbar\vk$, e i due non sono legati in modo semplice (se lo fossero, l'equazione
non avrebbe senso). Vicino a $\Gamma$ i termini $\hbar\vk\cdot\hp/m_0$ e $\hbar^2k^2/2m_0$ sono «piccoli» e si trattano con la teoria delle
perturbazioni.
\end{interpretazione}

\subsection{Dimostrazione principale (II): bande e massa efficace vicino a $\Gamma$}\label{sec:kpmassa}
La slide 40 dà già pronte le formule della teoria delle perturbazioni; prima le ricaviamo, con la notazione delle QP, poi le applichiamo.
\begin{approfondimento}[teoria delle perturbazioni non degenere fino al secondo ordine (\QP{§13, pp.~167--169})]
\passo{Impostazione.}
$H_0$ ha spettro noto e non degenere; il parametro $\tau$ serve solo a contare l'ordine ($\tau\to1$ alla fine):
\begin{gather*}
H=H_0+\tau V,\qquad H_0|n\rangle=E^0_n|n\rangle,\\
|E_n\rangle=|n\rangle+\tau|b_n\rangle+\tau^2|c_n\rangle+\dots,\qquad E_n=E^0_n+\tau E^1_n+\tau^2E^2_n+\dots
\end{gather*}
\passo{Sostituiamo in $H|E_n\rangle=E_n|E_n\rangle$ e uguagliamo le potenze di $\tau$.}
\[
\tau^1:\ \ \bigl(H_0-E^0_n\bigr)|b_n\rangle=-V|n\rangle+E^1_n|n\rangle,\qquad
\tau^2:\ \ \bigl(H_0-E^0_n\bigr)|c_n\rangle=-V|b_n\rangle+E^1_n|b_n\rangle+E^2_n|n\rangle .
\]
\passo{Scegliamo le correzioni ortogonali a $|n\rangle$.}
Il contributo di $|n\rangle$ è già nel primo termine: $\langle n|b_n\rangle=\langle n|c_n\rangle=0$, cioè $|b_n\rangle=\sum_{k\ne n}B_k(n)\,|k\rangle$.
\passo{Proiettiamo su $\langle n|$: le correzioni all'energia.}
$\langle n|(H_0-E^0_n)=0$ perché $H_0$ è hermitiana; restano
\[
E^1_n=\langle n|V|n\rangle,\qquad E^2_n=\langle n|V|b_n\rangle .
\]
\passo{Proiettiamo l'ordine $\tau^1$ su $\langle k|$, $k\ne n$: l'autovettore al primo ordine.}
\[
\bigl(E^0_k-E^0_n\bigr)B_k(n)=-\langle k|V|n\rangle\qquad\Longrightarrow\qquad B_k(n)=-\frac{\langle k|V|n\rangle}{E^0_k-E^0_n}.
\]
\passo{Sostituiamo in $E^2_n$.}
$V$ è hermitiana, $\langle n|V|k\rangle=\langle k|V|n\rangle^*$:
\[
E^2_n=\sum_{k\ne n}\langle n|V|k\rangle\,B_k(n)=\sum_{k\ne n}\frac{|\langle k|V|n\rangle|^2}{E^0_n-E^0_k}.
\]
\passo{Risultato ($\tau=1$): energia al secondo ordine, autovettore al primo.}
\[
E_n\simeq E^0_n+\langle n|V|n\rangle+\sum_{k\ne n}\frac{|\langle k|V|n\rangle|^2}{E^0_n-E^0_k},\qquad
|E_n\rangle\simeq|n\rangle+\sum_{k\ne n}\frac{\langle k|V|n\rangle}{E^0_n-E^0_k}\,|k\rangle .
\]
Per $E^2_n$ basta $|b_n\rangle$: la correzione $|c_n\rangle$ non entra. Nel $\vk\cdot\mathbf p$: $|n\rangle\to u_{n0}$, $E^0_n\to E_n(0)$, $\tau V\to H'_{\vk}$,
$k\to n'$.\qed
\end{approfondimento}
\begin{dimostrazione}[massa efficace al secondo ordine (slide 39--40 e 48)]
\passo{Problema imperturbato: $\vk=0$ (slide 39).}
Si risolve prima
\[
H_0u_{n0}=\Bigl[\frac{\hp^2}{2m_0}+V(\vr)\Bigr]u_{n0}=E_n(0)\,u_{n0}:
\]
$H_0$ è l'Hamiltoniana esatta in $\Gamma$, e le $u_{n0}$, autofunzioni della stessa $H_0$ hermitiana, formano un insieme completo e ortonormale di
funzioni periodiche (\QP{§4.1, Teorema~2, p.~56}; relazione di completezza in \QP{§3.1, p.~41}).
\passo{Perturbazione (slide 40).}
Sempre più piccola per $\vk\to0$: la teoria è tanto più accurata quanto più $k$ è piccolo.
\[
H=H_0+H'_{\vk},\qquad H'_{\vk}=\frac{\hbar^2k^2}{2m_0}+\frac{\hbar\vk\cdot\hp}{m_0}.
\]
\passo{Formule della teoria delle perturbazioni per una banda $n$ non degenere (slide 40; approfondimento precedente).}
\begin{gather*}
E_n(\vk)=E_n(0)+\langle u_{n0}|H'_{\vk}|u_{n0}\rangle+\sum_{n'\ne n}\frac{|\langle u_{n0}|H'_{\vk}|u_{n'0}\rangle|^2}{E_n(0)-E_{n'}(0)},\\
u_{n\vk}=u_{n0}+\sum_{n'\ne n}\frac{\langle u_{n'0}|H'_{\vk}|u_{n0}\rangle}{E_n(0)-E_{n'}(0)}\,u_{n'0}.
\end{gather*}
\passo{Primo ordine.}
Il secondo termine è $\hbar\vk\cdot\mathbf v_n(0)$, nullo in un estremo della banda (in $\Gamma$ per la simmetria di inversione):
\[
\langle u_{n0}|H'_{\vk}|u_{n0}\rangle=\frac{\hbar^2k^2}{2m_0}+\frac\hbar{m_0}\,\vk\cdot\langle u_{n0}|\hp|u_{n0}\rangle .
\]
\passo{Secondo ordine: nei termini misti conta solo $\vk\cdot\hp$.}
$\hbar^2k^2/2m_0$ è un numero e $\langle u_{n0}|u_{n'0}\rangle=0$ per $n'\ne n$:
\[
E_n(\vk)\simeq E_n(0)+\frac{\hbar^2k^2}{2m_0}+\frac{\hbar^2}{m_0^2}\sum_{n'\ne n}\frac{|\langle u_{n0}|\vk\cdot\hp|u_{n'0}\rangle|^2}{E_n(0)-E_{n'}(0)} .
\]
\passo{Conclusione: confronto con la parabola $E_n(0)+\hbar^2k^2/2m^*$ (slide 48).}
Dividendo per $\hbar^2k^2/2$:
\[
\frac1{m^*}=\frac1{m_0}+\frac2{m_0^2k^2}\sum_{n'\ne n}\frac{|\langle u_{n0}|\vk\cdot\hp|u_{n'0}\rangle|^2}{E_n(0)-E_{n'}(0)} .\qed
\]
\end{dimostrazione}

\begin{chiave}[Massa efficace dal metodo $\vk\cdot\mathbf p$]
\[
\boxed{\ E_n(\vk)\simeq E_n(0)+\frac{\hbar^2k^2}{2m^*},\qquad
\frac1{m^*}=\frac1{m_0}+\frac2{m_0^2k^2}\sum_{n'\ne n}\frac{|\langle u_{n0}|\vk\cdot\hp|u_{n'0}\rangle|^2}{E_n(0)-E_{n'}(0)}\ }
\]
Risultato del secondo ordine: le bande $n'$ più vicine a $n$ (denominatore piccolo) decidono la massa efficace di $n$. In generale $m^*$ dipende dalla
direzione di $\vk$ (tensore, Parte~3, §13.2).
\end{chiave}

\begin{osservazione}[slide 40]
Nella formula per $u_{i\vk}$ la somma contiene $u_{i0}$: va letto $u_{j0}$. Inoltre $\mathcal H_{\mathrm{so}}$ (§\ref{sec:so}) non dipende da $\vk$ e non
diventa piccolo per $\vk\to0$: va unito ad $H_0$ (si diagonalizza $H_0+H_{\mathrm{so}}$ in $\Gamma$), e la perturbazione resta
$\hbar\vk\cdot\hp/m_0+\hbar^2k^2/2m_0$.
\end{osservazione}

\begin{dimostrazione}[verifica: la massa al bordo dei quasi liberi]
Attorno a un punto $\vk_0\ne0$ vale la stessa formula, con $u_{n\vk_0}$ al posto di $u_{n0}$ e gli elementi di matrice
$\langle\psi_{n\vk_0}|\hp|\psi_{n'\vk_0}\rangle$. Due bande dei quasi liberi al bordo $k_0=\pi/a$ (Parte~3, §11): $\psi_c\propto\cos(\pi x/a)$ e
$\psi_s\propto\sin(\pi x/a)$, separate da $\Eg=2|V_G|$.
\passo{Elemento di matrice.}
$\hp\,\psi_s=-i\hbar\frac\pi a\,\psi_c$, quindi $|\langle\psi_c|\hp|\psi_s\rangle|=\hbar\pi/a$.
\passo{Formula con le sole due bande, $E_\pm-E_\mp=\pm\Eg$.}
\[
\frac{m_0}{m^*_\pm}=1\pm\frac2{m_0}\,\frac{(\hbar\pi/a)^2}{\Eg}=1\pm\frac{4E_0}{\Eg},\qquad E_0=\frac{\hbar^2\pi^2}{2m_0a^2}.
\]
\passo{Conclusione.}
$m^*_\pm=m_0\Eg/(\Eg\pm4E_0)$: la massa al bordo dei quasi liberi e del Kronig--Penney (Parte~3, §13.3).\qed
\end{dimostrazione}
Con tutte le bande la somma riproduce la curvatura esatta: per $a=\SI{3}{\angstrom}$ e $V_G=\SI{-0.5}{eV}$ si ha $m_0/m^*_+=17{,}731$ sia dalla
curvatura numerica sia dalla somma, e $17{,}72$ con le sole due bande (verificato con Python).

\subsection{Conseguenze: bande vicine, spin--orbita, lacune}\label{sec:so}
\domanda{Perché le bande vicine contano? Il segno dei contributi}
Una banda $n'$ \emph{sopra} $n$ ($E_{n'}>E_n$) dà un contributo negativo a $1/m^*$, una banda \emph{sotto} un contributo positivo. Per il fondo della
conduzione la banda più vicina è la valenza, più in basso di $\Eg$; tenendo solo quella,
\[
\frac{m_0}{m^*_c}\simeq1+\frac{E_P}{\Eg},\qquad E_P\equiv\frac2{m_0}\bigl|\langle u_{c0}|\hat p_x|u_{v0}\rangle\bigr|^2 ,
\]
e con $E_P\approx\SI{20}{eV}\gg\Eg$ (valore tipico, simile nei semiconduttori comuni) $m^*_c\simeq m_0\Eg/E_P\ll m_0$: \textbf{gap piccolo, massa
piccola}. È la tendenza $m^*\propto\Eg$ della tabella della Parte~3 (§13.3), da InSb a GaAs. Per la cima della valenza la conduzione sta sopra:
$m_0/m^*\simeq1-E_P/\Eg<0$, massa negativa, cioè lacune leggere.

\domanda{Che cosa aggiunge lo spin--orbita? (slide 38, 42, 44)}
L'interazione spin--orbita, che deriva dall'equazione di Dirac (qui senza dimostrazione; per un potenziale centrale è il termine
$H_{SL}=\frac1{2m_0^2c^2}\frac1r\frac{\dd U}{\dd r}\,\mathbf S\cdot\mathbf L$ di \QP{§12.2, p.~161}, con $U$ energia potenziale), è (slide 38, eq.~6.7--6.8)
\[
H_{\mathrm{so}}=-\frac{\hbar}{4m_0^2c^2}\,\boldsymbol\sigma\cdot\hp\times(\nabla V),\qquad
\sigma_x=\begin{pmatrix}0&1\\1&0\end{pmatrix},\ \ \sigma_y=\begin{pmatrix}0&-i\\i&0\end{pmatrix},\ \ \sigma_z=\begin{pmatrix}1&0\\0&-1\end{pmatrix};
\]
aggiunta all'Hamiltoniana $\vk\cdot\mathbf p$ dà (eq.~6.9)
\[
\Bigl[\frac{\hp^2}{2m_0}+V+\frac{\hbar^2k^2}{2m_0}+\frac\hbar{m_0}\vk\cdot\boldsymbol\pi+\frac{\hbar}{4m_0^2c^2}\,\hp\cdot\boldsymbol\sigma\times\nabla V\Bigr]|n\vk\rangle
=E_n(\vk)\,|n\vk\rangle,\qquad\boldsymbol\pi\equiv\hp+\frac\hbar{4m_0c^2}\,\boldsymbol\sigma\times\nabla V .
\]
Vicino ai nuclei ha la forma atomica $H_{\mathrm{so}}=\lambda\,\mathbf L\cdot\mathbf S$ (slide 42). L'effetto più importante è alla cima della valenza in
$\Gamma=(0,0,0)$: senza $H_{\mathrm{so}}$ gli stati sono di tipo $p$ ($L=1$), tre volte degeneri (sei volte con lo spin). Con $\mathbf J=\mathbf L+\mathbf S$,
$J^2=L^2+S^2+2\,\mathbf L\cdot\mathbf S$ (\QP{§11.2, p.~150}), e $L=1$, $S=\frac12$:
\[
\lambda\,\mathbf L\cdot\mathbf S=\frac\lambda2\bigl(J^2-L^2-S^2\bigr)=\frac{\lambda\hbar^2}2\Bigl[J(J+1)-2-\frac34\Bigr]=
\begin{cases}+\frac12\lambda\hbar^2 & J=\frac32\ \ (4\ \text{stati})\\[2pt] -\lambda\hbar^2 & J=\frac12\ \ (2\ \text{stati})\end{cases}
\]
mentre per $L=0$ è $\mathbf L\cdot\mathbf S=0$. $H_{\mathrm{so}}$ separa quindi in parte la degenerazione (figura~\ref{fig:kp}a): in alto i quattro stati
$J=\frac32$, le \textbf{lacune pesanti} \ingl{heavy holes}, $m_J=\pm\frac32$, e le \textbf{lacune leggere} \ingl{light holes}, $m_J=\pm\frac12$, ancora
degeneri in $\Gamma$; più in basso di
\[
\Delta=\tfrac32\lambda\hbar^2
\]
i due stati $J=\frac12$, la banda \textbf{split-off} \ingl{spin-orbit split-off}. Gli stati $s$ ($L=0$) non si separano.

\begin{figure}[tbp]
\centering
\figura{kp}
\caption{(a) Livelli in $\Gamma$: senza spin--orbita la cima della valenza $p$ è sei volte degenere; $\lambda\mathbf L\cdot\mathbf S$ la separa in $J=\frac32$
(lacune pesanti e leggere) e $J=\frac12$ (split-off), più in basso di $\Delta$; la conduzione $s$ non si separa. (b) Bande vicino a $\Gamma$ dalle formule
della slide 41 con parametri di tipo GaAs ($\Eg=\SI{1.52}{eV}$, $\Delta=\SI{0.34}{eV}$, $E_P=\SI{20}{eV}$): conduzione $c$, lacune pesanti $hh$,
leggere $lh$ (tratteggiata), split-off $so$ (schema della slide 43).}\label{fig:kp}
\end{figure}

\domanda{Che carattere hanno le funzioni d'onda ai bordi delle bande? (slide 44)}
La parte della funzione d'onda nella cella centrale decide la natura delle transizioni ottiche nei dispositivi optoelettronici. In un materiale a gap
diretto il minimo della conduzione ($\vk=0$) è di tipo $s$; in uno a gap indiretto (minimo in $X$) è una miscela $s+p$ (longitudinale) e $p$
(trasversale). La cima della valenza è fatta di stati $p$: con lo spin il momento angolare totale è $\frac32\hbar$, e con gli orbitali $|p_x\rangle$,
$|p_y\rangle$, $|p_z\rangle$:
\begin{align*}
\Phi_{\frac32,\frac32}&=-\tfrac1{\sqrt2}\bigl(|p_x\rangle+i|p_y\rangle\bigr){\uparrow}, &
\Phi_{\frac32,-\frac32}&=-\tfrac1{\sqrt2}\bigl(|p_x\rangle-i|p_y\rangle\bigr){\downarrow},\\
\Phi_{\frac32,\frac12}&=-\tfrac1{\sqrt6}\bigl[(|p_x\rangle+i|p_y\rangle){\downarrow}-2|p_z\rangle{\uparrow}\bigr], &
\Phi_{\frac32,-\frac12}&=\tfrac1{\sqrt6}\bigl[(|p_x\rangle-i|p_y\rangle){\uparrow}+2|p_z\rangle{\downarrow}\bigr]:
\end{align*}
nella prima riga le lacune pesanti ($m_J=\pm\frac32$), nella seconda le leggere ($m_J=\pm\frac12$).
I due stati split-off, ortogonali a questi, sono fatti degli stessi orbitali $p$ (Clebsch--Gordan, non nelle slide: gli stati $|j=\ell\pm\frac12,m_j\rangle$
di \QP{§11.2, p.~151} con $\ell=1$; verificati con Python):
$\Phi_{\frac12,\frac12}=-\frac1{\sqrt3}\bigl[(|p_x\rangle+i|p_y\rangle){\downarrow}+|p_z\rangle{\uparrow}\bigr]$ e
$\Phi_{\frac12,-\frac12}=-\frac1{\sqrt3}\bigl[(|p_x\rangle-i|p_y\rangle){\uparrow}-|p_z\rangle{\downarrow}\bigr]$.

\begin{osservazione}[slide 44]
L'etichetta $L=0$ accanto alla banda split-off, con le frecce verso $|s\rangle{\uparrow}$ e $|s\rangle{\downarrow}$, va letta $J=\frac12$: la split-off
nasce dagli orbitali $p$ ($L=1$), con $J=L-S$; si comporta come un doppietto di spin $\frac12$, come gli stati $s$. Il valore
$\lambda\mathbf L\cdot\mathbf S=\lambda$ «se $L=1$» dipende da $J$: $+\frac12\lambda\hbar^2$ per $J=\frac32$ e $-\lambda\hbar^2$ per $J=\frac12$; è nullo
per $L=0$, come scritto.
\end{osservazione}

\domanda{Come sono le bande vicino a $\Gamma$ in un semiconduttore a gap diretto? (slide 41, 43)}
Per un materiale a gap diretto come GaAs, con la conduzione $s$ e le tre valenze $p$ accoppiate da $\vk\cdot\hp$ e lo spin--orbita, il metodo dà
(Kittel, cap.~8; $P$, $\Eg$ e $\Delta$ sono parametri da adattare):
\begin{gather*}
E_c=\Eg+\frac{\hbar^2k^2}{2m_0}+\frac{P^2k^2}3\Bigl(\frac2\Eg+\frac1{\Eg+\Delta}\Bigr)\qquad(\text{conduzione più bassa}),\\
E_{hh}=-\frac{\hbar^2k^2}{2m_0},\qquad
E_{lh}=-\frac{\hbar^2k^2}{2m_0}-\frac{2P^2k^2}{3\Eg},\qquad
E_{so}=-\Delta-\frac{\hbar^2k^2}{2m_0}-\frac{P^2k^2}{3(\Eg+\Delta)}\\
(\text{le tre valenze più alte: lacune pesanti, leggere e split-off}),
\end{gather*}
con $P=\frac\hbar{m_0}|\langle S|\hat p_x|X\rangle|$, l'elemento di matrice «ottico» tra la conduzione $s$ e le valenze $p$. Sono parabole; con
$E_P=2m_0P^2/\hbar^2$ le masse (delle lacune, per le valenze) sono
\[
\frac{m_0}{m_c}=1+\frac{E_P}3\Bigl(\frac2\Eg+\frac1{\Eg+\Delta}\Bigr),\qquad
\frac{m_0}{m_{lh}}=1+\frac{2E_P}{3\Eg},\qquad
\frac{m_0}{m_{so}}=1+\frac{E_P}{3(\Eg+\Delta)},\qquad m_{hh}=m_0 .
\]
Le lacune pesanti non si accoppiano alla conduzione e restano «pesanti»; leggere e split-off si accoppiano e sono leggere come gli elettroni
(figura~\ref{fig:kp}b, schema della slide 43).

\begin{osservazione}[segno nella slide 41]
Nella formula delle lacune leggere manca il segno meno tra i due termini: $E_{lh}=-\frac{\hbar^2k^2}{2m_0}-\frac{2P^2k^2}{3\Eg}$.
\end{osservazione}

\begin{approfondimento}[da dove vengono le formule: un modello a due bande]
Senza spin--orbita e con $\vk$ lungo $z$. In $\Gamma$ ci sono la conduzione $|S\rangle$ (energia $\Eg$) e le valenze $|X\rangle$, $|Y\rangle$, $|Z\rangle$
(energia $0$); per simmetria $\langle S|\hat p_z|Z\rangle$ è l'unico elemento di $\hat p_z$ non nullo tra questi stati.
\passo{$\vk\cdot\hp=k\hat p_z$ accoppia $|S\rangle$ solo con $|Z\rangle$.}
Nella base $\{|S\rangle,|Z\rangle\}$, con $Pk=\frac{\hbar k}{m_0}\langle S|\hat p_z|Z\rangle$ (reale con una scelta di fase):
\[
H=\begin{pmatrix}\Eg&Pk\\Pk&0\end{pmatrix}+\frac{\hbar^2k^2}{2m_0}.
\]
\passo{Autovalori.}
\[
E=\frac{\hbar^2k^2}{2m_0}+\frac\Eg2\pm\sqrt{\frac{\Eg^2}4+P^2k^2}.
\]
\passo{Sviluppo per $Pk\ll\Eg$.}
$|X\rangle$ e $|Y\rangle$ non sono accoppiati e restano a $E=\hbar^2k^2/2m_0$, mentre
\[
E_c\simeq\Eg+\frac{\hbar^2k^2}{2m_0}+\frac{P^2k^2}\Eg,\qquad E_{lh}\simeq\frac{\hbar^2k^2}{2m_0}-\frac{P^2k^2}\Eg .
\]
\passo{Conclusione.}
Elettroni e lacune leggere hanno $|m^*|\simeq m_0\Eg/E_P$, piccola; le lacune pesanti non si accoppiano alla conduzione. Con lo spin--orbita
l'accoppiamento si divide tra lacune leggere ($\frac23$) e split-off ($\frac13$), e per $\Delta\to0$ la $E_c$ della slide 41 diventa
$\Eg+\hbar^2k^2/2m_0+P^2k^2/\Eg$, come qui.
\end{approfondimento}
Qui il termine libero delle valenze è $+\hbar^2k^2/2m_0$; la slide 41 (dal Kittel) scrive $-\hbar^2k^2/2m_0$. Per lacune leggere e split-off la
differenza conta poco ($E_P\gg\Eg$); per le pesanti la curvatura viene dalle bande lontane, e $m_{hh}=m_0$ è solo un ordine di grandezza
(misurata $0{,}3$--$0{,}5\,m_0$, tabella sotto).

\begin{esempio}[masse efficaci di GaAs e InSb]
Con $E_P=\SI{20}{eV}$: GaAs ($\Eg=\SI{1.43}{eV}$, $\Delta=\SI{0.34}{eV}$) dà $m_c=0{,}071\,m_0$ e $m_{lh}=0{,}097\,m_0$; InSb ($\Eg=\SI{0.18}{eV}$,
$\Delta=\SI{0.82}{eV}$) dà $m_c=0{,}012\,m_0$. Misurate: $0{,}067$ e $0{,}0145$ (tabella della Parte~3, §13.3), e $m_{lh}=0{,}082$ per il GaAs. Un
solo parametro, quasi uguale per tutti, spiega masse che variano di un fattore 5 seguendo il gap.
\end{esempio}

\begin{center}\small
\begin{tabular}{@{}lcccc@{}}
\toprule
& $m_{hh}/m_0$ & $m_{lh}/m_0$ & $m_{so}/m_0$ & $\Delta$ (eV)\\
\midrule
Si & 0,52 & 0,16 & & 0,044\\
Ge & 0,34 & 0,043 & & 0,29\\
GaAs & 0,5 & 0,082 & 0,17 & 0,34\\
InSb & 0,39 & 0,021 & 0,11 & 0,82\\
\bottomrule
\end{tabular}\\[3pt]
{\footnotesize Masse delle lacune e separazione spin--orbita misurate (Kittel, cap.~8).}
\end{center}

\domanda{Che cosa mostrano le bande calcolate di GaAs, Ge e Si? (slide 45--47)}
\begin{itemize}[leftmargin=*]
\item \textbf{GaAs} (slide 45): gap diretto in $\Gamma$, circa \SI{1.5}{eV}; in $\Gamma$ lacune pesanti e leggere degeneri e, più in basso
di $\Delta\simeq\SI{0.34}{eV}$, la split-off.
\item \textbf{Ge} (slide 46, calcolo di C.\,Y.~Fong): gap indiretto, con il minimo della conduzione in $L=\frac{2\pi}a\bigl(\frac12,\frac12,\frac12\bigr)$
(\SI{0.9}{eV} nel calcolo; misurato \SI{0.66}{eV} a \SI{300}{K}) e superfici di energia costante ellissoidali attorno a $L$; in $\Gamma$ la
conduzione è a circa \SI{1.1}{eV}; la struttura fine della cima della valenza è lo spin--orbita, $\Delta=\SI{0.29}{eV}$.
\item \textbf{Si} (slide 47): gap indiretto, con il minimo della conduzione («electrons») lungo $\Gamma X$, vicino a $X$, e la cima della valenza
(«holes») in $\Gamma$; la figura è senza spin--orbita (cima tre volte degenere, $\Gamma_{25'}$): nel Si la split-off è solo \SI{0.044}{eV} sotto.
\end{itemize}
Il metodo attorno a $\Gamma$ descrive la cima della valenza di tutti e tre; per la conduzione di Si e Ge si sviluppa attorno al minimo ($X$ o $L$).

\domanda{A che cosa serve il metodo? (slide 48)}
(A) Dà bande ragionevoli vicino ai punti di simmetria della zona con una parametrizzazione e un calcolo semplici («con una calcolatrice tascabile»).
(B) Dà masse efficaci ragionevolmente accurate (chiave del §\ref{sec:kpmassa}). Con abbastanza termini nello sviluppo può essere accurato in tutta la
zona di Brillouin (slide 40).

\begin{infobox}[Curiosità di laboratorio: lacune pesanti nei qubit di germanio]
Nei pozzi quantici di Ge compresso tra strati di SiGe, la piattaforma dei qubit di lacuna di Delft (Veldhorst, Scappucci), deformazione e
confinamento separano lacune pesanti e leggere di oltre \SI{100}{meV}: le lacune dei qubit hanno carattere di lacuna pesante ($J_z=\pm\frac32$). Per
un'«inversione di massa» la lacuna pesante, pesante lungo la direzione di crescita, nel piano del pozzo ha massa efficace molto piccola, circa
$0{,}05\,m_0$ ($0{,}061\,m_0$ misurata in dispositivi non drogati): i livelli orbitali dei punti quantici sono ben separati anche in punti grandi (circa
\SI{100}{nm}), con requisiti di litografia meno severi (Scappucci et al., «The germanium quantum information route», 2021).
\end{infobox}

% =================================================================
\section{Il metodo dello pseudopotenziale}\label{sec:pseudo}
% =================================================================
Le bande vere di Si, Ge e GaAs si ottengono con gli elettroni quasi liberi, ma con un potenziale efficace debole e liscio, lo
\textbf{pseudopotenziale} \ingl{pseudopotential}, al posto del potenziale vero, e senza trattarlo come una perturbazione (slide 1). Si ricava dalle
onde piane ortogonalizzate (OPW) e in pratica si adatta ai dati (Lezione 5e, slide 1--38; Ashcroft--Mermin, cap.~11; Yu--Cardona, cap.~2).

\begin{osservazione}[notazione]
Le slide scrivono $U(\vr)$ o $V(\vr)$ per il potenziale cristallino; $\psi^c_{\vk n}$, $\psi^c_n$ o $|\Phi_{j\vk}\rangle$ per gli stati di core e
$E_n$, $E(\vk)_n$ o $E_j$ per le loro energie; $\psi^O_{\vk}$, $\phi_{\vk}$ o $\chi_i$ per l'onda piana ortogonalizzata, con coefficienti
$\beta_n(\vk)$, $b_c$ o $\mu_{ij}$; $E_{\vk}$ o $\varepsilon(\vk)$ per la banda di valenza; $V'$, $V^R$ o $V_B$ per il potenziale repulsivo,
$V_{\mathrm{ps}}$ o $W$ per lo pseudopotenziale, $H'$ per la pseudo-Hamiltoniana; nell'esercizio (slide 35--38) $\hat p$ è il proiettore,
$\varphi$ la pseudo-funzione, $C_i$ i suoi coefficienti, $\vk_i$ i vettori reciproci, $\mu_{ij}=-\beta$, $W'$ ed $E'$ la scelta alternativa; unità
gaussiane. Qui: $V(\vr)$ potenziale vero e $H=\hp^2/2m_0+V$; stati di core $\psi_{c\vk}$ di energia $E_c$ ($c$ conta i livelli
di core); OPW $\chi_{\vk}$ con coefficienti $\beta_c(\vk)$; stato di valenza $\psi_{\vk}$ di energia $E(\vk)$ ($E_{\vk}$ resta la parabola libera);
parte liscia $\psi^{\mathrm f}_{\vk}$; $V'$ e $V_{\mathrm{ps}}=V+V'$; pseudo-Hamiltoniana $H_{\mathrm{ps}}$ ($H'_{\vk}$ è la perturbazione del
§\ref{sec:kp}); proiettore sul core $\hat\Pi_c$ ($\hp$ è l'impulso); $c_{\vG}$, $\vG_i$, $\tilde V_{\mathrm{ps}}$ e $\tilde E$ nell'esercizio; $Z$ carica del nucleo e $Z_{\mathrm{ion}}$ dello ione (Parte~3, §6.1); unità SI
($e^2\to e^2/4\pi\varepsilon_0$).
\end{osservazione}

\subsection{Perché il potenziale vero non è «quasi libero»}\label{sec:pseudomotivi}
Ashcroft--Mermin (cap.~11, slide 3) parte dai limiti dei due metodi visti: i legami forti, nella forma semplice, descrivono bene solo le bande dei
livelli di core; gli elettroni quasi liberi non si applicano direttamente a nessun solido reale. Tre premesse sul modello a elettroni indipendenti:
\begin{itemize}[leftmargin=*]
\item[\#1] scrivere per ogni elettrone l'equazione (11.1), $\bigl[-\frac{\hbar^2}{2m_0}\nabla^2+V(\vr)\bigr]\psi_{\vk}=E(\vk)\,\psi_{\vk}$, semplifica
già molto il problema di tanti elettroni interagenti, che in un trattamento esatto non hanno ciascuno una propria funzione d'onda;
\item[\#2] l'approssimazione a elettroni indipendenti non trascura le interazioni: le include con una scelta opportuna di $V$, che contiene il potenziale
degli ioni \emph{e} gli effetti periodici dell'interazione con gli altri elettroni;
\item[\#3] quest'ultima dipende dalle funzioni d'onda degli altri elettroni, soluzioni della stessa (11.1): per conoscere $V$ bisogna conoscere le
soluzioni, e il problema diventa autoconsistente. Lo si vedrà con lo pseudopotenziale, che dà risultati vicini ai quasi liberi.
\end{itemize}
I metodi di calcolo puntano alle bande più alte, le \textbf{bande di valenza} \ingl{valence bands}, piene, in parte piene o vuote, che decidono il
comportamento elettronico del solido; gli elettroni di core sono inerti per molti scopi (Parte~3, §6.1, con l'esempio del Si, slide 8--9:
core $1s^22s^22p^6$, valenza $3s^23p^2$, ibridi $sp^3$ con i 4 primi vicini del diamante). Lo pseudopotenziale cerca insieme una funzione d'onda migliore
e un potenziale efficace (slide 4).

\domanda{Come sono fatte le funzioni d'onda di core e di valenza? (slide 4--5)}
Le funzioni di core sono apprezzabili solo vicino allo ione, dove hanno le oscillazioni tipiche degli orbitali atomici (alta energia cinetica che, con
la grande energia potenziale negativa, dà l'energia del livello); tra gli ioni sono quasi nulle, e l'inviluppo è la $\mathrm{Re}\,e^{i\vk\cdot\vr}$ di
Bloch. Le funzioni di valenza hanno le stesse oscillazioni atomiche nel core, ma tra gli ioni non sono piccole: variano lentamente, come onde piane
(figura 11.2 di Ashcroft--Mermin). Poiché la velocità è $\frac\hbar{im_0}\nabla$, più rapida è la variazione, maggiore è la velocità dell'elettrone in
quella regione.

\begin{dimostrazione}[perché la funzione di valenza oscilla nel core (slide 6)]
\passo{Argomento energetico (\#1).}
Nel core un elettrone di valenza sente la stessa grande energia potenziale negativa di un elettrone di core, ma ha energia totale più alta: la sua
energia cinetica è ancora maggiore, quindi nel core oscilla \emph{più} delle funzioni di core.
\passo{Valenza e core sono autostati della stessa $H$ hermitiana con autovalori diversi (\#2).}
Le autofunzioni di un operatore hermitiano con autovalori diversi sono ortogonali (\QP{§4.1, Teorema~2, p.~56}):
\[
0=\int\dd\vr\;\psi_{c\vk}^*(\vr)\,\psi_{\vk}(\vr)\qquad\text{per ogni livello di core }c\qquad\text{(eq.~11.5)}.
\]
\passo{Basta una cella.}
Per il teorema di Bloch ($\psi_{n\vk}=e^{i\vk\cdot\vr}u_{n\vk}$, eq.~8.3) l'integrando $\psi^*_{c\vk}\psi_{\vk}=u^*_{c\vk}u_{\vk}$ è lo stesso in ogni cella:
l'integrale è $N$ volte quello su una cella.
\passo{Conclusione: solo la regione del core di un solo ione.}
Le funzioni di core sono apprezzabili solo vicino allo ione: perché la (11.5) valga per \emph{tutti} i livelli di core, nel core $\psi_{\vk}$ deve avere
oscillazioni che si intrecciano con quelle di ciascun $\psi_{c\vk}$ e ne cancellano l'integrale.\qed
\end{dimostrazione}

\domanda{Perché allora poche onde piane non bastano? (slide 7, 26)}
Una funzione che oscilla sulla scala del core ha nello sviluppo di Fourier (eq.~9.1 di Ashcroft--Mermin, l'equazione centrale della Parte~3, §11.3)
\[
\psi_{\vk}(\vr)=\sum_{\vG}c_{\vG}\,e^{i(\vk+\vG)\cdot\vr}
\]
molti termini con $|\vG|$ grande: il metodo dei quasi liberi, con poche onde piane, è insostenibile. Ci si aspetta bande di valenza di carattere
\textbf{atomico vicino agli ioni} e di \textbf{onda piana nella regione interstiziale} \ingl{interstitial region}: tutti i metodi in uso cercano di
riprodurre la struttura atomica nel core tenendo conto che i livelli di valenza non sono di tipo a legami forti.
Inoltre il potenziale vero non è piccolo. Nel core $V(\vr)\simeq-Ze^2/4\pi\varepsilon_0r$ (eq.~11.2), e il suo coefficiente di Fourier (eq.~11.3--11.4,
con $\Omega$ volume per atomo e $a_0$ raggio di Bohr) è
\[
V_{\vG}\simeq-\frac{Ze^2}{\varepsilon_0\Omega G^2},\qquad
|V_{\vG}|\simeq\frac{e^2}{8\pi\varepsilon_0a_0}\Bigl(\frac{a_0^3}\Omega\Bigr)\frac{8\pi Z}{(a_0G)^2},\qquad\frac{e^2}{8\pi\varepsilon_0a_0}=\SI{13.6}{eV}:
\]
alcuni eV per moltissimi $\vG$, quanto le energie cinetiche dell'equazione centrale. L'ipotesi $|V_{\vG}|\ll E_{\vk+\vG}$ non è lecita
(è piccolo solo $V_{\mathrm{ps}}$, §\ref{sec:pseudoconseg}).
\begin{esempio}[il Si]
$a=\SI{5.431}{\angstrom}$, $\Omega=a^3/8=\SI{20.0}{\angstrom^3}$, $Z=14$, $|\vG_{111}|=2\pi\sqrt3/a=\SI{2.00}{\angstrom^{-1}}$: $|V_{\vG}|\simeq\SI{31}{eV}$
($\SI{12}{eV}$ e \SI{9}{eV} per $|\vG|^2=8$ e $11$ in unità $(2\pi/a)^2$), contro un'energia cinetica libera di \SI{3.8}{eV} in $\vG_{111}/2$.
Il fattore di forma dello pseudopotenziale per lo stesso $\vG$ è \SI{-2.9}{eV} (tabella del §\ref{sec:pseudoconseg}): dieci volte più piccolo
(verificato con Python).
\end{esempio}

\domanda{Che cosa chiede Yu--Cardona? (slide 10)}
Perché un'approssimazione che parte dagli elettroni quasi liberi dovrebbe funzionare per elettroni di valenza legati fortemente nei legami? La
risposta è nel §\ref{sec:pseudoconseg}.

\subsection{Le onde piane ortogonalizzate}\label{sec:opw}
Per dare a $\psi_{\vk}$ il doppio carattere, oscillante nel core e di onda piana tra gli ioni, Herring propose di non usare onde piane semplici ma
onde piane già ortogonalizzate ai livelli di core (slide 11--13). Approssimazione «zero»: elettroni liberi, $\psi^{\mathrm f}_{\vk}=e^{i\vk\cdot\vr}$
(f per \emph{free}); approssimazione successiva, «quasi liberi»: $\psi_{\vk}$ simile a un'onda piana ma, per la regola della slide 6, ortogonale a tutti
gli stati di core.

\begin{modello}[onde piane ortogonalizzate, OPW]\label{mod:opw}
Gli stati di core sono noti, $H\psi_{c\vk}=E_c\psi_{c\vk}$, normalizzati ($\int|\psi_{c\vk}|^2\dd\vr=1$) e in pratica combinazioni a legami forti dei
livelli atomici. L'\textbf{onda piana ortogonalizzata} \ingl{orthogonalized plane wave, OPW} è (eq.~11.24)
\[
\chi_{\vk}(\vr)=\underbrace{e^{i\vk\cdot\vr}}_{\text{interstizi}}+\underbrace{\sum_c\beta_c(\vk)\,\psi_{c\vk}(\vr)}_{\text{core}},\qquad
\int\chi_{\vk}^*\,\psi_{c\vk}\,\dd\vr=0\ \ \text{per ogni }c,
\]
con la somma su tutti i livelli di core. Gli stati di valenza si sviluppano in OPW:
$\psi_{\vk}=\sum_{\vG}c_{\vG}\,\chi_{\vk+\vG}$.
\tcblower
All'onda piana si tolgono le sue componenti sugli stati di core. Nel core $\chi_{\vk}$ oscilla come le funzioni di core che contiene, tra gli ioni il
secondo termine è piccolo e $\chi_{\vk}\simeq e^{i\vk\cdot\vr}$: ha già il carattere di una funzione di valenza, e ne servono poche dove servivano
moltissime onde piane.
\end{modello}

\begin{dimostrazione}[i coefficienti $\beta_c(\vk)$ (slide 15 e 17)]
\passo{L'OPW.}
\[
\chi_{\vk}(\vr)=e^{i\vk\cdot\vr}+\sum_{c'}\beta_{c'}(\vk)\,\psi_{c'\vk}(\vr).
\]
\passo{Imponiamo l'ortogonalità al livello di core $c$.}
\[
\int\dd\vr\;\chi_{\vk}^*(\vr)\,\psi_{c\vk}(\vr)=0 .
\]
\passo{Sostituiamo il complesso coniugato dell'OPW.}
\[
\int e^{-i\vk\cdot\vr}\,\psi_{c\vk}\,\dd\vr+\sum_{c'}\beta_{c'}^*(\vk)\int\psi^*_{c'\vk}\,\psi_{c\vk}\,\dd\vr=0 .
\]
\passo{Gli stati di core sono ortonormali.}
Livelli diversi sono ortogonali per lo stesso teorema della slide 6 (\QP{§4.1, Teorema~2, p.~56}), e ciascuno è normalizzato:
$\int\psi^*_{c'\vk}\psi_{c\vk}\,\dd\vr=\delta_{c'c}$.
\passo{Conclusione.}
\[
\beta_c^*(\vk)=-\int\dd\vr\;e^{-i\vk\cdot\vr}\,\psi_{c\vk}(\vr),\qquad\text{cioè}\qquad
\beta_c(\vk)=-\int\psi^*_{c\vk}(\vr)\,e^{i\vk\cdot\vr}\,\dd\vr=-\langle\psi_{c\vk}|e^{i\vk\cdot\vr}\rangle .\qed
\]
\end{dimostrazione}
\formula{\chi_{\vk}=e^{i\vk\cdot\vr}-\sum_c\langle\psi_{c\vk}|e^{i\vk\cdot\vr}\rangle\,\psi_{c\vk}}
L'ultima forma è in notazione di Dirac, con $\langle\psi|\phi\rangle=\int\psi^*\phi\,\dd\vr$ (\QP{§10.2, pp.~128--130}).

\domanda{L'OPW è una funzione di Bloch? (slide 13 e 15)}
Sì, con lo stesso $\vk$.
\begin{dimostrazione}[$\chi_{\vk}$ soddisfa la condizione di Bloch]
\passo{L'onda piana la soddisfa banalmente.}
\[
e^{i\vk\cdot(\vr+\vR)}=e^{i\vk\cdot\vR}\,e^{i\vk\cdot\vr}.
\]
\passo{La somma di core è una funzione di Bloch, come nei legami forti.}
Le $\psi_{c\vk}$ hanno vettore $\vk$, quindi
\[
\sum_c\beta_c(\vk)\,\psi_{c\vk}(\vr)=e^{i\vk\cdot\vr}\,u^c_{\vk}(\vr),\qquad u^c_{\vk}(\vr+\vR)=u^c_{\vk}(\vr).
\]
\passo{Conclusione.}
\[
\chi_{\vk}(\vr)=e^{i\vk\cdot\vr}\bigl[1+u^c_{\vk}(\vr)\bigr]:
\]
una funzione di Bloch di parte periodica $1+u^c_{\vk}$.\qed
\end{dimostrazione}
Per la stessa ragione $\chi_{\vk}$ è ortogonale anche agli stati di core $\psi_{c\vk'}$ con $\vk'\ne\vk$ (nota 26 di Ashcroft--Mermin), e la
combinazione $\sum_{\vG}c_{\vG}\chi_{\vk+\vG}$ ha ancora la forma di Bloch.

\begin{osservazione}[slide 11]
La nota «questa condizione definisce un insieme completo di funzioni» va letta così: l'ortogonalità fissa i $\beta_c$; le OPW $\chi_{\vk+\vG}$ di tutti i
$\vG$, insieme agli stati di core, formano un insieme completo, perché le onde piane lo sono (\QP{§3.2, p.~44}) e l'ortogonalizzazione toglie solo la
loro parte di core. Le OPW con $\vG$ diversi non sono ortogonali tra loro.
\end{osservazione}

\subsection{Dimostrazione principale: l'equazione di pseudo-Schrödinger}\label{sec:pseudodim}
L'OPW suggerisce di separare lo stato di valenza in una parte liscia e una parte di core, e di scrivere un'equazione per la sola parte liscia
(slide 18--22). Al posto di $e^{i\vk\cdot\vr}$ si mette una funzione più generale $\psi^{\mathrm f}_{\vk}$, per esempio una combinazione di onde piane.

\begin{teorema}[equazione di pseudo-Schrödinger]\label{teo:pseudo}
Sia lo stato di valenza
\[
\psi_{\vk}(\vr)=\psi^{\mathrm f}_{\vk}(\vr)+\sum_c\beta_c(\vk)\,\psi_{c\vk}(\vr),\qquad \beta_c(\vk)=-\langle\psi_{c\vk}|\psi^{\mathrm f}_{\vk}\rangle,
\]
ortogonale a tutti gli stati di core. Allora l'equazione vera $H\psi_{\vk}=E(\vk)\psi_{\vk}$ equivale all'equazione per la sola parte liscia
\[
\Bigl[\frac{\hp^2}{2m_0}+V_{\mathrm{ps}}\Bigr]\psi^{\mathrm f}_{\vk}=E(\vk)\,\psi^{\mathrm f}_{\vk},\qquad V_{\mathrm{ps}}=V+V',\qquad
V'=\sum_c\bigl[E(\vk)-E_c\bigr]\,|\psi_{c\vk}\rangle\langle\psi_{c\vk}| .
\]
\tcblower
Le energie sono le stesse dell'equazione vera, ma la funzione d'onda è solo la parte liscia, quella che serve tra gli atomi: le grandi oscillazioni
nel core, necessarie per l'ortogonalità, sono tutte raccolte nel potenziale efficace $V'$. Poiché $E(\vk)>E_c$, $V'$ è repulsivo e compensa in buona
parte il potenziale attrattivo $V$: lo \textbf{pseudopotenziale} $V_{\mathrm{ps}}$ è piccolo e liscio.
\end{teorema}

\begin{dimostrazione}[dall'equazione vera all'equazione per $\psi^{\mathrm f}_{\vk}$ (slide 20)]
\passo{Equazione di Schrödinger del solido (1).}
$V$ è il potenziale cristallino vero:
\[
H\psi_{\vk}=E(\vk)\,\psi_{\vk},\qquad H=\frac{\hp^2}{2m_0}+V(\vr).
\]
\passo{Equazione degli stati di core (2).}
Le soluzioni, di tipo legami forti, sono note:
\[
H\psi_{c\vk}=E_c\,\psi_{c\vk}.
\]
\passo{Sostituiamo $\psi_{\vk}=\psi^{\mathrm f}_{\vk}+\sum_c\beta_c\psi_{c\vk}$ nella (1).}
$H$ è lineare:
\[
H\psi^{\mathrm f}_{\vk}+\sum_c\beta_c\,H\psi_{c\vk}=E\,\psi^{\mathrm f}_{\vk}+E\sum_c\beta_c\,\psi_{c\vk}.
\]
\passo{Usiamo la (2).}
\[
H\psi^{\mathrm f}_{\vk}+\sum_c\beta_c\,E_c\,\psi_{c\vk}=E\,\psi^{\mathrm f}_{\vk}+E\sum_c\beta_c\,\psi_{c\vk}.
\]
\passo{Portiamo i termini di core a sinistra.}
\[
H\psi^{\mathrm f}_{\vk}-\sum_c\beta_c\,(E-E_c)\,\psi_{c\vk}=E\,\psi^{\mathrm f}_{\vk}.
\]
\passo{Definiamo l'operatore $V'$ (3).}
È una definizione operatoriale: dice che cosa fa $V'$ quando agisce su $\psi^{\mathrm f}_{\vk}$.
\[
V'\psi^{\mathrm f}_{\vk}\equiv-\sum_c\bigl(E-E_c\bigr)\beta_c(\vk)\,\psi_{c\vk}\qquad\Longrightarrow\qquad(H+V')\,\psi^{\mathrm f}_{\vk}=E\,\psi^{\mathrm f}_{\vk}.
\]
\passo{Usiamo l'ortogonalità al core.}
Come per l'OPW, con $\psi^{\mathrm f}_{\vk}$ al posto di $e^{i\vk\cdot\vr}$: $0=\langle\psi_{c\vk}|\psi_{\vk}\rangle=\langle\psi_{c\vk}|\psi^{\mathrm f}_{\vk}\rangle+\beta_c$.
\passo{Conclusione.}
\[
V'\psi^{\mathrm f}_{\vk}=\sum_c\bigl(E-E_c\bigr)\langle\psi_{c\vk}|\psi^{\mathrm f}_{\vk}\rangle\,\psi_{c\vk},\qquad\text{cioè}\qquad
V'=\sum_c\bigl(E-E_c\bigr)|\psi_{c\vk}\rangle\langle\psi_{c\vk}| .\qed
\]
\end{dimostrazione}

\begin{chiave}[Equazione di pseudo-Schrödinger]
\[
\boxed{\ H_{\mathrm{ps}}\,\psi^{\mathrm f}_{\vk}=\Bigl[-\frac{\hbar^2\nabla^2}{2m_0}+V_{\mathrm{ps}}(\vr)\Bigr]\psi^{\mathrm f}_{\vk}=E(\vk)\,\psi^{\mathrm f}_{\vk},\qquad
V_{\mathrm{ps}}=V+\sum_c\bigl(E(\vk)-E_c\bigr)|\psi_{c\vk}\rangle\langle\psi_{c\vk}|\ }
\]
La (1) è equivalente alla (3), cioè alla (4) $H_{\mathrm{ps}}\psi^{\mathrm f}_{\vk}=E\psi^{\mathrm f}_{\vk}$: stessa banda $E(\vk)$. Lo stato vero è
$\psi_{\vk}=\psi^{\mathrm f}_{\vk}+{}$correzioni di core $\sum_c\beta_c\psi_{c\vk}$.
\end{chiave}

\begin{osservazione}[slide 19]
La slide deduce «$V'$ positivo» dalla condizione sui $\beta^*_n$; il motivo preciso è la forma con il proiettore: per ogni funzione $\phi$
\[
\langle\phi|V'|\phi\rangle=\sum_c\bigl(E-E_c\bigr)\,|\langle\psi_{c\vk}|\phi\rangle|^2\ \ge0\qquad(E>E_c),
\]
cioè $V'$ è un operatore positivo (repulsivo), non una funzione positiva di $\vr$: è \emph{non locale} (agisce proiettando sul core) e dipende
dall'energia. La slide avverte anche che Ashcroft--Mermin (pp.~208--209) usa un'altra notazione e ha un errore nella formula (11.34).
\end{osservazione}

\subsection{Conseguenze: pseudopotenziale debole, pochi coefficienti, bande vere}\label{sec:pseudoconseg}
\domanda{Che cos'è $\psi^{\mathrm f}_{\vk}$? (slide 19)}
La parte «liscia» di $\psi_{\vk}$, quella necessaria tra gli atomi; la parte $\sum_c(E-E_c)\beta_c\psi_{c\vk}$, con le grandi oscillazioni vicino agli
atomi richieste dall'ortogonalità, è conglobata nel potenziale efficace $V'$. L'equazione (3) contiene solo $\psi^{\mathrm f}_{\vk}$.

\domanda{Perché $V_{\mathrm{ps}}$ è piccolo? (slide 19, 24, 38)}
$E(\vk)>E_c$, perché la valenza sta sopra il core: $V'$ è repulsivo (osservazione precedente) e compensa in parte il potenziale vero $V$, attrattivo e
negativo. $V_{\mathrm{ps}}=V+V'$ è piccolo (\emph{schermatura} del potenziale). Nel Si il fattore di forma per $\vG_{111}$ scende da circa
\SI{31}{eV} (ione nudo, esempio del §\ref{sec:pseudomotivi}) a \SI{2.9}{eV}.

\domanda{Come si usa in pratica? (slide 22--23)}
In linea di principio si potrebbe usare l'espressione formale $V_{\mathrm{ps}}=V+V'$, liscia e piccola, con la somma sugli stati di core: \emph{non si fa
quasi mai}. Si sceglie invece fra:
\begin{enumerate}[leftmargin=*, label=\arabic*.]
\item \textbf{pseudopotenziale empirico} \ingl{empirical pseudopotential method}: $V_{\mathrm{ps}}$ si scrive con pochi parametri, adattati alle bande
e ad altre proprietà misurate;
\item \textbf{pseudopotenziale autoconsistente} \ingl{self-consistent pseudopotential method}: l'equazione di pseudo-Schrödinger
$[-\hbar^2\nabla^2/2m_0+V_{\mathrm{ps}}]\psi^{\mathrm f}_{\vk}=E\psi^{\mathrm f}_{\vk}$ si accoppia all'equazione di Poisson per la carica degli
elettroni, $\nabla^2V_{\mathrm{ps}}=-4\pi\rho=-4\pi e|\psi^{\mathrm f}_{\vk}|^2$ (unità gaussiane, slide 23), e si risolvono insieme.
\end{enumerate}
\begin{osservazione}[l'equazione di Poisson della slide 23]
È una scrittura schematica: in SI il potenziale elettrostatico $\phi$ soddisfa $\nabla^2\phi=-\rho/\varepsilon_0$, l'energia potenziale
dell'elettrone è $-e\phi$, e $\rho$ contiene la carica degli ioni e quella di \emph{tutti} gli stati di valenza occupati, $-e\sum_{\text{occ}}|\psi|^2$.
È la premessa \#3: il potenziale dipende dalle soluzioni.
\end{osservazione}

\domanda{Perché bastano pochi coefficienti di Fourier? (slide 24, 29--30)}
Lo ione è fortemente schermato: vicino al nucleo $V_{\mathrm{ps}}$ non è più una «sfera dura» coulombiana, ma tende dolcemente a zero
(pseudopotenziale \emph{soft core}) o a una costante (\emph{hard core}) entro un raggio $r_c$ di circa metà della lunghezza di legame, e fuori
coincide con il potenziale dello ione (figura~\ref{fig:pseudo}a). La sua trasformata di Fourier ha \textbf{pochi coefficienti non nulli}: è per questo
che gli elettroni quasi liberi davano bande «buone» (slide 24). Con i coefficienti sui vettori reciproci (Parte~3, §11.2)
\[
V_{\mathrm{ps}}(\vr)=\sum_{\vG}V_{\vG}\,e^{i\vG\cdot\vr},\qquad V_{\vG}=\frac1{\Omega_{\mathrm{cella}}}\int_{\text{cella}}V_{\mathrm{ps}}(\vr)\,e^{-i\vG\cdot\vr}\,\dd\vr ,
\]
e per un solo ione la trasformata continua $V_{\mathrm{ps}}(\vr)\sim\int v(q)\,e^{i\vq\cdot\vr}\,\dd\vq$: $v(q)$ parte da $-\frac23\EF$ per $q\to0$
(ione schermato, sotto), risale con l'andamento $\propto-1/q^2$ dello ione (slide 29), cambia segno e tende a zero (figura~\ref{fig:pseudo}b). Nel diamante e nella zincoblenda i $\vG$
più corti hanno $|\vG|^2=3,4,8,11$ in unità $(2\pi/a)^2$, e il fattore di struttura dei due atomi in $\pm\boldsymbol\tau$,
$\boldsymbol\tau=\frac a8(1,1,1)$ (Parte~1, §14), lascia i fattori di forma \ingl{form factors} simmetrici $V^s_3$, $V^s_8$, $V^s_{11}$ (moltiplicati
per $\cos\vG\cdot\boldsymbol\tau$, nullo per $|\vG|^2=4$) e, nei composti, gli antisimmetrici $V^a_3$, $V^a_4$, $V^a_{11}$ (per
$\sin\vG\cdot\boldsymbol\tau$, nullo per $|\vG|^2=8$). Sono i parametri del metodo empirico (slide 30, in Ry $=\SI{13.6}{eV}$):
\begin{center}\small
\begin{tabular}{@{}lcccccccc@{}}
\toprule
& Si & Ge & GaAs & GaP & InAs & InSb & ZnSe & CdTe\\
\midrule
$V^s_3$ & $-0{,}211$ & $-0{,}269$ & $-0{,}252$ & $-0{,}249$ & $-0{,}27$ & $-0{,}25$ & $-0{,}23$ & $-0{,}245$\\
$V^s_8$ & $0{,}04$ & $0{,}038$ & $0$ & $0{,}017$ & $0{,}02$ & $0{,}01$ & $0{,}01$ & $-0{,}015$\\
$V^s_{11}$ & $0{,}08$ & $0{,}035$ & $0{,}08$ & $0{,}083$ & $0{,}041$ & $0{,}044$ & $0{,}06$ & $0{,}073$\\
\bottomrule
\end{tabular}\\[3pt]
{\footnotesize Fattori di forma simmetrici (Yu--Cardona); la slide non riporta le colonne antisimmetriche, il cui segno dipende dalla scelta di
anione e catione.}
\end{center}

\begin{figure}[tbp]
\centering
\figura{pseudo}
\caption{(a) Potenziale dello ione nudo ($\propto-1/r$) e pseudopotenziale: dentro il core, di raggio $r_c$ (circa metà della lunghezza di
legame), tende a zero (\emph{soft core}) o a una costante (\emph{hard core}, tratteggiato); fuori coincide con lo ione. (b) Fattore di forma
$v(q)$ di uno ione schermato (curva di modello, «empty core» di Ashcroft con la schermatura di Thomas--Fermi, parametri di tipo Si): parte da
$-\frac23\EF$, cambia segno e tende a zero; nei $\vG$ del diamante ($|\vG|^2=3$, $8$, $11$): $V_3<0$ e $V_8$, $V_{11}$ piccoli e positivi, come nella
tabella (slide 24 e 29).}\label{fig:pseudo}
\end{figure}

\begin{osservazione}[slide 24]
«$V_3$, $V_8$, $V_9$ vengono dal fattore di struttura» va letto $V_3$, $V_8$, $V_{11}$ ($|\vG|^2=9$ non è un vettore dell'FCC); «$V_8$, $V_{11}\sim0$»
vuol dire piccoli rispetto a $V_3$, non nulli (tabella).
\end{osservazione}

\domanda{Quanto vale lo pseudopotenziale per $q\to0$? (slide 25, metalli)}
Lo ione nudo diverge per $q\to0$ come il potenziale coulombiano; la schermatura elimina la divergenza e lascia un valore finito, fissato dalla densità
di stati al livello di Fermi.
\begin{dimostrazione}[$v(0)=-\frac23\EF$]
\passo{Ione nudo di carica $Z_{\mathrm{ion}}e$.}
Per $q$ piccolo la trasformata di Fourier è coulombiana ($\Omega$ volume per ione) e diverge:
\[
\vion(q)\to-\frac{Z_{\mathrm{ion}}e^2}{\varepsilon_0\Omega q^2}\qquad\Bigl(\text{gaussiane: }-\frac{4\pi Z_{\mathrm{ion}}e^2}{\Omega q^2}\Bigr).
\]
\passo{Schermatura del gas di elettroni.}
Il gas risponde e riduce il potenziale, $v(q)=\vion(q)/\varepsilon(q)$, con la costante dielettrica di Thomas--Fermi. Per $q\ll\kF$ la funzione di
Lindhard coincide con quella di Thomas--Fermi, quindi il limite è esatto:
\[
\varepsilon(q)=1+\frac{k_{\mathrm{TF}}^2}{q^2}\ \xrightarrow{q\to0}\ \frac{k_{\mathrm{TF}}^2}{q^2}.
\]
\passo{Il fattore $q^2$ si semplifica.}
\[
v(0)=-\frac{Z_{\mathrm{ion}}e^2}{\varepsilon_0\Omega\,k_{\mathrm{TF}}^2}.
\]
\passo{Vettore d'onda di Thomas--Fermi.}
$k_{\mathrm{TF}}^2=e^2g(\EF)/\varepsilon_0$, con $g(\EF)$ densità di stati per unità di volume al livello di Fermi; per il gas libero con
$n=Z_{\mathrm{ion}}/\Omega$ (Parte~2, §3), $g(\EF)=3n/2\EF$:
\[
k_{\mathrm{TF}}^2=\frac{3e^2Z_{\mathrm{ion}}}{2\varepsilon_0\Omega\EF}.
\]
\passo{Conclusione: $Z_{\mathrm{ion}}$, $\Omega$ ed $e^2$ si semplificano.}
\[
v(0)=-\frac{Z_{\mathrm{ion}}e^2}{\varepsilon_0\Omega}\cdot\frac{2\varepsilon_0\Omega\EF}{3e^2Z_{\mathrm{ion}}}=-\frac23\EF .\qed
\]
\end{dimostrazione}
\begin{interpretazione}[schermatura perfetta]
Per $q\to0$ la carica indotta, $-g(\EF)\,v$ elettroni per unità di volume, compensa esattamente la carica dello ione, $Z_{\mathrm{ion}}/\Omega$:
$v(0)=-Z_{\mathrm{ion}}/\Omega g(\EF)=-\frac23\EF$. Il potenziale sposta il livello di Fermi quanto basta perché nella regione dello ione entri la
carica che lo neutralizza. Il risultato non dipende dal modello di ione: nell'«empty core» di Ashcroft $\vion(q)=-Z_{\mathrm{ion}}e^2\cos(qR_c)/
\varepsilon_0\Omega q^2$, e $\cos(qR_c)\to1$.
\end{interpretazione}
Conseguenza: lo pseudopotenziale non diverge e in $q=0$ vale una frazione definita di $\EF$; nei metalli semplici è debole e la teoria perturbativa dei
quasi liberi funziona bene. Per esempio Na ($\EF=\SI{3.24}{eV}$): $v(0)=\SI{-2.16}{eV}$; Al ($\SI{11.7}{eV}$): \SI{-7.8}{eV} (verificato con Python).
\begin{osservazione}[slide 25]
«TF = temperatura di Fermi, $k$ costante di Boltzmann» va letto: TF = Thomas--Fermi, e $k_{\mathrm{TF}}$ è il vettore d'onda di schermatura di
Thomas--Fermi (l'inverso della lunghezza di schermatura). «Per elettrone» vuol dire che $v(0)$ è un'energia potenziale, il coefficiente di Fourier
per un elettrone.
\end{osservazione}

\domanda{Perché allora gli elettroni di valenza sono «quasi liberi»? (slide 10, 27--28)}
Risposta di Yu--Cardona: sono quasi liberi nel senso che gran parte della carica del nucleo è schermata dagli elettroni di core, legati molto
fortemente. Riassumendo: l'equazione di pseudo-Schrödinger dà la banda $E(\vk)$ cercata; di solito si assume $V_{\mathrm{ps}}$ debole (non è
indispensabile, ma giustifica l'approssimazione originale), e schematicamente
\[
\psi_{\vk}\simeq\psi^{\mathrm f}_{\vk}+\text{correzioni}\ \Bigl(\sum_c\beta_c\psi_{c\vk}\Bigr).
\]
Così si torna agli elettroni quasi liberi (Parte~3, §11), ma \textbf{con $V_{\mathrm{ps}}$ al posto del potenziale vero $V$}: l'equazione centrale con i
$V_{\vG}$ della tabella. Nel metodo empirico non la si tratta con la teoria delle perturbazioni (slide 1): si diagonalizza con abbastanza onde piane.

\domanda{Che cosa danno le bande calcolate? (slide 2, 31--34)}
\begin{itemize}[leftmargin=*]
\item \textbf{Si} (slide 2 e 33): bande con pseudopotenziale non locale (continue) e locale (tratteggiate), quasi uguali; cima della valenza in
$\Gamma$ ($\Gamma_{25'}$), minimo della conduzione lungo $\Gamma X$ vicino a $X$: gap \textbf{indiretto}.
\item \textbf{Ge} (slide 33): gap \textbf{indiretto}, con il minimo della conduzione in $L$; la cima della valenza in $\Gamma$ mostra lo spin--orbita.
\item \textbf{GaAs} e \textbf{ZnSe} (slide 34): gap \textbf{diretto} in $\Gamma$ (circa \SI{1.5}{eV} e \SI{2.7}{eV}).
\end{itemize}
Le masse efficaci della conduzione in $\Gamma$ si confrontano con le misure (slide 31; $\Eg(\Gamma)$ è il gap diretto in $\Gamma$, $E_0$ nella slide):
\begin{center}\small
\begin{tabular}{@{}lccccccc@{}}
\toprule
& Ge & GaAs & InP & InAs & GaSb & InSb & CdTe\\
\midrule
$\Eg(\Gamma)$ (eV) & 0,89 & 1,55 & 1,34 & 0,45 & 0,81 & 0,24 & 1,59\\
$m^*_c/m_0$ misurata & 0,041 & 0,067 & 0,073 & 0,026 & 0,047 & 0,015 & 0,11\\
$m^*_c/m_0$ calcolata & 0,04 & 0,078 & 0,067 & 0,023 & 0,04 & 0,012 & 0,08\\
\bottomrule
\end{tabular}
\end{center}
La riga «calcolata» (eq.~2.44 di Yu--Cardona) coincide, entro gli arrotondamenti, con $m^*_c/m_0\simeq\Eg(\Gamma)/E_P$, $E_P=\SI{20}{eV}$: è la formula a
due bande del §\ref{sec:so}, $m_0/m^*_c\simeq1+E_P/\Eg$, senza l'1 (verificato con Python). $E_P$ è quasi uguale per tutti perché le pseudo-funzioni
sono quasi onde piane: con $|\langle u_{c0}|\hat p_x|u_{v0}\rangle|\sim\hbar\,2\pi/a$ si stima $E_P\sim4\,\frac{\hbar^2}{2m_0}\bigl(\frac{2\pi}a\bigr)^2
\approx\SIrange{14}{20}{eV}$ per $a=\SIrange{6.5}{5.4}{\angstrom}$ (stima, non nelle slide).

\subsection{Altre dimostrazioni: l'esercizio delle slide 35--38, con il proiettore}\label{sec:pseudoes}
\textbf{Testo} (SUNY Buffalo). Il metodo delle onde piane ortogonalizzate, legato a quello dello pseudopotenziale, si usa spesso per le bande dei
metalli e spiega perché vale l'approssimazione degli elettroni quasi liberi anche se il potenziale vero non è affatto debole. Discuti le idee fisiche e
i passaggi matematici: (a) costruzione delle funzioni di Bloch degli stati di core e loro proprietà; (b) costruzione delle onde piane ortogonalizzate e
loro proprietà; (c) costruzione dello pseudopotenziale e sua non unicità; (d) perché lo pseudopotenziale si può trattare come piccolo.

\begin{dimostrazione}[pseudopotenziale con il proiettore sul core]
Fuori dai core l'onda è quasi piana: $|\vk+\vG_i\rangle=\frac1{\sqrt{N\Omega}}e^{i(\vk+\vG_i)\cdot\vr}$, con $N$ celle di volume $\Omega$ e $\vG_i$
vettore reciproco (\QP{§3.2, p.~44}: onde piane normalizzate in una scatola con condizioni periodiche); tutte hanno lo stesso $\vk$ nella zona ridotta.
\passo{(a) Stati di core a legami forti.}
Nel core gli elettroni sono fortemente legati; con $\varphi_c(\vr-\vR_l)$ orbitale di core dell'atomo in $\vR_l$ e la somma su tutti gli atomi:
\[
|\psi_{c\vk}\rangle=\frac1{\sqrt N}\sum_le^{i\vk\cdot\vR_l}\,\varphi_c(\vr-\vR_l),\qquad H|\psi_{c\vk}\rangle=E_c|\psi_{c\vk}\rangle .
\]
\passo{(b) OPW: togliamo all'onda piana le componenti di core.}
Come nel §\ref{sec:opw}, con $|\vk+\vG_i\rangle$ al posto di $e^{i\vk\cdot\vr}$:
\[
\chi_{\vk+\vG_i}=|\vk+\vG_i\rangle+\sum_c\beta_c(\vk+\vG_i)\,|\psi_{c\vk}\rangle,\qquad
\langle\psi_{c\vk}|\chi_{\vk+\vG_i}\rangle=0\ \Longrightarrow\ \beta_c(\vk+\vG_i)=-\langle\psi_{c\vk}|\vk+\vG_i\rangle ;
\]
lo stato di valenza è una combinazione di OPW, $\psi_{\vk}=\sum_ic_{\vG_i}\,\chi_{\vk+\vG_i}$.
\passo{(c) Il proiettore sugli stati di core.}
$\hat\Pi_c$ è la parte della relazione di completezza $\hat I=\sum_n|n\rangle\langle n|$ (\QP{§3.1, p.~41; §10.2, eq.~(115), p.~130}) ristretta al core:
\[
\hat\Pi_c=\sum_c|\psi_{c\vk}\rangle\langle\psi_{c\vk}|,\qquad\chi_{\vk+\vG_i}=\bigl(1-\hat\Pi_c\bigr)|\vk+\vG_i\rangle .
\]
\passo{Pseudo-funzione e stato vero.}
\[
\psi^{\mathrm f}_{\vk}\equiv\sum_ic_{\vG_i}\,|\vk+\vG_i\rangle\qquad\Longrightarrow\qquad\psi_{\vk}=\bigl(1-\hat\Pi_c\bigr)\psi^{\mathrm f}_{\vk}.
\]
\passo{Sostituiamo nell'equazione di Schrödinger (slide 37).}
\[
H\bigl(1-\hat\Pi_c\bigr)\psi^{\mathrm f}_{\vk}=E\bigl(1-\hat\Pi_c\bigr)\psi^{\mathrm f}_{\vk}\qquad\Longrightarrow\qquad
H\psi^{\mathrm f}_{\vk}-H\hat\Pi_c\psi^{\mathrm f}_{\vk}+E\hat\Pi_c\psi^{\mathrm f}_{\vk}=E\psi^{\mathrm f}_{\vk}.
\]
\passo{$H\hat\Pi_c$: gli stati di core sono autostati di $H$.}
\[
H\hat\Pi_c=\sum_cH|\psi_{c\vk}\rangle\langle\psi_{c\vk}|=\sum_cE_c\,|\psi_{c\vk}\rangle\langle\psi_{c\vk}| .
\]
\passo{Conclusione: lo pseudopotenziale.}
Con $H=-\frac{\hbar^2}{2m_0}\nabla^2+V$:
\[
-\frac{\hbar^2}{2m_0}\nabla^2\psi^{\mathrm f}_{\vk}+V_{\mathrm{ps}}\psi^{\mathrm f}_{\vk}=E\psi^{\mathrm f}_{\vk},\qquad
V_{\mathrm{ps}}=V-H\hat\Pi_c+E\hat\Pi_c=V+\sum_c(E-E_c)\,|\psi_{c\vk}\rangle\langle\psi_{c\vk}| ,
\]
lo stesso $V_{\mathrm{ps}}$ del Teorema~\ref{teo:pseudo}.\qed
\end{dimostrazione}

\domanda{(c) Perché lo pseudopotenziale non è unico?}
Con una funzione arbitraria $f(E,E_c)$ al posto di $E-E_c$ si ottiene la stessa energia.
\begin{dimostrazione}[non unicità: $\tilde E=E$]
\passo{Un'altra scelta.}
\[
\tilde V_{\mathrm{ps}}=V+\sum_cf(E,E_c)\,|\psi_{c\vk}\rangle\langle\psi_{c\vk}|,\qquad
-\frac{\hbar^2}{2m_0}\nabla^2\tilde\psi^{\mathrm f}+\tilde V_{\mathrm{ps}}\tilde\psi^{\mathrm f}=\tilde E\,\tilde\psi^{\mathrm f}.
\]
\passo{Prodotto scalare con lo stato vero $\psi_{\vk}$.}
$\psi_{\vk}$ è ortogonale al core, $\langle\psi_{\vk}|\psi_{c\vk}\rangle=0$: il termine con $f$ sparisce,
\[
\Bigl\langle\psi_{\vk}\Bigl|-\frac{\hbar^2}{2m_0}\nabla^2+V\Bigr|\tilde\psi^{\mathrm f}\Bigr\rangle=\langle\psi_{\vk}|H|\tilde\psi^{\mathrm f}\rangle
=\tilde E\,\langle\psi_{\vk}|\tilde\psi^{\mathrm f}\rangle .
\]
\passo{$H$ è hermitiana: agisce a sinistra su $\psi_{\vk}$.}
$\langle\psi_{\vk}|H\tilde\psi^{\mathrm f}\rangle=\langle H\psi_{\vk}|\tilde\psi^{\mathrm f}\rangle$ (\QP{§2.4, p.~24}), e $H\psi_{\vk}=E\psi_{\vk}$ con $E$
reale (\QP{§4.1, Teorema~1, p.~55}):
\[
E\,\langle\psi_{\vk}|\tilde\psi^{\mathrm f}\rangle=\tilde E\,\langle\psi_{\vk}|\tilde\psi^{\mathrm f}\rangle .
\]
\passo{Conclusione.}
Se $\langle\psi_{\vk}|\tilde\psi^{\mathrm f}\rangle\ne0$, $\tilde E=E$: anche $\tilde V_{\mathrm{ps}}$ è uno pseudopotenziale.\qed
\end{dimostrazione}
Anche la pseudo-funzione non è unica: $\psi^{\mathrm f}_{\vk}+\alpha\,\psi_{c\vk}$ dà lo stesso $\psi_{\vk}$, perché $(1-\hat\Pi_c)\psi_{c\vk}=0$. La
libertà si usa per rendere $V_{\mathrm{ps}}$ il più debole e liscio possibile. Equivalenza ($E$ autovalore di $H+\sum_c(E-E_c)|\psi_{c\vk}\rangle
\langle\psi_{c\vk}|$, $\psi_{\vk}=(1-\hat\Pi_c)\psi^{\mathrm f}_{\vk}$) e non unicità verificate con Python su un modello unidimensionale.

\domanda{(d) Perché lo pseudopotenziale è piccolo?}
$V_{\mathrm{ps}}=V+V'$, con $V'=\sum_c(E-E_c)|\psi_{c\vk}\rangle\langle\psi_{c\vk}|$. Nel cristallo l'energia $E$ di un elettrone di valenza è maggiore
di quella $E_c$ di un elettrone interno: $V'>0$ respinge gli elettroni, mentre il potenziale vero $V$ è attrattivo, cioè negativo. $V'$ tende a
cancellare $V$, e la cancellazione è di solito apprezzabile: $V_{\mathrm{ps}}$ è debole e abbastanza liscio.

\begin{osservazione}[slide 35--38]
Nella slide 36 la
scrittura $\psi_{\vk}=(1-\hat p)|\vk+\vk_i\rangle$, segnata come errore nella slide stessa, è l'OPW $\chi_{\vk+\vG_i}$; nella condizione di ortogonalità
va letto $\Phi^*_{j\vk}$. L'annotazione «$H\psi^{\mathrm f}_{\vk}=E(\vk)\psi^{\mathrm f}_{\vk}$» accanto all'equazione di Schrödinger va letta
$\bigl(-\frac{\hbar^2}{2m_0}\nabla^2+V_{\mathrm{ps}}\bigr)\psi^{\mathrm f}_{\vk}=E\psi^{\mathrm f}_{\vk}$: la pseudo-funzione non risolve l'equazione
con $H$. Nella dimostrazione a mano (slide 37) $|\vk-\vk_i\rangle$ è $|\vk+\vk_i\rangle$.
\end{osservazione}

% =================================================================
\section{Le bande dei materiali veri: metalli e semiconduttori}\label{sec:reali}
% =================================================================
Con i metodi dei §§\ref{sec:tb3d}--\ref{sec:pseudo} si leggono le bande calcolate e misurate dei solidi in tre dimensioni: l'alluminio, metallo
quasi libero, e Si e GaAs, semiconduttori con un gap assoluto. Contando gli stati di una banda e gli elettroni per cella si capisce quando un cristallo
è un metallo (Lezione 6a, prima parte, slide 1--23; Hofmann, cap.~6; Ashcroft--Mermin, cap.~8; Kittel, cap.~9). In questa sezione «slide» indica la
Lezione 6a.

\begin{osservazione}[notazione]
Le slide scrivono $\varepsilon_{\vk}$, $\varepsilon_{n,\vk'}$ o $E_{n\vk}$ per la banda; $U_G$, $U_{2G}$ per i coefficienti di Fourier del potenziale;
$g$ per il vettore reciproco in 1D ($k=k'\pm ng$) e $\vG'$ nella (6.41); $g(E)$ con il volume $V$ del campione; nella fotoemissione $E_f$, $\vk_f$
(stato finale), $E_i$, $\vk_i$ (iniziale), $h\nu$, $\mathbf K$ (elettrone nel vuoto), $V_0$ (fondo della banda finale). Qui: $E_n(\vk)$; $V_G$,
$V_{2G}$ come nella Parte~3, §11; $\vG$; $g(E)$ \emph{per unità di volume}, come nel §\ref{sec:pseudoconseg}; $E_{\mathrm{fin}}$, $\vk_{\mathrm{fin}}$
($E_f$ si confonde con $\EF$), $E_{\mathrm i}$, $\vk_{\mathrm i}$, $\hbar\omega$, $\mathbf K$ con le componenti $K_\parallel$, $K_\perp$ ($K$ da solo è il
punto della zona), $V_{\mathrm{int}}$ potenziale interno ($V_0$ è già usato), $\phi_{\mathrm e}$ lavoro di estrazione; $N$ celle primitive,
$n_{\mathrm{val}}$ elettroni di valenza per cella primitiva.
\end{osservazione}

\subsection{Tre modi di disegnare le bande}\label{sec:schemi}
Nei tre metodi (quasi liberi, legami forti, pseudopotenziale; slide 1) il potenziale periodico ha due effetti.
\begin{enumerate}[leftmargin=*, label=\arabic*.]
\item \textbf{Simmetria delle bande} (eq.~6.41; Parte~3, Teorema~9.1): basta la prima zona di Brillouin, perché le bande sono identiche attorno
agli altri punti del reticolo reciproco; al bordo della zona le bande si \textbf{ripiegano} \ingl{back-folding}.
\formula{E_n(\vk+\vG)=E_n(\vk)}
\item \textbf{Gap} dove le bande dell'elettrone libero erano degeneri.
\end{enumerate}

\begin{definizione}[i tre schemi di zona]
Lo stesso spettro si disegna in tre modi (slide 3; Kittel, cap.~9):
\begin{itemize}[leftmargin=*]
\item \textbf{zona estesa}: la banda $n$ sta nella zona $n$; $E(k)$ segue la parabola libera e salta di $2|V_G|$, $2|V_{2G}|$, \ldots{} ai bordi;
\item \textbf{zona ridotta}: tutte le bande nella prima zona; un $\vk$ fuori si riporta in $\vk'$ con una traslazione di $\pm n\vG$,
\[
\vk=\vk'+\vG,\qquad E_{n}(\vk')\ \ \text{con }\vk'\ \text{nella prima zona};
\]
\item \textbf{zona periodica}: la zona ridotta ripetuta con periodo $\vG$, comoda per seguire una banda attraverso il bordo.
\end{itemize}
\end{definizione}

\begin{figure}[tbp]
\centering
\figura{schemi}
\caption{Gli stessi elettroni quasi liberi in 1D nei tre schemi (slide 3; $V_G=0{,}45\,E_0$, $V_{2G}=0{,}3\,E_0$, $E_0=\hbar^2\pi^2/2m_0a^2$,
calcolati con 15 onde piane). Il gap tra le bande 1 e 2 si apre in $\pm\pi/a$; quello tra le bande 2 e 3 in $\pm2\pi/a$, cioè, ripiegato, in
$k=0$. In grigio la prima zona.}\label{fig:schemi}
\end{figure}

\subsection{L'alluminio: un metallo quasi libero}\label{sec:al}
L'alluminio ($1s^22s^22p^63s^23p^1$) ha un core di neon e \textbf{3 elettroni di valenza} $s$ e $p$; il reticolo è fcc con un atomo per cella
primitiva. La zona di Brillouin è l'ottaedro troncato (Parte~1, §13), con i punti di alta simmetria (unità $2\pi/a$):
\begin{center}\small
\begin{tabular}{@{}lcccccc@{}}
\toprule
& $\Gamma$ & $X$ & $L$ & $K$ & $W$\\
\midrule
coordinate & $(0,0,0)$ & $(1,0,0)$ & $(\h,\h,\h)$ & $(\frac34,\frac34,0)$ & $(1,\h,0)$\\
distanza da $\Gamma$ & 0 & 1 & 0,866 & 1,061 & 1,118\\
\bottomrule
\end{tabular}
\end{center}
$\Gamma$ è il centro della zona; $X$, $L$, $K$, $W$ stanno sul bordo.

\domanda{Che cosa mostra la banda lungo $\Gamma$--$X$--$\Gamma$? (slide 2, 4, 6)}
In zona periodica (Hofmann, fig.~6.12a): da $\Gamma$ parte la parabola libera; in $X$, sul bordo, si apre un gap (circa \SI{1.5}{eV}) tra la banda
$s$ sotto e la banda $p$ sopra; sopra il gap una banda torna verso $\Gamma$: è la parabola centrata nel $\Gamma$ della zona successiva,
$(2\pi/a)(2,0,0)$, come nel modello 1D. Unica differenza dal modello 1D: lo zero delle energie, che nei metalli si mette quasi sempre in $\EF$.

\domanda{E nelle altre direzioni? (slide 2, 5, 7)}
In zona ridotta lungo $\Gamma$--$X$--$W$--$K$--$L$--$\Gamma$ (fig.~6.12b) compaiono le bande ripiegate dalle zone vicine e i gap ai bordi. L'alluminio
è descritto molto bene dagli elettroni quasi liberi: \emph{sembra} complicato solo perché in 3D molte zone vicine si ripiegano nella prima. Lo si vede
confrontando la fig.~6.12b con le parabole libere ripiegate (fig.~\ref{fig:al}).

\begin{figure}[tbp]
\centering
\figura{al_libero}
\caption{Alluminio nel reticolo vuoto ($V=0$): lungo $\Gamma$--$X$--$W$--$K$--$L$--$\Gamma$, in zona ridotta, ogni curva è una parabola
$\hbar^2|\vk+\vG|^2/2m_0$ con un diverso $\vG$. Con $a=\SI{4.05}{\angstrom}$ e 3 elettroni per cella $\EF$ sta \SI{11.65}{eV} sopra il fondo. Le bande
vere (slide 5, fig.~6.12b) sono quasi uguali, con piccoli gap dove qui le curve si incrociano ($X$, $W$, $K$, $L$).}\label{fig:al}
\end{figure}

\domanda{Perché l'alluminio è un metallo? (slide 5--7)}
Le bande sono piene fino a $\EF$ e \textbf{più bande attraversano $\EF$}: un elettrone appena sotto $\EF$ passa a uno stato libero appena sopra, nella
stessa banda, con un'energia piccola a piacere, e così trasporta corrente elettrica e termica, in più direzioni. In 3D non c'è un gap assoluto:
$\EF$ taglia alcune bande.

\begin{esempio}[la sfera di Fermi contiene tutta la prima zona]
Slide 9: $\EF$ dista circa \SI{12}{eV} dal fondo della banda. Il gas libero (Parte~2, §3) con $n=3\cdot4/a^3=\SI{1.81e29}{m^{-3}}$ dà
$\kF=\SI{1.75}{\angstrom^{-1}}=1{,}127\,(2\pi/a)$ ed $\EF=\SI{11.65}{eV}$. La sfera di Fermi ha volume $\frac32$ di quello della zona (3 elettroni, 2
per stato) e $\kF>|\Gamma W|$, il punto più lontano della zona: la prima banda è tutta piena e la mezza zona che resta riempie in parte la seconda e la
terza banda. Per questo più bande tagliano $\EF$ (verificato con Python).
\end{esempio}

\domanda{E con un solo elettrone per cella? (slide 9, 15)}
$\EF$ scende di un fattore $3^{2/3}$: $\EF'=\SI{5.6}{eV}$ sopra il fondo, $\kF'=\kF/3^{1/3}=0{,}78\,(2\pi/a)$, minore della distanza del punto più
vicino del bordo, $|\Gamma L|=0{,}866\,(2\pi/a)$. La sfera di Fermi resta dentro la prima zona: le bande tagliano $\EF'$ alla stessa distanza da
$\Gamma$ in tutte le direzioni e la superficie di Fermi è una sfera, come nel gas di elettroni liberi. Sarebbe un metallo davvero libero, come il
sodio.

\begin{osservazione}[slide 2, 9 e 15: la linea tratteggiata]
Nella fig.~6.12b $\EF'$ è disegnata a \SI{-8}{eV}, circa \SI{3}{eV} sopra il fondo: corrisponderebbe a circa mezzo elettrone per cella. Con un
elettrone il gas libero dà \SI{5.6}{eV} sopra il fondo, cioè circa \SI{-6}{eV} (fig.~\ref{fig:al}); la linea è indicativa e la conclusione non cambia.
\end{osservazione}

\begin{osservazione}[slide 5 e 12: $\Gamma$ non è sul bordo]
«$\Gamma$, $X$, $W$, $K$, $L$ are all BZ boundary» (slide 5) e la nota della slide 12 («$L$, $\Gamma$, $X$, $W$, $K$ BZ zone boundaries»): $\Gamma$ è
il centro della zona; sul bordo stanno $X$, $W$, $K$, $L$, come scrive la slide 6.
\end{osservazione}

\subsection{Come si misurano le bande: la fotoemissione}\label{sec:arpes}
Le bande si misurano con la \textbf{fotoemissione risolta in angolo} \ingl{angle-resolved photoemission spectroscopy, ARPES} (slide 8). Il campione,
in ultra alto vuoto ($p<\SI{e-7}{mbar}$), è illuminato con fotoni ultravioletti monocromatici di energia $\hbar\omega$; per effetto fotoelettrico
emette elettroni, che un analizzatore separa per energia cinetica $E_{\mathrm{kin}}$ e direzione (angolo $\theta$ dalla normale). Dalle leggi di
conservazione si risale a energia e $\vk$ dentro il campione, cioè a $E(\vk)$ (la copertina di Hofmann è la versione misurata della fig.~6.12a).

\begin{risultato}[leggi di conservazione nella fotoemissione]
Con le energie misurate dal livello del vuoto ($E_{\mathrm{vuoto}}=0$, figura (a) della slide 8):
\[
E_{\mathrm{fin}}=E_{\mathrm i}+\hbar\omega,\qquad \vk_{\mathrm{fin}}=\vk_{\mathrm i},\qquad K_\parallel=k_{\mathrm{fin}\parallel},\qquad
K_\perp\ne k_{\mathrm{fin}\perp},
\]
con $\vk_{\mathrm i}$ ed $E_{\mathrm i}$ dello stato occupato nella banda, $\vk_{\mathrm{fin}}$ ed $E_{\mathrm{fin}}=E_{\mathrm{kin}}$ dello stato finale nel
solido, $\mathbf K$ dell'elettrone nel vuoto. La transizione è «verticale» (l'impulso del fotone è trascurabile); attraversando la superficie si
conserva solo la componente parallela, per la simmetria di traslazione nel piano.
\end{risultato}

\begin{dimostrazione}[$k_\parallel$ dall'angolo di emissione]
\passo{Nel vuoto l'elettrone è libero.}
\[
E_{\mathrm{kin}}=\frac{\hbar^2K^2}{2m_0}\qquad\Longrightarrow\qquad K=\frac{\sqrt{2m_0E_{\mathrm{kin}}}}\hbar .
\]
\passo{Componente parallela alla superficie.}
\[
K_\parallel=K\sin\theta=\frac{\sqrt{2m_0E_{\mathrm{kin}}}}\hbar\,\sin\theta .
\]
\passo{Conclusione: conservazione di $\vk$ nella transizione e di $K_\parallel$ alla superficie.}
\[
k_{\mathrm i\parallel}=k_{\mathrm{fin}\parallel}=K_\parallel=\frac{\sqrt{2m_0E_{\mathrm{kin}}}}\hbar\,\sin\theta,\qquad
E_{\mathrm i}=E_{\mathrm{kin}}-\hbar\omega .\qed
\]
\end{dimostrazione}

\begin{chiave}[Fotoemissione risolta in angolo]
\[
\boxed{\ k_{\mathrm i\parallel}=\frac{\sqrt{2m_0E_{\mathrm{kin}}}}\hbar\,\sin\theta,\qquad
E_{\mathrm{kin}}=\hbar\omega-\phi_{\mathrm e}-\bigl(\EF-E_{\mathrm i}\bigr)\ }
\]
(la seconda è $E_{\mathrm i}=E_{\mathrm{kin}}-\hbar\omega$ con lo zero spostato da $E_{\mathrm{vuoto}}$ a $\EF=-\phi_{\mathrm e}$.) Misurando
$E_{\mathrm{kin}}$ a vari $\theta$ si traccia $E(k_\parallel)$; con $\hbar\omega=\SI{21.2}{eV}$ bastano poche decine di gradi per attraversare la zona
(esercizio E12).
\end{chiave}

\begin{approfondimento}[impulso del fotone e componente perpendicolare]
Per $\hbar\omega=\SI{21.2}{eV}$ il fotone ha $k=\omega/c=\SI{0.011}{\angstrom^{-1}}$, lo $0{,}7\,\%$ di $2\pi/a$: per questo $\vk_{\mathrm{fin}}=
\vk_{\mathrm i}$ (a meno di un $\vG$). Per $k_\perp$ serve un modello dello stato finale: una parabola libera con il fondo a $-V_{\mathrm{int}}$
sotto il vuoto (figura (a) della slide 8; $V_{\mathrm{int}}$, potenziale interno, si stima o si misura).
\passo{Stato finale quasi libero.}
\[
E_{\mathrm{kin}}=\frac{\hbar^2\bigl(k_{\mathrm{fin}\parallel}^2+k_{\mathrm{fin}\perp}^2\bigr)}{2m_0}-V_{\mathrm{int}} .
\]
\passo{Sostituiamo $k_{\mathrm{fin}\parallel}^2=2m_0E_{\mathrm{kin}}\sin^2\theta/\hbar^2$ e risolviamo.}
\[
k_{\mathrm{fin}\perp}=\frac1\hbar\sqrt{2m_0\bigl(E_{\mathrm{kin}}\cos^2\theta+V_{\mathrm{int}}\bigr)} .
\]
Variando $\hbar\omega$ (luce di sincrotrone) si esplora anche la direzione perpendicolare.
\end{approfondimento}

\subsection{Si e GaAs: il gap assoluto}\label{sec:gapass}
Si e GaAs hanno la stessa zona di Brillouin dell'alluminio (reticolo fcc), ma bande molto più complicate: due atomi e otto elettroni di valenza per
cella (slide 10--11, fig.~\ref{fig:epm}). Le caratteristiche principali:
\begin{itemize}[leftmargin=*]
\item in fondo alle bande si riconosce la \textbf{parabola libera}: con 8 elettroni per cella liberi il fondo starebbe \SI{12.5}{eV} sotto $\EF$
(Si); il fondo della valenza calcolato è a \SI{-12.6}{eV};
\item ci sono \textbf{gap assoluti}, come in 1D;
\item il gap di GaAs in $\Gamma$ somiglia a quello che in 1D apre il secondo coefficiente di Fourier, $2|V_{2G}|$ in $k=0$ (fig.~\ref{fig:schemi}b);
\item Si: cima della valenza in $\Gamma$, minimo della conduzione lungo $\Gamma X$ vicino a $X$, a un $\vk$ diverso: \textbf{gap indiretto}
(\SI{1.12}{eV}); GaAs: entrambi in $\Gamma$, \textbf{gap diretto} (\SI{1.42}{eV}).
\end{itemize}

\begin{definizione}[gap assoluto]
Intervallo di energia in cui non c'è nessuno stato, per \emph{nessun} $\vk$ della zona (slide 12). Non basta un gap a un bordo, lungo una direzione.
\end{definizione}

\begin{figure}[tbp]
\centering
\figura{epm}
\caption{Bande di Si e GaAs con lo pseudopotenziale empirico del §\ref{sec:pseudoconseg} (fattori di forma della tabella, per GaAs anche
$V^a_3=0{,}068$, $V^a_4=0{,}066$, $V^a_{11}=\SI{0.012}{Ry}$; 113 onde piane, senza spin--orbita). In blu le 4 bande di valenza, in rosa il gap
assoluto, tratteggiato $\mu$. Il modello locale dà gap di \SI{0.8}{eV} (Si, indiretto, minimo a $0{,}85$ di $\Gamma X$) e \SI{1.9}{eV} (GaAs,
diretto): la forma delle bande è quella delle slide 10--11, i gap sono indicativi (misurati \SI{1.12}{eV} e \SI{1.42}{eV}).}\label{fig:epm}
\end{figure}

\domanda{Che cosa succede riempiendo le bande? (slide 12)}
Con il principio di Pauli le bande sotto il gap si riempiono \textbf{esattamente}; sopra non resta nessun elettrone. Il potenziale chimico cade quindi
nel gap, quasi a metà (Parte~5, §4.7: $\mu_i=\frac12\Eg+\frac34\kB T\ln(m_h/m_e)$); lo zero delle energie si mette di solito alla cima della
valenza. Questi materiali non sono metalli: nessuna banda attraversa $\mu$, e per aumentare l'energia di un elettrone bisogna dargli almeno $\Eg$.

\subsection{Metalli e isolanti: la densità degli stati}\label{sec:dos}
La differenza tra metallo e isolante si vede meglio nella densità degli stati (slide 13).

\begin{definizione}[densità degli stati]
$g(E)\,\dd E$ è il numero di stati (spin compreso) per unità di volume con energia tra $E$ ed $E+\dd E$:
\[
g(E)=\frac2{(2\pi)^3}\sum_n\oint_{E_n(\vk)=E}\frac{\dd S_E}{|\nabla_{\vk}E_n(\vk)|},
\]
con l'integrale sulla superficie a energia costante nella prima zona (slide 13, che la scrive per tutto il campione, con $2V$ al numeratore).
\end{definizione}

\needspace{12\baselineskip}
\begin{dimostrazione}[la densità degli stati da una superficie a energia costante]
\passo{Stati per unità di volume nello spazio $\vk$.}
Con le condizioni periodiche ogni $\vk$ occupa $(2\pi)^3/N\Omega$, con $N\Omega$ volume del campione (§\ref{sec:riempimento}; \QP{§3.2, p.~44}); per unità di volume e con
lo spin:
\[
\dd n=2\,\frac{\dd^3k}{(2\pi)^3}.
\]
\passo{Guscio tra le superfici $E$ ed $E+\dd E$: area per spessore.}
\[
\dd^3k=\dd S_E\,\dd k_\perp .
\]
\passo{Lo spessore viene dal gradiente, normale alla superficie a energia costante.}
\[
\dd E=|\nabla_{\vk}E|\,\dd k_\perp\qquad\Longrightarrow\qquad\dd k_\perp=\frac{\dd E}{|\nabla_{\vk}E|}.
\]
\passo{Conclusione: sommiamo su tutta la superficie e su tutte le bande.}
\[
g(E)\,\dd E=\frac2{(2\pi)^3}\sum_n\oint_{E_n(\vk)=E}\frac{\dd S_E}{|\nabla_{\vk}E_n|}\,\dd E .\qed
\]
\end{dimostrazione}

\begin{dimostrazione}[verifica: elettrone libero]
\passo{Superficie sferica di raggio $k$, gradiente costante su di essa.}
$E=\hbar^2k^2/2m_0$, $|\nabla_{\vk}E|=\hbar^2k/m_0$, $\oint\dd S_E=4\pi k^2$:
\[
g(E)=\frac2{(2\pi)^3}\,\frac{4\pi k^2m_0}{\hbar^2k}=\frac{m_0k}{\pi^2\hbar^2}.
\]
\passo{Esprimiamo $k$ con $E$.}
\[
g(E)=\frac1{2\pi^2}\Bigl(\frac{2m_0}{\hbar^2}\Bigr)^{3/2}\sqrt E,
\]
la densità della Parte~2, §3 (per unità di volume).\qed
\end{dimostrazione}

\begin{interpretazione}
$g$ è grande dove le bande sono piatte ($|\nabla_{\vk}E|$ piccolo) e va a zero ai bordi di una banda. Le forme della slide 13: $\sqrt E$ in fondo
(gas libero); una cuspide dove la superficie a energia costante tocca il bordo della zona; zero nel gap.
\end{interpretazione}

\begin{figure}[tbp]
\centering
\figura{dos_metsc}
\caption{Stessa densità degli stati (schematica) e stati occupati a $T=0$ (in blu), con un numero diverso di elettroni (slide 13--14). (a) Metallo:
$\mu$ cade dentro una banda, $g(\mu)\ne0$. (b) Semiconduttore o isolante: gli elettroni riempiono esattamente le bande fino a un gap assoluto, $\mu$
cade nel gap e $g(\mu)=0$.}\label{fig:dos}
\end{figure}

\begin{definizione}[metallo e semiconduttore, slide 14]
Un \textbf{metallo} ha una densità degli stati finita al potenziale chimico (energia di Fermi): $g(\mu)\ne0$. Un \textbf{semiconduttore} (o
isolante) deve avere un \textbf{gap assoluto} (condizione necessaria, non sufficiente) e un numero di elettroni per cella che riempie
\emph{esattamente} le bande fino al gap: $g(\mu)=0$.
\end{definizione}
Isolante e semiconduttore si distinguono poi per l'ampiezza del gap (Parte~5, §1.3).

\domanda{Basta che $\EF$ cada in un gap lungo una direzione? (slide 15)}
No. Se nell'alluminio, mettendo meno elettroni, $\EF$ cadesse nel gap in $X$, sembrerebbe un semiconduttore; ma quel gap non è assoluto: verso $W$ e $K$
altre bande attraversano la stessa energia e $g(\EF)\ne0$.

\subsection{Il riempimento delle bande: contare gli elettroni}\label{sec:riempimento}
Si può prevedere dalla struttura cristallina se un materiale è un metallo? Sì, contando bande, stati $\vk$ ed elettroni di valenza per cella
primitiva (slide 16--19).

\domanda{Perché basta una cella primitiva? (slide 16)}
Dall'equazione di Schrödinger a un elettrone (1) con la forma di Bloch (2), $\psi_{n\vk}=e^{i\vk\cdot\vr}u_{n\vk}$, si ottiene l'equazione per la
parte periodica, la stessa del §\ref{sec:kpu}:
\[
\Bigl[\frac{\hp^2}{2m_0}+\frac{\hbar\vk\cdot\hp}{m_0}+\frac{\hbar^2k^2}{2m_0}+V(\vr)\Bigr]u_{n\vk}(\vr)=E_n(\vk)\,u_{n\vk}(\vr),\qquad
u_{n\vk}(\vr+\vR)=u_{n\vk}(\vr).
\]
Per ogni $\vk$ è un problema in una sola cella primitiva, con autovalori discreti $E_1(\vk)\le E_2(\vk)\le\dots$, uno per banda. Lo stato
fondamentale degli elettroni di Bloch si costruisce come quello del gas libero (slide 17, Ashcroft--Mermin p.~141): si occupano i livelli $(n,\vk)$
più bassi, con $\vk$ in una sola cella del reticolo reciproco perché ogni livello va contato una volta. Ne escono due configurazioni: (a) alcune bande
del tutto piene e le altre vuote; (b) alcune bande piene solo in parte.

\begin{dimostrazione}[quanti stati contiene una banda (slide 19)]
\passo{Condizioni periodiche su una catena di $N$ celle (Parte~2, §2).}
\[
\psi(x+Na)=\psi(x).
\]
\passo{Applichiamo il teorema di Bloch $N$ volte.}
\[
\psi(x+Na)=e^{ikNa}\,\psi(x)\qquad\Longrightarrow\qquad e^{ikNa}=1 .
\]
\passo{I valori permessi di $k$.}
\[
k=\frac{2\pi}{aN}\,n,\qquad\Delta k=k_{n+1}-k_n=\frac{2\pi}{aN}.
\]
\passo{Contiamo i $k$ in una zona, larga $2\pi/a$.}
$(2\pi/a)/\Delta k=N$ valori: $n=1,\dots,N$ come nella slide, o $-N/2<n\le N/2$ nella prima zona. In 3D, con $N=N_1N_2N_3$ celle, la zona
(volume $(2\pi)^3/\Omega$) contiene $N$ valori di $\vk$, ciascuno di volume $(2\pi)^3/N\Omega$.
\passo{Conclusione: due stati di spin per ogni $(n,\vk)$.}
Ogni banda ha $N$ valori di $\vk$, e $\vk$ è un numero quantico anche nel senso di Pauli: contiene $2N$ elettroni, 2 per cella primitiva.\qed
\end{dimostrazione}

\begin{chiave}[Riempimento delle bande]
\[
\boxed{\ \text{una banda}=N\ \text{valori di }\vk\times2\ \text{spin}=2N\ \text{elettroni}=2\ \text{per cella primitiva}\ }
\]
Con $n_{\mathrm{val}}$ elettroni di valenza per cella primitiva:
\begin{itemize}[leftmargin=*, itemsep=1pt, topsep=1pt]
\item $n_{\mathrm{val}}$ \textbf{dispari}: almeno una banda resta piena a metà $\Rightarrow$ \textbf{metallo};
\item $n_{\mathrm{val}}$ \textbf{pari}: le bande possono essere tutte piene o vuote $\Rightarrow$ isolante o semiconduttore \emph{solo se} c'è un gap
assoluto (condizione necessaria, non sufficiente).
\end{itemize}
\end{chiave}

\begin{figure}[tbp]
\centering
\figura{riempimento1d}
\caption{Una banda di una catena di $N=18$ celle: 18 valori di $k$, ciascuno con due elettroni (slide 19). (a) Con 2 elettroni per cella tutti i
$k$ sono occupati: la banda è piena e $\EF$ cade sopra. (b) Con 1 elettrone per cella se ne occupano 9, per $|k|\le\kF=\pi/2a$: banda piena a metà,
metallo.}\label{fig:riemp}
\end{figure}

\begin{osservazione}[slide 18: stati per cella, non per atomo]
«\# of states in the band = \# primitive cells = \# number of atoms in the solid»: la seconda uguaglianza vale solo con un atomo per cella primitiva
(Al, Na). Nel Si, con due atomi per cella, una banda ha $N=N_{\mathrm{at}}/2$ valori di $\vk$: si conta per cella, non per atomo.
\end{osservazione}

\domanda{Che cosa dà il conteggio nei materiali delle slide? (slide 18, 20--21)}
\begin{center}\small
\begin{tabular}{@{}llccl@{}}
\toprule
& struttura & atomi per cella & $n_{\mathrm{val}}$ & risultato\\
\midrule
Al & fcc & 1 & 3 & una banda piena e mezza: metallo\\
Si & diamante & 2 & $4+4=8$ & 4 bande piene, gap assoluto: semiconduttore\\
GaAs & zincoblenda & 2 & $3+5=8$ & 4 bande piene, gap assoluto: semiconduttore\\
Be & hcp & 2 & $2+2=4$ & bande $2s$ e $2p$ sovrapposte: metallo\\
\bottomrule
\end{tabular}
\end{center}
Nell'alluminio la mezza banda è divisa tra seconda e terza banda (esempio del §\ref{sec:al}). In Si e GaAs le 4 bande di valenza sono quelle blu
della fig.~\ref{fig:epm}: una in fondo, di tipo $s$, e tre di tipo $p$, degeneri in $\Gamma$ (le «degenerate bands» della slide 20; §\ref{sec:so}).

\domanda{Perché il berillio, con un numero pari di elettroni, è un metallo? (slide 22--23)}
Il Be ($1s^22s^2$) ha 2 elettroni di valenza e cristallizza hcp con 2 atomi per cella: 4 elettroni per cella, abbastanza per riempire esattamente le
due bande $2s$. Ma la banda $2p$ (vuota) si sovrappone alla $2s$: in alcune regioni della zona di Brillouin la $2p$ scende sotto la cima della $2s$. Gli
elettroni si dividono tra le due bande, entrambe piene in parte, e bastano energie piccole a piacere per eccitarli: metallo. Con 2 elettroni si può
riempire una banda o riempirne a metà due: il numero pari è necessario per un isolante, non sufficiente (così anche Mg e Ca). Il litio, con un solo
elettrone $2s$, è il caso semplice: banda $2s$ piena a metà.

\begin{osservazione}[slide 20--22]
(1) Slide 20, «Si: fcc (It is not a Bravais Lattice)»: il reticolo di Bravais del Si è l'fcc; a non essere di Bravais è la struttura del diamante,
fcc con una base di due atomi (Parte~1, §10). (2) Slide 21, «As has 5 valence electrons (column IV)»: l'arsenico è nella colonna V. (3) Slide 22:
la banda più bassa, indicata come «1s», viene come la seconda dagli orbitali $2s$ (2 atomi per cella danno 2 bande $s$); il livello $1s$ del Be è di
core, circa \SI{111}{eV} sotto $\EF$, fuori dalla figura.
\end{osservazione}

% =================================================================
\section{Semimetalli: il bismuto e il grafene}\label{sec:semimetalli}
% =================================================================
Il conteggio del §\ref{sec:riempimento} dice quando un cristallo \emph{deve} essere un metallo; con un numero pari di elettroni per cella decide la
struttura a bande. Due casi limite (Lezione 6a, seconda parte; Hofmann capp.~5--6 e note integrative): il bismuto, quasi un isolante ma con pochi
elettroni e lacune, e il grafene, in cui banda piena e banda vuota si toccano in un punto. Le slide 1--12 della seconda parte ripetono le slide 12--23
della prima (§§\ref{sec:gapass}--\ref{sec:riempimento}); in questa sezione «slide» indica la seconda parte.

\begin{osservazione}[notazione]
Le slide chiamano $\mu_e$, $\mu_h$ le mobilità: qui $b_e$, $b_h$ come nella Parte~5 ($\mu$ è il potenziale chimico). Per il grafene la slide 34
chiama $\mathbf b_1$, $\mathbf b_2$ i vettori primitivi del reticolo diretto, $\mathbf d_1$, $\mathbf d_2$ quelli del reciproco, $\mathbf r_i$ i vettori
verso i primi vicini e $a_0=\SI{1.42}{\angstrom}$ la distanza C--C; qui, come nella slide 20, $\va_1$, $\va_2$ (diretto), $\vb_1$, $\vb_2$
(reciproco), $\bm\delta_i$ (primi vicini), $a_{\mathrm{CC}}$ (distanza C--C). I coefficienti degli ibridi $u_l$, $v_l$, $w_l$ diventano $\alpha_l$, $\beta_l$,
$\gamma_l$ ($u$ e $v$ sono già la parte periodica di Bloch e la velocità). La velocità degli elettroni del grafene $c_g$ (slide 30) è $v_{\mathrm F}$;
la massa dell'elettrone libero $m_e$ (slide 29--30) è $m_0$; la massa a riposo relativistica $m_0$ (slide 31--32) è $m_{\mathrm r}$.
\end{osservazione}

\subsection{Il bismuto: quasi un isolante}\label{sec:bismuto}
\domanda{Perché l'effetto Hall del bismuto è «anomalo»? (slide 13--15)}
Se ogni atomo di Bi contribuisse con i suoi 5 elettroni di valenza, il gas libero darebbe $\RH=-1/ne$ con $n=5n_{\mathrm{at}}$: circa
\SI{-4e-11}{m^3/C}. Il valore misurato (Hofmann, tab.~5.1) è più grande di diversi ordini di grandezza: i portatori sono molti meno.

\domanda{Che cosa dice il conteggio degli elettroni?}
Il Bi ($6s^26p^3$, gruppo V) ha 5 elettroni di valenza e \textbf{2 atomi per cella primitiva}: 10 elettroni per cella, numero pari. Le due bande più
basse vengono dai $6s$ (4 elettroni), ben separate dalle bande $6p$; i 6 elettroni $p$ bastano a riempire esattamente altre tre bande. In tutto
\textbf{5 bande piene}: il Bi potrebbe essere un isolante. Non lo è, ma per poco (fig.~\ref{fig:bismuto}):
\begin{itemize}
\item la cima della banda di «valenza» nel punto $T$ della zona supera $\EF$: pochissime \textbf{lacune};
\item il fondo della banda di «conduzione» nel punto $L$ scende sotto $\EF$: pochissimi \textbf{elettroni};
\item le due bande si sovrappongono in energia di circa \SI{40}{meV}; per la neutralità il numero di elettroni e di lacune è lo stesso, $n=p$;
\item la densità degli stati vicino a $\EF$ sembra nulla, come in un gap, ma è solo piccolissima: il Bi è un metallo, con portatori elettroni
\emph{e} lacune, in numero molto minore di uno per atomo ($n=p\approx\SI{3e23}{m^{-3}}$, un portatore ogni $10^5$ atomi circa).
\end{itemize}

\begin{figure}[tbp]
\centering\figura{bismuto}
\caption{Bismuto, schema delle bande che attraversano $\EF$ (slide 13 e 17, Hofmann): la «valenza» supera $\EF$ in $T$ (lacune), la «conduzione»
scende sotto $\EF$ in $L$ (elettroni); la sovrapposizione è di circa \SI{40}{meV}, e questi stati si eccitano anche termicamente.}\label{fig:bismuto}
\end{figure}

\begin{definizione}[semimetallo]
Un \textbf{semimetallo} \ingl{semimetal} ha un numero pari di elettroni per cella, che basterebbe a riempire esattamente delle bande, ma una banda
piena e una vuota si sovrappongono (Bi) o si toccano (grafene, §\ref{sec:grafenebande}) in una piccola regione della zona: $g(\EF)$ è nulla o
piccolissima e i portatori sono pochi, elettroni e lacune insieme.
\end{definizione}

\domanda{Come si legge allora il coefficiente di Hall? (slide 14--17)}
Con elettroni e lacune insieme vale la formula a due portatori (dimostrata nella Parte~5, §8.1, Kittel cap.~8; Hofmann problema 7.6):
\[
\RH=\frac{p\,b_h^2-n\,b_e^2}{e\,(p\,b_h+n\,b_e)^2}\qquad\xrightarrow{\ n=p\ }\qquad
\RH=\frac{b_h^2-b_e^2}{e\,n\,(b_h+b_e)^2}=\frac{b_h-b_e}{e\,n\,(b_h+b_e)} .
\]
\begin{itemize}
\item Il conteggio sbagliava perché attribuiva all'effetto Hall \emph{tutti} gli elettroni: contribuiscono solo i pochi portatori vicino a $\EF$.
Con $n$ piccolo il denominatore è piccolo e $\RH$ è grande.
\item Il segno dipende dalle mobilità: contributo negativo dagli elettroni in $L$, positivo dalle lacune in $T$, entrambi grandi; per $b_e=b_h$ si
annulla il numeratore.
\item Bande strette vicino a $\EF$: scale di energia piccole (decine di meV) e masse efficaci piccole; spin--orbita forte, ma ogni banda contiene
sempre due elettroni (slide 13).
\end{itemize}

\begin{esempio}[ordine di grandezza]
Con $n=p=\SI{3e23}{m^{-3}}$ e $b_e=2b_h$: $\RH=-1/(3en)=\SI{-6.9e-6}{m^3/C}$, circa $1{,}6\cdot10^5$ volte il valore del gas libero con 5 elettroni
per atomo (\SI{-4.4e-11}{m^3/C}, $n_{\mathrm{at}}=\SI{2.82e28}{m^{-3}}$).
\end{esempio}

\begin{riepilogo}[la conclusione della slide 14]
Nei metalli complessi (come il Bi) e soprattutto nei semiconduttori la densità dei portatori è molto minore di uno per atomo e il loro segno può
essere positivo. Questi effetti Hall «anomali» li spiega la teoria delle bande: pochi portatori, eccitati anche termicamente, che possono essere
elettroni (carica efficace negativa) o lacune (carica efficace positiva).
\end{riepilogo}

\begin{osservazione}[slide 13]
«Two atoms per primitive cell (HCP not bravais)»: il Bi non è hcp; ha una struttura romboedrica (tipo A7, Kittel cap.~1) con 2 atomi per cella
primitiva. Per il conteggio conta solo questo numero.
\end{osservazione}

\subsection{Il grafene: reticolo e zona di Brillouin}\label{sec:grafenereticolo}
Il grafene è un singolo strato di atomi di carbonio disposti a nido d'ape (slide 18--20, 33--34). Il nido d'ape \textbf{non è un reticolo di
Bravais}: due atomi vicini, $A$ e $B$, hanno intorni ruotati di \ang{180}. Il reticolo di Bravais è esagonale (triangolare), con una base di due atomi.

\begin{risultato}[vettori del grafene (slide 20)]
Con $a=|\va_1|=|\va_2|=\SI{2.46}{\angstrom}$ e la distanza C--C $a_{\mathrm{CC}}=a/\sqrt3=\SI{1.42}{\angstrom}$:
\[
\va_1=a\Bigl(\tfrac{\sqrt3}2,\,-\tfrac12\Bigr),\quad \va_2=a\Bigl(\tfrac{\sqrt3}2,\,\tfrac12\Bigr);\qquad
\vb_1=\frac{2\pi}a\Bigl(\tfrac1{\sqrt3},\,-1\Bigr),\quad \vb_2=\frac{2\pi}a\Bigl(\tfrac1{\sqrt3},\,1\Bigr),\quad |\vb_i|=\frac{4\pi}{\sqrt3\,a}.
\]
\end{risultato}
\begin{dimostrazione}[i $\vb_i$ sono i vettori reciproci]
\passo{Condizione che definisce il reticolo reciproco (Parte~1, §13; slide 20, eq.~1.3 con $\va_3=c\,\zh$).}
\[ \va_i\cdot\vb_j=2\pi\,\delta_{ij}. \]
\passo{Prodotti diagonali.}
\[ \va_1\cdot\vb_1=2\pi\Bigl(\tfrac{\sqrt3}2\cdot\tfrac1{\sqrt3}+\tfrac12\Bigr)=2\pi,\qquad \va_2\cdot\vb_2=2\pi\Bigl(\tfrac12+\tfrac12\Bigr)=2\pi. \]
\passo{Prodotti misti.}
\[ \va_1\cdot\vb_2=2\pi\Bigl(\tfrac12-\tfrac12\Bigr)=0,\qquad \va_2\cdot\vb_1=0.\qed \]
\end{dimostrazione}

\begin{itemize}
\item \textbf{Cella primitiva}: il rombo di lati $\va_1$, $\va_2$, con un atomo $A$ e uno $B$; equivalente la cella esagonale, che contiene
$6\times\frac13=2$ atomi (slide 19).
\item \textbf{Primi vicini}: ogni atomo $A$ ha tre vicini $B$ a distanza $a_{\mathrm{CC}}$, nelle direzioni $\bm\delta_1$, $\bm\delta_2$, $\bm\delta_3$ a
\ang{120} l'una dall'altra (slide 34, eq.~1: $a_0(0,-1)$, $a_0(\pm\frac{\sqrt3}2,\frac12)$ nel suo sistema di assi).
\item \textbf{Zona di Brillouin}: un esagono (la cella di Wigner--Seitz del reticolo reciproco esagonale). Punti di alta simmetria: $\Gamma$ al centro;
$M$ al centro di un lato, $|\Gamma M|=|\vb_i|/2=2\pi/(\sqrt3a)$; $K$ e $K'$ ai vertici, $|\Gamma K|=4\pi/(3a)$. I vertici sono di due tipi non
equivalenti, $K$ e $K'$, alternati.
\end{itemize}

\begin{figure}[tbp]
\centering\figura{grafene_reticolo}
\caption{Grafene. (a) Nido d'ape con gli atomi $A$ e $B$, i vettori primitivi $\va_1$, $\va_2$ con la cella (in verde) e i vettori $\bm\delta_i$ verso
i primi vicini. (b) Reticolo reciproco (punti neri) e prima zona di Brillouin, con i punti $\Gamma$, $M$, $K$, $K'$.}\label{fig:grafenereticolo}
\end{figure}

\begin{osservazione}[slide 18, 20 e 34]
(1) La nota della slide 18 dice che nel grafene «la cella primitiva non è propriamente definita»: lo è, è il rombo della fig.~\ref{fig:grafenereticolo}a
con due atomi non equivalenti; non di Bravais è il nido d'ape, non la cella. (2) La slide 34 chiama i vettori primitivi diretti «basis for the first
Brillouin zone» (corretto a mano in «primitive cell») e li disegna ruotati di \ang{90} rispetto alla slide 20; la fisica non cambia. (3) Slide 20, per la
grafite: con $\vb_3=(2\pi/c)\zh$ i piani della zona sono a $k_z=\pm\pi/c$ (metà di $|\vb_3|$) e il prisma è alto $2\pi/c$, non $\pm2\pi/c$ e $4\pi/c$.
\end{osservazione}

\subsection{L'ibridazione sp$^2$: legami $\sigma$ e $\pi$}\label{sec:ibridi}
Il carbonio ($1s^22s^22p^2$) ha 4 elettroni di valenza. Nel grafene l'orbitale $2s$ e due orbitali $2p$ ($2p_x$, $2p_y$) si combinano in tre
\textbf{orbitali ibridi sp$^2$}, nel piano e alla stessa energia, che formano con i tre vicini i forti \textbf{legami $\sigma$}; l'orbitale $2p_z$,
perpendicolare al piano, forma i \textbf{legami $\pi$} (slide 21--22, 33--37).

\begin{dimostrazione}[gli orbitali ibridi sp$^2$ (slide 35--37)]
\passo{Combinazione lineare per ogni ibrido (eq.~6--7 della slide 35, notazione di Dirac).}
\[ |\mathrm{sp}^2_l\rangle=\alpha_l\,|2s\rangle+\beta_l\,|2p_x\rangle+\gamma_l\,|2p_y\rangle,\qquad l=1,2,3. \]
\passo{Condizioni di ortonormalità: $|2s\rangle$, $|2p_x\rangle$, $|2p_y\rangle$ sono ortonormali, quindi i prodotti scalari si calcolano sui coefficienti.}
\[ \langle\mathrm{sp}^2_l|\mathrm{sp}^2_l\rangle=\alpha_l^2+\beta_l^2+\gamma_l^2=1,\qquad
   \langle\mathrm{sp}^2_l|\mathrm{sp}^2_{l'}\rangle=\alpha_l\alpha_{l'}+\beta_l\beta_{l'}+\gamma_l\gamma_{l'}=0\quad(l\ne l'). \]
\passo{L'orbitale $2s$ si divide in parti uguali tra i tre ibridi (eq.~10).}
\[ \alpha_1^2+\alpha_2^2+\alpha_3^2=1,\quad \alpha_1=\alpha_2=\alpha_3\quad\Longrightarrow\quad \alpha_l=\frac1{\sqrt3}. \]
\passo{Scegliamo il primo ibrido orientato lungo $\yh$: niente $2p_x$, e la normalizzazione fissa $\gamma_1$.}
\[ \beta_1=0,\qquad \gamma_1^2=1-\frac13\quad\Longrightarrow\quad\gamma_1=\sqrt{\frac23}. \]
\passo{Ortogonalità con il primo ibrido.}
\[ \frac1{\sqrt3}\cdot\frac1{\sqrt3}+0\cdot\beta_2+\sqrt{\frac23}\,\gamma_2=0\quad\Longrightarrow\quad\gamma_2=-\frac1{\sqrt6},\qquad
   \text{e allo stesso modo}\quad\gamma_3=-\frac1{\sqrt6}. \]
\passo{Normalizzazione del secondo e del terzo.}
\[ \frac13+\beta_{2,3}^2+\frac16=1\quad\Longrightarrow\quad\beta_{2,3}=\pm\frac1{\sqrt2}. \]
\passo{Ortogonalità tra il secondo e il terzo: i segni sono opposti.}
\[ \frac13+\beta_2\beta_3+\frac16=0\quad\Longrightarrow\quad\beta_2\beta_3=-\frac12,\qquad \beta_2=\frac1{\sqrt2},\ \ \beta_3=-\frac1{\sqrt2}.\qed \]
\end{dimostrazione}

\begin{chiave}[Orbitali ibridi sp$^2$ (eq.~11 della slide 37)]
\[
|\mathrm{sp}^2_1\rangle=\frac1{\sqrt3}|2s\rangle+\sqrt{\frac23}\,|2p_y\rangle,\qquad
|\mathrm{sp}^2_{2,3}\rangle=\frac1{\sqrt3}|2s\rangle\pm\frac1{\sqrt2}|2p_x\rangle-\frac1{\sqrt6}|2p_y\rangle
\]
La parte $p$ di ciascun ibrido punta a \ang{90}, \ang{-30} e \ang{210}: tre direzioni nel piano a \ang{120}, quelle dei legami con i tre vicini
(fig.~\ref{fig:ibridi}a). Ogni legame $\sigma$ ha due elettroni (spin su e giù, Pauli); resta un elettrone per atomo nel $2p_z$ (legame $\pi$),
simmetrico per rotazioni attorno a $z$.
\end{chiave}

\domanda{E gli ibridi sp e sp$^3$? (slide 37--38)}
Allo stesso modo. \textbf{sp}: $(|2s\rangle\pm|2p_x\rangle)/\sqrt2$, a \ang{180}. \textbf{sp$^3$} (diamante, Si, Ge): un $2s$ e tre $2p$ in quattro ibridi
equivalenti, ciascuno con un lobo grande e uno piccolo,
\begin{align*}
|\mathrm{sp}^3_1\rangle&=\tfrac12\bigl(|2s\rangle+|2p_x\rangle+|2p_y\rangle+|2p_z\rangle\bigr), &
|\mathrm{sp}^3_2\rangle&=\tfrac12\bigl(|2s\rangle+|2p_x\rangle-|2p_y\rangle-|2p_z\rangle\bigr),\\
|\mathrm{sp}^3_3\rangle&=\tfrac12\bigl(|2s\rangle-|2p_x\rangle-|2p_y\rangle+|2p_z\rangle\bigr), &
|\mathrm{sp}^3_4\rangle&=\tfrac12\bigl(|2s\rangle-|2p_x\rangle+|2p_y\rangle-|2p_z\rangle\bigr),
\end{align*}
diretti verso i vertici di un tetraedro: l'angolo tra due ibridi è $\arccos(-\tfrac13)=\ang{109.5}$. I quattro ibridi sono degeneri, con energia
tra quella del $2s$ e quella del $2p$, e ognuno ospita due elettroni.

\begin{figure}[tbp]
\centering\figura{ibridi}
\caption{Direzioni degli orbitali ibridi: (a) i tre sp$^2$ nel piano, con il lobo grande verso i tre legami $\sigma$; (b) i quattro sp$^3$ verso i vertici
di un tetraedro. Il colore dei lobi indica il segno della funzione d'onda.}\label{fig:ibridi}
\end{figure}

\begin{osservazione}[slide 38]
Nella figura della slide i colori dei lobi «rappresentano la carica, positiva e negativa»: rappresentano il \emph{segno} della funzione d'onda (la
densità di carica $-e|\psi|^2$ ha un solo segno). «All four sp$^3$ hybrid orbitals are delocalized» va letto «sono degeneri» (stessa energia): sono
orbitali atomici, localizzati sull'atomo.
\end{osservazione}

\subsection{Le bande del grafene: un semiconduttore a gap nullo}\label{sec:grafenebande}
\domanda{Quante bande ci sono e quante sono piene? (slide 22--25)}
Per cella: 2 atomi $\times$ 4 elettroni $=$ \textbf{8 elettroni}. Ogni atomo ha 3 orbitali sp$^2$ e un $2p_z$: 8 orbitali per cella, quindi 8 bande
(§\ref{sec:tb3d}):
\begin{itemize}
\item dai 6 sp$^2$: 3 bande $\sigma$ (leganti) e 3 bande $\sigma^*$ (antileganti, più in alto, non disegnate nella slide); 6 elettroni $\sigma$ riempiono
le 3 bande $\sigma$, due per banda;
\item dai 2 $p_z$: la banda $\pi$ (legante) e la $\pi^*$ (antilegante); i 2 elettroni $\pi$ riempiono la $\pi$.
\end{itemize}
In tutto \textbf{4 bande piene}: il grafene potrebbe essere un isolante. Ma tra la banda $\pi$ piena e la $\pi^*$ vuota non c'è gap: si toccano
esattamente ai vertici $K$ (e $K'$) della zona, e $\EF$ passa proprio lì (fig.~\ref{fig:grafenebande}).

\begin{figure}[tbp]
\centering\figura{grafene_bande}
\caption{Bande $\pi$ e $\pi^*$ del grafene a legami forti, $E=\pm t|f(\vk)|$ con $t=\SI{2.8}{eV}$ (approfondimento qui sotto), e densità degli stati: lineare e nulla in
$\EF$, con i picchi in $\pm t$ dai punti $M$. Le bande $\sigma$ (nella slide 24, da $-4$ a $-20$ eV circa) non sono disegnate.}\label{fig:grafenebande}
\end{figure}

\begin{approfondimento}[le bande $\pi$ a legami forti (Wallace; Ashcroft--Mermin cap.~10)]
Le slide mostrano le bande senza ricavarle. Si ottengono come quelle della catena biatomica (§\ref{sec:biat}) con un orbitale $p_z$ per atomo.
\passo{Funzione di Bloch con un orbitale $p_z$ su $A$ e uno su $B$, coefficienti $b_A$, $b_B$.}
\[ \psi_{\vk}=\sum_{\vR}e^{i\vk\cdot\vR}\bigl[b_A\,\varphi(\vr-\vR)+b_B\,\varphi(\vr-\vR-\bm\delta_1)\bigr]. \]
\passo{Proiezione sugli orbitali, come nel §\ref{sec:biat}: atomi uguali, $E_A=E_B=\bar E_p$, e i tre vicini $B$ di ogni $A$ danno tre fasi.}
\[ (\bar E_p-E)\,b_A-t\,f(\vk)\,b_B=0,\qquad -t\,f^*(\vk)\,b_A+(\bar E_p-E)\,b_B=0,\qquad f(\vk)=\sum_{i=1}^3e^{i\vk\cdot\bm\delta_i}. \]
\passo{Determinante nullo, con lo zero dell'energia in $\bar E_p$.}
\[ E^2=t^2|f(\vk)|^2\quad\Longrightarrow\quad E_\pm(\vk)=\pm t\,|f(\vk)|\qquad(\pi^*:\ +,\ \ \pi:\ -). \]
\passo{Valori nei punti della zona.}
In $\Gamma$ le tre fasi valgono 1: $|f|=3$, $E=\pm3t$; in $M$ $|f|=1$, $E=\pm t$; in $K$ le tre fasi sono $1$, $e^{2\pi i/3}$, $e^{-2\pi i/3}$ e la loro
somma è nulla: $E_+=E_-=0$. \textbf{Le due bande si toccano in $K$ per la simmetria tra $A$ e $B$} (come nella catena biatomica con $E_A=E_B$, senza gap).
\passo{Sviluppo vicino a $K$: $\vk=\mathbf K+\vq$ con $q\ll1/a_{\mathrm{CC}}$.}
\[ |f(\mathbf K+\vq)|\simeq\frac32\,a_{\mathrm{CC}}\,q\quad\Longrightarrow\quad E_\pm=\pm\hbar\vF\,q,\qquad \vF=\frac{3\,t\,a_{\mathrm{CC}}}{2\hbar}. \]
\passo{Numeri: con $t\simeq\SI{2.8}{eV}$.}
$\vF=\SI{9.1e5}{m/s}\approx c/330$; la banda $\pi$ va da $-3t=\SI{-8.4}{eV}$ in $\Gamma$ a $0$ in $K$, come nella slide 27.
\end{approfondimento}

\begin{chiave}[Il cono di Dirac]
Vicino a $K$ e $K'$, con $\vq=\vk-\mathbf K$:
\[ \boxed{E_\pm(\vq)=\pm\hbar\vF\,|\vq|},\qquad \vF\approx\SI{1e6}{m/s}. \]
Dispersione lineare, come per i fotoni ($\hbar\omega=\hbar ck$), con $\vF$ al posto di $c$: nello spazio $\vk$ bidimensionale le due bande formano due coni
con il vertice in $\EF$ (\textbf{cono di Dirac}). La velocità di gruppo $\vF$ è la stessa a ogni energia.
\end{chiave}

\domanda{Com'è la densità degli stati vicino a $\EF$? (slide 26, Hofmann problema 6.2)}
Lineare in $E$ e nulla in $\EF$ (fig.~\ref{fig:grafenebande}).
\begin{dimostrazione}[densità degli stati del cono di Dirac]
\passo{Stati con $|\vq|<q$ in un cono, per unità di area: un cerchio nello spazio $\vk$ 2D, $(2\pi)^2$ per stato (Parte~2, §2).}
\[ N(q)=\frac{\pi q^2}{(2\pi)^2}=\frac{q^2}{4\pi}. \]
\passo{In funzione dell'energia, $q=|E|/\hbar\vF$.}
\[ N(E)=\frac{E^2}{4\pi\hbar^2\vF^2}\quad\Longrightarrow\quad\frac{\dd N}{\dd E}=\frac{|E|}{2\pi\hbar^2\vF^2}. \]
\passo{Degenerazioni: 2 di spin e 2 di valle ($K$ e $K'$).}
\[ g(E)=4\cdot\frac{|E|}{2\pi\hbar^2\vF^2}=\frac{2|E|}{\pi\hbar^2\vF^2}\quad(\text{per unità di area}).\qed \]
\end{dimostrazione}
Il calcolo numerico con tutta la banda della fig.~\ref{fig:grafenebande} coincide con questa formula entro lo \SI{0.5}{\percent} fino a \SI{0.3}{eV}.

\begin{chiave}[Il grafene è un semimetallo]
$g(\EF)=0$, quindi secondo la definizione del §\ref{sec:dos} il grafene \emph{non è un metallo}; ma non ha nemmeno un gap tra l'ultimo stato
occupato e il primo vuoto, come un isolante o un semiconduttore. Si chiama \textbf{semimetallo} o \textbf{semiconduttore a gap nullo}
\ingl{zero band gap semiconductor}.
\end{chiave}

\begin{osservazione}[slide 24 e 27]
(1) Slide 24, «6 available bands per atom, 12 available states per unit cell»: le 6 bande $\sigma+\sigma^*$ sono per \emph{cella} (3 orbitali sp$^2$
per 2 atomi) e contengono 12 stati; allo stesso modo $\pi+\pi^*$ sono 2 bande per cella, 4 stati. (2) Slide 27, «group velocity
$\sim\SI{e7}{ms^{-1}}$, close to speed of light»: $\vF\approx\SI{1e6}{m/s}$, circa $c/300$. Gli elettroni del grafene si dicono «relativistici» per la
\emph{forma} della dispersione, non per la velocità (lo dice la slide 31).
\end{osservazione}

\subsection{La massa efficace nel grafene}\label{sec:massagrafene}
\domanda{Gli elettroni del grafene hanno massa infinita o nulla? (slide 28--32)}
Il grafene ha la conducibilità a temperatura ambiente più alta di ogni materiale noto, per due ragioni: la temperatura di Debye alta, che riduce la
diffusione da fononi, e la struttura a bande lineare vicino a $\EF$. Ma con la definizione della Parte~3, §13.2,
\[ m^*=\hbar^2\Bigl(\frac{\dd^2E}{\dd k^2}\Bigr)^{-1}, \]
una dispersione lineare dà $\dd^2E/\dd k^2=0$ e $m^*\to\infty$: sembrerebbe impossibile accelerare gli elettroni con un campo. Eppure si dicono
«senza massa». La definizione con la curvatura presuppone una banda parabolica; ne serve una che non lo presupponga.

\begin{dimostrazione}[una definizione alternativa della massa efficace]
\passo{Definiamo la massa dalla relazione tra quantità di moto e velocità (eq.~3).}
Si identifica, in senso semiclassico, la quantità di moto con $\hbar k$ (sappiamo che non è esatto: §\ref{sec:trasporto}); per il grafene $k$ si misura da $K$.
\[ p=\hbar k=\tilde m\,v_g. \]
\passo{Velocità di gruppo (eq.~4).}
\[ v_g=\frac1\hbar\frac{\dd E}{\dd k}. \]
\passo{Sostituiamo e risolviamo per $\tilde m$ (eq.~2).}
\[ \tilde m=\hbar^2k\Bigl(\frac{\dd E}{\dd k}\Bigr)^{-1}.\qed \]
\end{dimostrazione}
\begin{dimostrazione}[verifica: elettrone libero]
\passo{Banda parabolica (eq.~5--6).}
\[ E=\frac{\hbar^2k^2}{2m_0},\qquad\frac{\dd E}{\dd k}=\frac{\hbar^2k}{m_0}. \]
\passo{Sostituiamo.}
\[ \tilde m=\hbar^2k\cdot\frac{m_0}{\hbar^2k}=m_0:\ \text{per una parabola coincide con }m^*.\qed \]
\end{dimostrazione}
(In $k=0$ la formula è indeterminata, $0/0$: si prende il limite.)

\begin{chiave}[Massa efficace del grafene (eq.~7)]
Con $E=\hbar\vF k$:
\[ \boxed{\tilde m=\frac{\hbar k}{\vF}}\ \xrightarrow{\ k\to0\ }\ 0 . \]
La massa dipende da $k$ e si annulla al punto di Dirac: per questo gli elettroni del grafene si dicono \emph{senza massa} \ingl{massless}.
\end{chiave}

\domanda{Perché «senza massa» e «relativistici»? (slide 31--32)}
Per una particella relativistica di massa a riposo $m_{\mathrm r}$:
\[ E^2=(pc)^2+(m_{\mathrm r}c^2)^2=(\hbar kc)^2+(m_{\mathrm r}c^2)^2 . \]
\begin{itemize}
\item \textbf{$m_{\mathrm r}=0$} (fotone): $E=\pm\hbar kc$. È esattamente la dispersione del grafene vicino a $\EF$ con $c\to\vF$: la parte con il segno $-$ è la
banda $\pi$ (legante), quella con $+$ la $\pi^*$ (antilegante). Eccitare un elettrone dalla $\pi$ alla $\pi^*$ crea un elettrone e una lacuna, come una
particella e la sua antiparticella; tutta l'energia va in energia cinetica, perché non c'è massa a riposo da creare (fig.~\ref{fig:dirac}a).
\item \textbf{$m_{\mathrm r}>0$}: servono almeno $2m_{\mathrm r}c^2$ per creare la coppia, anche a energia cinetica nulla. È un semiconduttore normale, con
gap $\Eg=2m_{\mathrm r}c^2$ (fig.~\ref{fig:dirac}b).
\item Con la definizione $\tilde m$:
\[ \tilde m=\hbar^2k\,\frac{E}{\hbar^2kc^2}=\frac{E}{c^2}=\frac{\sqrt{(\hbar kc)^2+(m_{\mathrm r}c^2)^2}}{c^2}\qquad(\text{eq.~9}), \]
che in $k=0$ restituisce la massa a riposo $m_{\mathrm r}$ (e $0$ per $m_{\mathrm r}=0$).
\end{itemize}

\begin{figure}[H]
\centering\figura{dirac}
\caption{Dispersione relativistica (slide 31--32): (a) particella senza massa, come gli elettroni del grafene con $c\to\vF$: l'eccitazione di una coppia
elettrone (rosso)--lacuna (arancione) non ha soglia; (b) particella con massa: serve almeno $2m_{\mathrm r}c^2$, come il gap di un semiconduttore.}\label{fig:dirac}
\end{figure}

% =================================================================
\section{Pacchetti d'onde, diffusione e quasi-impulso}\label{sec:trasporto}
% =================================================================
Come si muove un elettrone di Bloch e che cosa ne limita il moto: il pacchetto d'onde, il lavoro del campo elettrico, la diffusione da difetti e
vibrazioni e il ruolo di $\hbar\vk$ negli urti (Lezione 5d, slide 49--54). Molto è già nella Parte~3 (§§8.3 e 13); qui le dimostrazioni delle slide 50--51 e il
quasi-impulso negli urti.

\subsection{Il pacchetto d'onde di Bloch}\label{sec:pacchetto}
Le onde di Bloch sono delocalizzate su tutto il solido: non sono un buon quadro per un elettrone che viaggia nel cristallo (slide 49--50, da Hofmann).
Come per le vibrazioni reticolari, l'elettrone si descrive come un \textbf{pacchetto d'onde} \ingl{wave packet}, sovrapposizione di onde di Bloch
con $k$ in un intervallo $\Delta k$ attorno al valore di interesse $k_0$, pesate da una funzione $\Phi(k)$:
\[
\psi_k(x,t)=e^{ikx-i\omega(k)t}\,u_k(x),\qquad
\Psi(x,t)=\frac1{\sqrt{2\pi}}\int_{-\infty}^{\infty}\Phi(k)\,e^{i[kx-\omega(k)t]}\,u_k(x)\,\dd k .
\]
Questa «particella» ha velocità di gruppo $v_g=\frac1\hbar\frac{\dd E}{\dd k}$ (eq.~6.47 di Hofmann) e porta a una descrizione quasi classica della
conduzione elettrica.

\begin{dimostrazione}[il pacchetto gaussiano si muove con $v_g$ senza deformarsi (slide 51)]
\passo{Peso gaussiano in $k$.}
$\Phi(k)=e^{-(k-k_0)^2/b^2}$ (costanti di normalizzazione omesse, come nella slide):
\[
\Psi(x,t)=\int_{-\infty}^{\infty}e^{-(k-k_0)^2/b^2}\,e^{ikx-i\omega(k)t}\,u_k(x)\,\dd k .
\]
\passo{$b$ piccolo rispetto alla zona: $u_k\simeq u_{k_0}$.}
Per descrivere un pacchetto in moto la gaussiana deve essere stretta, $b\ll2\pi/a$: è concentrata attorno a $k_0$, e $u_k(x)\simeq u_{k_0}(x)$ esce
dall'integrale.
\passo{Sviluppo di Taylor di $\omega(k)$ al primo ordine.}
\[
\omega(k)\simeq\omega(k_0)+v_g(k-k_0)=v_gk+\omega_0,\qquad v_g=\frac{\dd\omega}{\dd k}\Big|_{k_0}=\frac1\hbar\frac{\dd E}{\dd k}\Big|_{k_0},
\qquad\omega_0=\omega(k_0)-v_gk_0 .
\]
\passo{Poniamo $q=k-k_0$.}
\[
\Psi(x,t)\simeq u_{k_0}(x)\,e^{i[k_0x-\omega(k_0)t]}\int_{-\infty}^{\infty}e^{-q^2/b^2}\,e^{iq(x-v_gt)}\,\dd q .
\]
\passo{Integrale gaussiano: $\int e^{-q^2/b^2}e^{iqX}\,\dd q=\sqrt\pi\,b\,e^{-b^2X^2/4}$.}
\[
\Psi(x,t)\simeq\sqrt\pi\,b\;u_{k_0}(x)\,e^{i[k_0x-\omega(k_0)t]}\,e^{-b^2(x-v_gt)^2/4}.
\]
\passo{Conclusione.}
Un inviluppo gaussiano largo $1/b$, centrato in $x=v_gt$, che trasla rigidamente con velocità $v_g$; $|u_{k_0}|^2$ dà la struttura dentro ogni cella:
\[
|\Psi(x,t)|^2=\pi b^2\,|u_{k_0}(x)|^2\,e^{-b^2(x-v_gt)^2/2}.\qed
\]
\end{dimostrazione}

\begin{figure}[tbp]
\centering
\figura{pacchetto}
\caption{Pacchetto d'onde gaussiano (slide 50--51). (a) Peso $\Phi(k)$ concentrato attorno a $k_0$, di larghezza $b$. (b) Parte reale di $\Psi$ a
$t=0$ (senza la struttura di $u_{k_0}$) con il suo inviluppo, e l'inviluppo al tempo $t$, traslato di $v_gt$ senza deformarsi.}\label{fig:pacchetto}
\end{figure}

\begin{interpretazione}[descrizione quasi classica]
Il pacchetto si muove con velocità $v_g$ e non si allarga (figura~\ref{fig:pacchetto}): è la «particella» del modello semiclassico (Parte~3,
Modello~13.2), e il suo moto genera una corrente quasi classica nel solido. Poiché $b\ll2\pi/a$, la larghezza $1/b$ è molto maggiore del passo: il
pacchetto si estende su molte celle.
\end{interpretazione}
\begin{osservazione}[il termine successivo dello sviluppo]
«Non si allarga» vale al primo ordine. Con il termine $\frac12\omega''(k-k_0)^2$, dove $\omega''=\hbar/m^*$, la larghezza cresce come
$\frac1b\sqrt{1+(\hbar b^2t/2m^*)^2}$ (verificato numericamente): l'allargamento è trascurabile per $t\ll2m^*/\hbar b^2$. Per $u_k$ costante e
$m^*=m_0$ è il pacchetto gaussiano della particella libera, con la velocità del picco e l'allargamento, di \QP{§2.6, pp.~30--33}.
\end{osservazione}

\subsection{Il lavoro del campo: $\hbar\dot k=-e\mathcal E$}
Una particella con velocità $v_g$ e carica $-e$ in un campo elettrico $\mathcal E$ (slide 50, eq.~6.63 di Hofmann).
\begin{dimostrazione}[equazione del moto dal bilancio di energia]
\passo{Lavoro del campo nel tempo $\dd t$.}
\[
\dd E=-e\mathcal E\,v_g\,\dd t .
\]
\passo{Variazione dell'energia lungo la banda.}
\[
\dd E=\frac{\dd E}{\dd k}\,\dd k=\hbar v_g\,\dd k .
\]
\passo{Uguagliamo e semplifichiamo $v_g$.}
\[
\hbar\frac{\dd k}{\dd t}=-e\mathcal E .\qed
\]
\end{dimostrazione}
È l'equazione del moto del modello semiclassico (Parte~3, Modello~13.2), qui ricavata dalla conservazione dell'energia.

\subsection{Il reticolo perfetto non fa resistenza}
\begin{risultato}[elettroni di Bloch senza urti (slide 49, 50, 52)]
Le onde di Bloch sono soluzioni stazionarie: in un cristallo perfetto l'elettrone nel livello $(n,\vk)$ ha velocità media
$\mathbf v_n(\vk)=\frac1\hbar\nabla_{\vk}E_n(\vk)$ (Ashcroft--Mermin, eq.~8.51), che si conserva senza degradarsi pur interagendo con gli ioni: le onde
di Bloch si propagano liberamente. Questo toglie molti problemi del modello di Drude: per esempio non è un problema spiegare cammini liberi medi
\ingl{mean free path} lunghissimi a bassa temperatura.
\end{risultato}
La resistenza finita nasce da qualunque deviazione dal reticolo perfetto, cioè dalla rottura della simmetria di traslazione (slide 52):
\begin{itemize}[leftmargin=*]
\item difetti di punto ed estesi, vacanze, atomi sostituzionali nel reticolo;
\item vibrazioni termiche, che sulla scala dei tempi dell'elettrone sono come distorsioni statiche.
\end{itemize}
Queste cause danno il tempo di rilassamento $\tau$ \ingl{relaxation time} del modello di Drude (Parte~2).

\subsection{Il quasi-impulso negli urti}
$\vk$ è insieme vettore d'onda e (tre) numeri quantici. È forte la tentazione di leggere $\hbar\vk$ come quantità di moto dell'elettrone, ma è sbagliato
(slide 53; Parte~3, §8.3):
\[
\psi_{\vk}(\vr)=e^{i\vk\cdot\vr}u_{\vk}(\vr),\qquad -i\hbar\nabla\psi_{\vk}(\vr)=\hbar\vk\,\psi_{\vk}(\vr)-e^{i\vk\cdot\vr}\,i\hbar\nabla u_{\vk}(\vr):
\]
$\hbar\vk$ sarebbe la quantità di moto, con $E_{\mathrm{kin}}=\hbar^2k^2/2m$, solo per $u_{\vk}$ costante, cioè per l'onda di un elettrone libero.

Ma $\hbar\vk$ resta utile: è l'\textbf{impulso cristallino} \ingl{crystal momentum}, una grandezza dell'intero cristallo che si conserva negli urti a
meno di $\hbar\mathbf G$ (slide 54). Come la quantità di moto $\mathbf P$ si conserva per la simmetria di traslazione continua dello spazio, $\hbar\vk$
si conserva per quella discreta del cristallo.

\begin{risultato}[urto anelastico elettrone--fonone (slide 54)]
Un elettrone in $(\vk,E(\vk))$ emette un fonone $(\vq,\hbar\omega)$ e passa in $\vk_f$ (figura~\ref{fig:urto}):
\[
\vk=\vk_f+\vq+\mathbf G,\qquad E(\vk_f)=E(\vk)-\hbar\omega ,
\]
con $\mathbf G$ vettore del reticolo reciproco: $\mathbf G=0$ (processo normale) oppure $\mathbf G\ne0$ (processo \emph{umklapp}, che riporta $\vk_f$
nella prima zona). $\hbar\vk$ è una quantità di moto che si scambia negli urti.
\end{risultato}

\begin{figure}[tbp]
\centering
\figura{urto}
\caption{Urto anelastico di un elettrone con un fonone (slide 54): il quasi-impulso si conserva a meno di un vettore $\mathbf G$ del reticolo reciproco,
l'energia perde $\hbar\omega$.}\label{fig:urto}
\end{figure}

% =================================================================
\section{Esercizi svolti delle slide}\label{sec:eserslide}
% =================================================================
Due esercizi con la soluzione nelle slide 55--65 della Lezione 5d: un ripasso degli elettroni liberi in una scatola (Parte~2, §2) e un esercizio sugli elettroni quasi
liberi in due dimensioni (Parte~3, §11). Il terzo, sui legami forti nel cubico semplice (slide 66--68), è nel §\ref{sec:bandas}.

\subsection{Elettrone libero in una scatola rettangolare (slide 55--62)}
\textbf{Testo.} Un elettrone di massa $m$ ha energia potenziale nulla dentro una scatola rettangolare di lati $a$, $b$, $c$ lungo $Ox$, $Oy$, $Oz$, e
infinita fuori (condizioni al contorno fisse, FBC \ingl{fixed boundary conditions}). (1) Qual è l'equazione differenziale della funzione d'onda?
(2) Cercala per separazione delle variabili, $\phi(x,y,z)=\phi_x(x)\phi_y(y)\phi_z(z)$: mostra che l'equazione si scrive
\[
\frac1{\phi_x}\frac{\partial^2\phi_x}{\partial x^2}+\frac1{\phi_y}\frac{\partial^2\phi_y}{\partial y^2}+\frac1{\phi_z}\frac{\partial^2\phi_z}{\partial z^2}
=-\frac{2m}{\hbar^2}E
\]
e che basta risolvere $\phi_x''+k_x^2\phi_x=0$, $\phi_y''+k_y^2\phi_y=0$, $\phi_z''+k_z^2\phi_z=0$, con $k_x^2+k_y^2+k_z^2=2mE/\hbar^2$. (3) Integra
queste equazioni. (4) Mostra che le condizioni sulle pareti impongono $\phi_x=A\sin(n_x\pi x/a)$, $\phi_y=B\sin(n_y\pi y/b)$,
$\phi_z=C\sin(n_z\pi z/c)$ con $n_x,n_y,n_z$ interi maggiori di zero, e scrivi la funzione d'onda completa. (5) Esprimi la quantizzazione
dell'energia in funzione di $n_x,n_y,n_z$ e delle dimensioni. (6) Per una cavità cubica ($a=b=c=L$) trova (a) le funzioni d'onda dei primi tre livelli
distinti, (b) la loro energia, (c) la loro degenerazione, trascurando lo spin.

\begin{esempio}[soluzione]
\passo{Equazione di Schrödinger (domanda 1).}
$-\frac{\hbar^2}{2m}\nabla^2\psi+V\psi=E\psi$, con $V=\infty$ fuori dalla scatola (quindi $\psi=0$) e $V=0$ dentro:
\[
-\frac{\hbar^2}{2m}\Bigl(\frac{\partial^2\psi}{\partial x^2}+\frac{\partial^2\psi}{\partial y^2}+\frac{\partial^2\psi}{\partial z^2}\Bigr)=E\psi .
\]
\passo{Separazione delle variabili (domanda 2).}
Sostituendo $\phi=\phi_x\phi_y\phi_z$, e poi dividendo per $\phi$, si ha la forma $f(x)+g(y)+h(z)=\text{cost}<0$:
\[
\phi_y\phi_z\frac{\dd^2\phi_x}{\dd x^2}+\phi_x\phi_z\frac{\dd^2\phi_y}{\dd y^2}+\phi_x\phi_y\frac{\dd^2\phi_z}{\dd z^2}=-\frac{2mE}{\hbar^2}\,\phi_x\phi_y\phi_z .
\]
Le variabili sono indipendenti: l'uguaglianza vale solo se $f$, $g$, $h$ sono separatamente costanti, $-k_x^2$, $-k_y^2$, $-k_z^2$:
\[
\frac{\dd^2\phi_x}{\dd x^2}+k_x^2\phi_x=0,\quad\frac{\dd^2\phi_y}{\dd y^2}+k_y^2\phi_y=0,\quad\frac{\dd^2\phi_z}{\dd z^2}+k_z^2\phi_z=0,\qquad
\frac{\hbar^2}{2m}\bigl(k_x^2+k_y^2+k_z^2\bigr)=E .
\]
\passo{Soluzioni (domanda 3).}
$\phi_x=A\sin k_xx+D\cos k_xx$, o equivalentemente $\phi_x=A'e^{ik_xx}+D'e^{-ik_xx}$; lo stesso per $y$ e $z$.
\passo{Condizioni al contorno (domanda 4).}
$\phi_x(0)=\phi_y(0)=\phi_z(0)=0$ impone $D=0$; $\phi_x(a)=\phi_y(b)=\phi_z(c)=0$ impone $k_xa=n_x\pi$, $k_yb=n_y\pi$, $k_zc=n_z\pi$ con $n_x,n_y,n_z$
interi positivi ($n=0$ dà $\phi=0$, $n<0$ lo stesso stato). La normalizzazione $\int_0^a\!\int_0^b\!\int_0^c\phi^2\,\dd x\,\dd y\,\dd z=1$ dà
$ABC=(8/abc)^{1/2}$:
\[
\phi(x,y,z)=\Bigl(\frac8{abc}\Bigr)^{1/2}\sin\frac{n_x\pi x}a\,\sin\frac{n_y\pi y}b\,\sin\frac{n_z\pi z}c .
\]
\passo{Energia (domanda 5).}
\[
E=\frac{\hbar^2\pi^2}{2m}\Bigl(\frac{n_x^2}{a^2}+\frac{n_y^2}{b^2}+\frac{n_z^2}{c^2}\Bigr).
\]
\passo{Cavità cubica (domanda 6).}
\[
\phi(x,y,z)=\Bigl(\frac8{L^3}\Bigr)^{1/2}\sin\frac{n_x\pi x}L\,\sin\frac{n_y\pi y}L\,\sin\frac{n_z\pi z}L,\qquad
E=\frac{\hbar^2\pi^2}{2mL^2}\bigl(n_x^2+n_y^2+n_z^2\bigr).
\]
I primi tre livelli distinti hanno $n_x^2+n_y^2+n_z^2=3$, $6$, $9$:
\begin{center}\small
\begin{tabular}{@{}cccl@{}}
\toprule
livello & $E$ (unità $\hbar^2\pi^2/2mL^2$) & degenerazione & stati $(n_xn_yn_z)$\\
\midrule
1 & 3 & 1 & $(111)$\\
2 & 6 & 3 & $(211)$, $(121)$, $(112)$\\
3 & 9 & 3 & $(221)$, $(212)$, $(122)$\\
\bottomrule
\end{tabular}
\end{center}
Per esempio $\phi(211)=\sqrt{8/L^3}\,\sin(2\pi x/L)\sin(\pi y/L)\sin(\pi z/L)$; i livelli successivi hanno $n^2=11$ (3 stati) e $12$ (1 stato).
\end{esempio}
Con le condizioni fisse gli stati sono onde stazionarie con $n>0$; con quelle periodiche (Parte~2) sono onde progressive. La slide 58 scrive due
volte $\partial^2\psi/\partial x^2$ (l'ultimo termine è $\partial^2\psi/\partial z^2$); la slide 59 chiama le tre costanti $k_x^2$, $k_y^2$, $k_z^2$: con
il segno delle equazioni sono $-k_x^2$, $-k_y^2$, $-k_z^2$.

\subsection{Elettroni quasi liberi su un reticolo quadrato (slide 63--65)}
\textbf{Testo.} Reticolo quadrato bidimensionale di passo $a$. (a) Mostra che l'energia cinetica di un elettrone libero in un vertice della prima zona
di Brillouin è maggiore di quella nel punto medio di un lato di un fattore $b$, e calcola $b$. (b) Il potenziale cristallino è
$V(x,y)=-2V_0\bigl[\cos(2\pi x/a)+\cos(2\pi y/a)\bigr]$, con $V_0$ costante: trova approssimativamente il gap nel punto medio del lato della zona.
(c) Supponendo esatto il risultato di (b) e il materiale bivalente, scrivi la condizione perché il sistema sia un metallo.

\begin{figure}[H]
\centering
\figura{zona2d}
\caption{Prima zona di Brillouin del reticolo quadrato di passo $a$ (figura~1.16 dell'esercizio): $X$ è il punto medio di un lato, $M$ un vertice.}\label{fig:zona2d}
\end{figure}

\begin{esempio}[soluzione]
\passo{Rapporto delle energie cinetiche (domanda a).}
$E=p^2/2m=\hbar^2k^2/2m\propto k^2$. La prima zona è un quadrato di lato $2\pi/a$ (figura~\ref{fig:zona2d}): il punto medio del lato è $X$, a distanza
$\overline{\Gamma X}=\pi/a$, il vertice è $M$, a distanza $\overline{\Gamma M}=\sqrt2\,\overline{\Gamma X}$:
\[
b=\frac{E_M}{E_X}=\frac{\overline{\Gamma M}^2}{\overline{\Gamma X}^2}=2 .
\]
\passo{Coefficienti di Fourier e gap in $X$ (domanda b).}
\[
V(x,y)=-V_0\bigl(e^{i2\pi x/a}+e^{-i2\pi x/a}+e^{i2\pi y/a}+e^{-i2\pi y/a}\bigr):
\]
i coefficienti dei quattro $\mathbf G$ più corti, $(\pm2\pi/a,0)$ e $(0,\pm2\pi/a)$, valgono $V_{\mathbf G}=V_{-\mathbf G}=-V_0$. In $X$,
$\vk=\frac\pi a\xh=\frac12\mathbf G_1$ (condizione di Bragg), e il gap è (Parte~3, §11.4)
\[
\Delta_X=2|V_{\mathbf G_1}|=2|V_0| .
\]
\passo{Condizione per un metallo (domanda c).}
Bivalente: 2 elettroni per cella, quanti ne contiene una banda. Il solido è un metallo solo se le bande si sovrappongono, cioè se la seconda banda
comincia sotto la cima della prima:
\[
E_X+\Delta_X\le E_M\ \Longrightarrow\ \frac{\hbar^2}{2m}\Bigl(\frac\pi a\Bigr)^2+2|V_0|\le2\,\frac{\hbar^2}{2m}\Bigl(\frac\pi a\Bigr)^2
\ \Longrightarrow\ |V_0|\le\frac{\pi^2\hbar^2}{4ma^2}.
\]
\end{esempio}
\begin{osservazione}[slide 64--65: una stima più accurata]
La slide 64 scrive $E=-p^2/2m$: è $p^2/2m$. La condizione (c) mette il gap tutto sopra $E_X$ e ignora il potenziale in $M$. Con i quasi liberi il gap in
$X$ è centrato in $E_X$ ($E_X\pm|V_0|$), e in $M$ le quattro onde piane degeneri danno $E_M-2|V_0|$, $E_M$ (doppio), $E_M+2|V_0|$: serve
$E_X+|V_0|<E_M-2|V_0|$, cioè $|V_0|<E_X/3=\pi^2\hbar^2/6ma^2$ (con 121 onde piane: $0{,}35\,E_X$). La stima della slide, $E_X/2$, è dello stesso ordine
ma meno restrittiva.
\end{osservazione}

% =================================================================
\section{Riepilogo finale}\label{sec:finale}
% =================================================================
\begin{formulario}
\gruppo{Legami forti (§\ref{sec:tb3d})}
\voce[\ref{sec:tb3d}]{Modello LCAO}{H_{1e}=H_{\mathrm{at}}+\Delta V,\qquad\psi_{\vk}(\vr)=\sum_{\vR}e^{i\vk\cdot\vR}\sum_nb_n\,\varphi_n(\vr-\vR)}
\voce[\ref{sec:dimI}]{Equazione (I)}{\bigl[E(\vk)-\bar E_m\bigr]b_m+\sum_n\sum_{\vR\ne0}e^{i\vk\cdot\vR}\,t_{mn}(\vR)\,b_n=0}
\voce{\quad sovrapposizione}{t_{mn}(\vR)=-\int\varphi_m^*(\vr)\,\Delta V(\vr)\,\varphi_n(\vr-\vR)\,\dd\vr}
\voce{\quad numero di bande}{M\ \text{orbitali}\ \Rightarrow\ M\ \text{bande};\qquad n_{\mathrm{at}}\ \text{atomi per cella}\ \Rightarrow\ n_{\mathrm{at}}M}
\voce[\ref{sec:bandas}]{Catena, orbitale $s$}{E(k)=\bar E_s-2t\cos ka=\bar E_s-2t+4t\sin^2\frac{ka}2,\qquad W=4t}
\voce{\quad cubico semplice}{E(\vk)=\bar E_s-2t\,(\cos k_xa+\cos k_ya+\cos k_za),\qquad W=12t}
\voce{\quad punti della zona}{\Gamma:\ \bar E_s-6t,\qquad X:\ \bar E_s-2t,\qquad M:\ \bar E_s+2t,\qquad R:\ \bar E_s+6t}
\gruppo{Catena biatomica (§\ref{sec:biat})}
\voce[\ref{sec:biat}]{Equazione (1a)}{(E_A-E)\,b_A-t\,(1+e^{-ika})\,b_B=0}
\voce{\quad equazione (2a)}{-t\,(1+e^{ika})\,b_A+(E_B-E)\,b_B=0}
\voce{\quad bande}{E_\pm(k)=\frac{E_A+E_B}2\pm\sqrt{\Bigl(\frac{E_B-E_A}2\Bigr)^2+4t^2\cos^2\frac{ka}2}}
\voce{\quad gap in $k=\pm\pi/a$}{\Eg=E_B-E_A}
\voce{\quad distanza in $k=0$}{E_+(0)-E_-(0)=2\sqrt{\tfrac14(E_B-E_A)^2+4t^2}}
\voce{\quad per $t\ll E_B-E_A$}{E_\pm\simeq E_{B,A}\pm W\cos^2\frac{ka}2,\qquad W=\frac{4t^2}{E_B-E_A}}
\voce{\quad masse al gap}{m^*_c=-m^*_v=\frac{\hbar^2(E_B-E_A)}{2t^2a^2}}
\gruppo{Metodo $\vk\cdot\mathbf p$ (§\ref{sec:kp})}
\voce[\ref{sec:kpu}]{Equazione per $u_{n\vk}$}{\Bigl[\frac{\hp^2}{2m_0}+V+\frac{\hbar\vk\cdot\hp}{m_0}+\frac{\hbar^2k^2}{2m_0}\Bigr]u_{n\vk}=E_n(\vk)\,u_{n\vk}}
\voce[\ref{sec:kpmassa}]{Perturbazione}{H_0=\frac{\hp^2}{2m_0}+V\ \ (\text{esatta in }\Gamma),\qquad H'_{\vk}=\frac{\hbar^2k^2}{2m_0}+\frac{\hbar\vk\cdot\hp}{m_0}}
\voce{\quad massa efficace}{\frac1{m^*}=\frac1{m_0}+\frac2{m_0^2k^2}\sum_{n'\ne n}\frac{|\langle u_{n0}|\vk\cdot\hp|u_{n'0}\rangle|^2}{E_n(0)-E_{n'}(0)}}
\voce[\ref{sec:so}]{\quad due bande}{\frac{m_0}{m^*_c}\simeq1+\frac{E_P}\Eg,\qquad E_P=\frac2{m_0}\bigl|\langle u_{c0}|\hat p_x|u_{v0}\rangle\bigr|^2\approx\SI{20}{eV}}
\voce{Spin--orbita}{\lambda\,\mathbf L\cdot\mathbf S=\frac\lambda2\bigl(J^2-L^2-S^2\bigr)}
\voce{\quad livelli $p$ in $\Gamma$}{+\tfrac12\lambda\hbar^2\ \ (J=\tfrac32),\qquad-\lambda\hbar^2\ \ (J=\tfrac12),\qquad\Delta=\tfrac32\lambda\hbar^2}
\voce{\quad conduzione}{E_c=\Eg+\frac{\hbar^2k^2}{2m_0}+\frac{P^2k^2}3\Bigl(\frac2\Eg+\frac1{\Eg+\Delta}\Bigr)}
\voce{\quad lacune pesanti e leggere}{E_{hh}=-\frac{\hbar^2k^2}{2m_0},\qquad E_{lh}=-\frac{\hbar^2k^2}{2m_0}-\frac{2P^2k^2}{3\Eg}}
\voce{\quad split-off}{E_{so}=-\Delta-\frac{\hbar^2k^2}{2m_0}-\frac{P^2k^2}{3(\Eg+\Delta)}}
\voce{\quad masse}{\frac{m_0}{m_c}=1+\frac{E_P}3\Bigl(\frac2\Eg+\frac1{\Eg+\Delta}\Bigr),\qquad\frac{m_0}{m_{lh}}=1+\frac{2E_P}{3\Eg}}
\voce{}{\frac{m_0}{m_{so}}=1+\frac{E_P}{3(\Eg+\Delta)},\qquad m_{hh}=m_0,\qquad E_P=\frac{2m_0P^2}{\hbar^2}}
\gruppo{Pseudopotenziale (§\ref{sec:pseudo})}
\voce[\ref{sec:opw}]{OPW}{\chi_{\vk}=e^{i\vk\cdot\vr}+\sum_c\beta_c(\vk)\,\psi_{c\vk},\qquad\beta_c(\vk)=-\langle\psi_{c\vk}|e^{i\vk\cdot\vr}\rangle}
\voce{\quad forma di Bloch}{\chi_{\vk}=e^{i\vk\cdot\vr}\bigl[1+u^c_{\vk}(\vr)\bigr],\qquad\psi_{\vk}=\sum_{\vG}c_{\vG}\,\chi_{\vk+\vG}}
\voce[\ref{sec:pseudodim}]{Pseudo-Schrödinger}{\Bigl[\frac{\hp^2}{2m_0}+V_{\mathrm{ps}}\Bigr]\psi^{\mathrm f}_{\vk}=E(\vk)\,\psi^{\mathrm f}_{\vk},\qquad
\psi_{\vk}=\psi^{\mathrm f}_{\vk}+\sum_c\beta_c\psi_{c\vk}=(1-\hat\Pi_c)\,\psi^{\mathrm f}_{\vk}}
\voce{\quad pseudopotenziale}{V_{\mathrm{ps}}=V+V',\qquad V'=\sum_c\bigl(E-E_c\bigr)|\psi_{c\vk}\rangle\langle\psi_{c\vk}|,\qquad\langle\phi|V'|\phi\rangle\ge0}
\voce[\ref{sec:pseudomotivi}]{Ione nudo (stima)}{|V_{\vG}|\simeq\SI{13.6}{eV}\,\Bigl(\frac{a_0^3}\Omega\Bigr)\frac{8\pi Z}{(a_0G)^2}}
\voce[\ref{sec:pseudoconseg}]{Coefficienti}{V_{\vG}=\frac1{\Omega_{\mathrm{cella}}}\int_{\text{cella}}V_{\mathrm{ps}}\,e^{-i\vG\cdot\vr}\,\dd\vr}
\voce{\quad diamante, zincoblenda}{V^s_3,\ V^s_8,\ V^s_{11}\ \ (\text{composti: anche }V^a_3,\ V^a_4,\ V^a_{11})}
\voce{\quad ione schermato}{\vion(q)=-\frac{Z_{\mathrm{ion}}e^2}{\varepsilon_0\Omega q^2},\qquad v(q)=\frac{\vion(q)}{\varepsilon(q)},\qquad\varepsilon(q)\simeq\frac{k_{\mathrm{TF}}^2}{q^2}}
\voce{\quad limite $q\to0$}{k_{\mathrm{TF}}^2=\frac{e^2g(\EF)}{\varepsilon_0},\qquad g(\EF)=\frac{3n}{2\EF}\ \Longrightarrow\ v(0)=-\frac23\EF}
\voce{\quad masse in $\Gamma$}{m^*_c/m_0\simeq\Eg(\Gamma)/E_P,\qquad E_P\approx\SI{20}{eV}}
\gruppo{Bande reali e riempimento (§\ref{sec:reali})}
\voce[\ref{sec:schemi}]{Simmetria delle bande}{E_n(\vk+\vG)=E_n(\vk),\qquad\vk=\vk'+\vG\ \ (\vk'\ \text{nella prima zona})}
\voce{\quad gap in 1D}{2|V_G|\ \text{in}\ k=\pm\frac\pi a,\qquad2|V_{2G}|\ \text{in}\ k=0\ \ (\text{zona ridotta})}
\voce[\ref{sec:al}]{Zona dell'fcc}{|\Gamma L|:|\Gamma X|:|\Gamma K|:|\Gamma W|=0{,}866:1:1{,}061:1{,}118\quad(\text{unità }2\pi/a)}
\voce{\quad Al, 3 elettroni}{\kF=1{,}127\,\frac{2\pi}a>|\Gamma W|,\qquad\frac{V_{\text{sfera}}}{V_{\text{zona}}}=\frac32,\qquad\EF=\SI{11.65}{eV}}
\voce{\quad Al, 1 elettrone}{\kF'=\frac{\kF}{3^{1/3}}=0{,}78\,\frac{2\pi}a<|\Gamma L|,\qquad\EF'=\frac{\EF}{3^{2/3}}=\SI{5.6}{eV}}
\voce[\ref{sec:arpes}]{Fotoemissione}{E_{\mathrm{fin}}=E_{\mathrm i}+\hbar\omega,\quad\vk_{\mathrm{fin}}=\vk_{\mathrm i},\quad K_\parallel=k_{\mathrm{fin}\parallel}}
\voce{\quad dall'angolo}{k_{\mathrm i\parallel}=\frac{\sqrt{2m_0E_{\mathrm{kin}}}}\hbar\sin\theta,\qquad E_{\mathrm{kin}}=\hbar\omega-\phi_{\mathrm e}-\bigl(\EF-E_{\mathrm i}\bigr)}
\voce{\quad perpendicolare}{k_{\mathrm{fin}\perp}=\frac1\hbar\sqrt{2m_0\bigl(E_{\mathrm{kin}}\cos^2\theta+V_{\mathrm{int}}\bigr)}}
\voce[\ref{sec:dos}]{Densità degli stati}{g(E)=\frac2{(2\pi)^3}\sum_n\oint_{E_n(\vk)=E}\frac{\dd S_E}{|\nabla_{\vk}E_n|}\quad(\text{per unità di volume})}
\voce{\quad criterio}{\text{metallo: }g(\mu)\ne0;\qquad\text{semiconduttore: gap assoluto e }g(\mu)=0}
\voce[\ref{sec:riempimento}]{Stati di una banda}{k=\frac{2\pi}{aN}\,n,\quad\Delta k=\frac{2\pi}{aN};\qquad N\ \text{valori di }\vk\ \Rightarrow\ 2N\ \text{elettroni}}
\voce{\quad conteggio}{n_{\mathrm{val}}\ \text{dispari}\Rightarrow\text{metallo};\qquad n_{\mathrm{val}}\ \text{pari}+\text{gap assoluto}\Rightarrow\text{isolante}}
\gruppo{Semimetalli: bismuto e grafene (§\ref{sec:semimetalli})}
\voce[\ref{sec:bismuto}]{Hall con due portatori}{\RH=\frac{p\,b_h^2-n\,b_e^2}{e\,(p\,b_h+n\,b_e)^2}\ \xrightarrow{\ n=p\ }\ \frac{b_h-b_e}{e\,n\,(b_h+b_e)}}
\voce[\ref{sec:grafenereticolo}]{Grafene}{\va_{1,2}=a\bigl(\tfrac{\sqrt3}2,\mp\tfrac12\bigr),\quad\vb_{1,2}=\tfrac{2\pi}a\bigl(\tfrac1{\sqrt3},\mp1\bigr),\quad a_{\mathrm{CC}}=\frac a{\sqrt3}}
\voce{\quad zona esagonale}{|\Gamma M|=\frac{2\pi}{\sqrt3\,a},\qquad|\Gamma K|=\frac{4\pi}{3a}}
\voce[\ref{sec:ibridi}]{Ibridi sp$^2$}{|\mathrm{sp}^2_1\rangle=\tfrac1{\sqrt3}|2s\rangle+\sqrt{\tfrac23}|2p_y\rangle,\quad|\mathrm{sp}^2_{2,3}\rangle=\tfrac1{\sqrt3}|2s\rangle\pm\tfrac1{\sqrt2}|2p_x\rangle-\tfrac1{\sqrt6}|2p_y\rangle}
\voce{\quad sp$^3$}{\tfrac12\bigl(|2s\rangle\pm|2p_x\rangle\pm|2p_y\rangle\pm|2p_z\rangle\bigr)\ (\text{4 combinazioni}),\qquad\ang{109.5}}
\voce[\ref{sec:grafenebande}]{Bande $\pi$}{E_\pm(\vk)=\pm t\,|f(\vk)|,\quad f(\vk)=\sum_{i=1}^3e^{i\vk\cdot\bm\delta_i};\quad\Gamma:\pm3t,\ M:\pm t,\ K:0}
\voce{\quad cono di Dirac}{E_\pm=\pm\hbar\vF|\vq|,\qquad\vF=\frac{3ta_{\mathrm{CC}}}{2\hbar}\approx\SI{1e6}{m/s}}
\voce{\quad densità degli stati}{g(E)=\frac{2|E|}{\pi\hbar^2\vF^2}\quad(\text{per unità di area})}
\voce[\ref{sec:massagrafene}]{Massa efficace}{\tilde m=\hbar^2k\Bigl(\frac{\dd E}{\dd k}\Bigr)^{-1}:\ \ \text{grafene }\frac{\hbar k}{\vF},\ \ \text{relat. }\frac{\sqrt{(\hbar kc)^2+(m_{\mathrm r}c^2)^2}}{c^2}}
\gruppo{Trasporto (§\ref{sec:trasporto})}
\voce[\ref{sec:pacchetto}]{Pacchetto gaussiano}{\Psi(x,t)=\int e^{-(k-k_0)^2/b^2}e^{i[kx-\omega(k)t]}u_k(x)\,\dd k}
\voce{\quad al primo ordine}{|\Psi|^2=\pi b^2\,|u_{k_0}(x)|^2\,e^{-b^2(x-v_gt)^2/2},\qquad v_g=\frac1\hbar\frac{\dd E}{\dd k}}
\voce{\quad lavoro del campo}{\dd E=-e\mathcal E\,v_g\,\dd t=\hbar v_g\,\dd k\ \Longrightarrow\ \hbar\dot k=-e\mathcal E}
\voce{Quasi-impulso}{-i\hbar\nabla\psi_{\vk}=\hbar\vk\,\psi_{\vk}-e^{i\vk\cdot\vr}\,i\hbar\nabla u_{\vk}}
\voce{\quad urto con un fonone}{\vk=\vk_f+\vq+\mathbf G,\qquad E(\vk_f)=E(\vk)-\hbar\omega}
\gruppo{Esercizi delle slide (§\ref{sec:eserslide})}
\voce{Scatola}{E=\frac{\hbar^2\pi^2}{2m}\Bigl(\frac{n_x^2}{a^2}+\frac{n_y^2}{b^2}+\frac{n_z^2}{c^2}\Bigr),\qquad n_x,n_y,n_z=1,2,\dots}
\voce{\quad cubo di lato $L$}{n_x^2+n_y^2+n_z^2=3,\ 6,\ 9:\ \ \text{degenerazioni }1,\ 3,\ 3}
\voce{Quasi liberi 2D}{E_M=2E_X,\qquad\Delta_X=2|V_0|}
\voce{\quad metallo (slide)}{E_X+\Delta_X\le E_M\ \Longrightarrow\ |V_0|\le\frac{\pi^2\hbar^2}{4ma^2}}
\end{formulario}

\begin{riepilogo}[concetti fondamentali]\
\begin{enumerate}[itemsep=1pt]
\item Funzione di Bloch LCAO con $M$ orbitali $\Rightarrow$ sistema lineare $M\times M$, la (I) $\Rightarrow$ $M$ bande ($n_{\mathrm{at}}M$ con $n_{\mathrm{at}}$
atomi per cella); i $t_{mn}$ sono parametri adattati agli esperimenti.
\item Ogni livello atomico diventa una banda attorno a $\bar E$, con la dispersione data dalle sovrapposizioni con i vicini: per un orbitale $s$ un coseno
per direzione, $W=2zt$.
\item Due atomi diversi per cella $\Rightarrow$ due bande e gap $E_B-E_A$ al bordo ($A=B$ $\Rightarrow$ nessun gap); la sovrapposizione respinge i
livelli: valenza sotto $E_A$ (carattere $A$), conduzione sopra $E_B$; 2 elettroni per cella $\Rightarrow$ semiconduttore in 1D.
\item $\vk\cdot\mathbf p$: equazione per $u_{n\vk}$ con $H_0$ esatta in $\Gamma$ $\Rightarrow$ teoria delle perturbazioni $\Rightarrow$ $1/m^*$ dalle
bande vicine (secondo ordine); banda sotto $\Rightarrow$ più leggera, banda sopra $\Rightarrow$ più pesante o negativa; $m^*\propto\Eg$.
\item Spin--orbita: valenza $p$ ($L=1$) $\Rightarrow$ $J=\frac32$ (lacune pesanti e leggere, degeneri in $\Gamma$) e $J=\frac12$ (split-off, più in basso di
$\Delta=\frac32\lambda\hbar^2$); la conduzione $s$ non si separa.
\item Kane: elettroni, lacune leggere e split-off sono accoppiati dall'elemento ottico $P$ $\Rightarrow$ masse $\sim m_0\Eg/E_P$; le lacune pesanti no
$\Rightarrow$ pesanti.
\item Valenza ortogonale al core (stessa $H$ hermitiana) $\Rightarrow$ oscilla nel core $\Rightarrow$ troppe onde piane; OPW: onda piana meno le sue
componenti di core $\Rightarrow$ equazione per la sola parte liscia con $V_{\mathrm{ps}}=V+\sum_c(E-E_c)|c\rangle\langle c|$.
\item $E>E_c$ $\Rightarrow$ $V'$ repulsivo cancella in parte $V$ $\Rightarrow$ $V_{\mathrm{ps}}$ debole, pochi $V_{\vG}$ ($V_3$, $V_8$, $V_{11}$), finito in
$q=0$ ($-\frac23\EF$) $\Rightarrow$ quasi liberi con $V_{\mathrm{ps}}$; non unico (stesse energie per ogni $f(E,E_c)$).
\item Bande vere: $E_n(\vk+\vG)=E_n(\vk)$ e gap; una banda $=2N$ stati (2 per cella) $\Rightarrow$ $n_{\mathrm{val}}$ dispari $\Rightarrow$ metallo
(Al), $n_{\mathrm{val}}$ pari $\Rightarrow$ isolante solo con gap assoluto (Si, GaAs; Be: $2s$ e $2p$ sovrapposte, metallo); metallo $\Leftrightarrow
g(\mu)\ne0$.
\item Numero pari ma bande che si sovrappongono o si toccano $\Rightarrow$ semimetallo: Bi (pochi elettroni in $L$ e lacune in $T$, $\RH$ enorme, segno
deciso dalle mobilità); grafene (bande $\sigma$ e $\pi$ piene, $\pi$ e $\pi^*$ si toccano in $K$: cono di Dirac, $g(\EF)=0$ senza gap, $\tilde m=\hbar k/\vF\to0$).
\item Pacchetto gaussiano di onde di Bloch $\Rightarrow$ trasla con $v_g$ senza deformarsi (al primo ordine); il lavoro del campo dà $\hbar\dot k=-e\mathcal E$.
\item Reticolo perfetto $\Rightarrow$ nessuna resistenza; difetti, impurezze e vibrazioni termiche $\Rightarrow$ urti e tempo di rilassamento $\tau$.
\item $\hbar\vk$ non è la quantità di moto ($u_{\vk}$ non costante), ma si conserva negli urti a meno di $\hbar\mathbf G$ (processi normali e
\emph{umklapp}).
\end{enumerate}
\end{riepilogo}

\subsection*{Domande di verifica (con traccia di risposta)}
\addcontentsline{toc}{subsection}{Domande di verifica}
\begin{domande}
\item \textbf{Ricava l'equazione (I) dei legami forti.} $H_{1e}=H_{\mathrm{at}}+\Delta V$, LCAO con $M$ orbitali; proiezione su $\varphi_m$; hermiticità di
$H_{\mathrm{at}}$ $\Rightarrow E_m$; ortonormalità; sovrapposizioni tra atomi diversi trascurate; il termine $\vR=0$ sposta $E_m\to\bar E_m$; definizione di
$t_{mn}$; sistema omogeneo $M\times M$, $M$ bande.
\item \textbf{Banda $s$ in una catena e nel cubico semplice.} (I) con un orbitale; primi vicini, $t$ uguale per tutti; $\bar E_s-2t\cos ka=\bar E_s-2t
+4t\sin^2(ka/2)$, $W=4t$, $t>0$; cubico: tre coseni, valori in $\Gamma$, $X$, $M$, $R$, $W=12t$; dispersione lungo $[100]$, $[110]$, $[111]$.
\item \textbf{Ricava le bande della catena biatomica.} LCAO con $b_A$, $b_B$; proiezione su $\varphi_{An}$, quattro termini; i vicini di $A_n$
sono $B_{n-1}$ e $B_n$; (1a), (2a); determinante $2\times2$ con $|1+e^{ika}|^2=4\cos^2(ka/2)$; $E_\pm$.
\item \textbf{Gap, larghezze e carattere degli stati della catena biatomica.} $\Eg=E_B-E_A$ al bordo; distanza massima in $k=0$; sviluppo per $t$
piccolo, $W=4t^2/(E_B-E_A)$, bande fuori da $[E_A,E_B]$; valenza $A$, conduzione $B$, 2 elettroni per cella; $A=B$ senza gap; $m^*\propto E_B-E_A$.
\item \textbf{Ricava l'equazione $\vk\cdot\mathbf p$ per $u_{n\vk}$.} Forma di Bloch; $\hp\psi=e^{i\vk\cdot\vr}\hbar(\vk+\nabla/i)u$; quadrato; sviluppo con
$\vk\cdot\hp=\hp\cdot\vk$; $H_0+\hbar\vk\cdot\hp/m_0+\hbar^2k^2/2m_0$; perché $\hp\ne\hbar\vk$.
\item \textbf{Ricava e commenta la massa efficace al secondo ordine.} $H_0$ esatta in $\Gamma$; $H'_{\vk}$ piccola; primo ordine nullo in un estremo;
secondo ordine con il solo $\vk\cdot\hp$; $1/m^*$; segno dei contributi; due bande $m_0/m^*\simeq1+E_P/\Eg$; verifica con i quasi liberi al bordo.
\item \textbf{Spin--orbita alla cima della valenza.} $H_{\mathrm{so}}$ con le matrici di Pauli; forma atomica $\lambda\mathbf L\cdot\mathbf S=\frac\lambda2(J^2-L^2-S^2)$;
$L=1$: $J=\frac32$ ($+\frac12\lambda\hbar^2$, 4 stati) e $J=\frac12$ ($-\lambda\hbar^2$, 2 stati), $\Delta=\frac32\lambda\hbar^2$; lacune pesanti $m_J=\pm\frac32$,
leggere $\pm\frac12$, split-off; stati $\Phi$; $L=0$ non si separa.
\item \textbf{Scrivi e commenta le bande di Kane vicino a $\Gamma$.} $E_c$, $E_{hh}$, $E_{lh}$, $E_{so}$; parametri $P$, $\Eg$, $\Delta$; masse con $E_P$;
lacune pesanti non accoppiate; modello a due bande e ripartizione $\frac23$, $\frac13$; GaAs, InSb; bande calcolate di GaAs, Ge, Si.
\item \textbf{Perché gli elettroni quasi liberi non si applicano al potenziale vero?} Limiti di legami forti e quasi liberi; funzioni di core e di
valenza; oscillazioni nel core (energia cinetica, ortogonalità per il Teorema~2 delle QP, una sola cella per Bloch); molte onde piane; $|V_{\vG}|$
dello ione nudo di alcuni eV (Si: \SI{31}{eV}).
\item \textbf{Onde piane ortogonalizzate.} Definizione; $\beta_c=-\langle\psi_{c\vk}|e^{i\vk\cdot\vr}\rangle$ dall'ortogonalità e dall'ortonormalità del
core; forma di Bloch $e^{i\vk\cdot\vr}(1+u^c_{\vk})$; proprietà nel core e negli interstizi; sviluppo $\psi_{\vk}=\sum c_{\vG}\chi_{\vk+\vG}$.
\item \textbf{Ricava l'equazione di pseudo-Schrödinger.} $\psi=\psi^{\mathrm f}+\sum\beta_c\psi_c$ in $H\psi=E\psi$; $H\psi_c=E_c\psi_c$; definizione di
$V'$; $\beta_c=-\langle\psi_c|\psi^{\mathrm f}\rangle$; $V'=\sum(E-E_c)|\psi_c\rangle\langle\psi_c|$ positivo e non locale; stesse energie. Con il
proiettore: $\psi=(1-\hat\Pi_c)\psi^{\mathrm f}$ e $V_{\mathrm{ps}}=V-H\hat\Pi_c+E\hat\Pi_c$; non unicità con $f(E,E_c)$.
\item \textbf{Perché lo pseudopotenziale è debole e che cosa se ne fa.} Cancellazione di $V$ da parte di $V'$; soft/hard core; pochi fattori di forma
e fattore di struttura del diamante; $v(0)=-\frac23\EF$ con Thomas--Fermi; metodo empirico o autoconsistente; risposta di Yu--Cardona; bande di Si,
Ge, GaAs, ZnSe; masse in $\Gamma$ e $E_P$.
\item \textbf{Schemi di zona e bande dell'alluminio.} I due effetti del potenziale periodico, $E(\vk+\vG)=E(\vk)$ e gap dove le bande libere sono
degeneri; zona estesa, ridotta (ripiegamento con $\vk=\vk'+\vG$) e periodica; Al: fcc, 3 elettroni, banda $s$ e banda $p$ in $X$, banda che torna
verso $\Gamma$ dalla zona vicina; reticolo vuoto lungo $\Gamma XWKL\Gamma$; più bande tagliano $\EF$ $\Rightarrow$ metallo; $\kF>|\Gamma W|$; con un
elettrone sfera di Fermi dentro la zona.
\item \textbf{Come si misura $E(\vk)$: la fotoemissione.} Fotoni UV monocromatici, effetto fotoelettrico, analizzatore in energia e angolo;
$E_{\mathrm i}=E_{\mathrm{kin}}-\hbar\omega$; $\vk_{\mathrm{fin}}=\vk_{\mathrm i}$ (impulso del fotone trascurabile); alla superficie si conserva
$K_\parallel$: $k_\parallel=\sqrt{2m_0E_{\mathrm{kin}}}\sin\theta/\hbar$; $k_\perp$ con il potenziale interno.
\item \textbf{Densità degli stati e definizione di metallo.} $2\,\dd^3k/(2\pi)^3$ stati per unità di volume; guscio $\dd S_E\,\dd k_\perp$ con
$\dd E=|\nabla_{\vk}E|\dd k_\perp$; verifica sul gas libero; metallo $g(\mu)\ne0$; semiconduttore: gap assoluto ed elettroni che riempiono esattamente
le bande; un gap in una sola direzione non basta.
\item \textbf{Riempimento delle bande e conteggio degli elettroni.} Equazione per $u_{n\vk}$ in una cella; condizioni periodiche $\Rightarrow$
$k=2\pi n/aN$, $N$ valori $\Rightarrow$ $2N$ elettroni per banda; dispari $\Rightarrow$ metallo, pari $\Rightarrow$ serve il gap assoluto; Al (3), Si
($4+4$), GaAs ($3+5$), Be ($2+2$, sovrapposizione $2s$--$2p$); si conta per cella, non per atomo.
\item \textbf{Perché il bismuto è un semimetallo? Che cosa dice l'effetto Hall?} 10 elettroni per cella, 5 bande piene; sovrapposizione di
\SI{40}{meV} tra valenza in $T$ e conduzione in $L$; $n=p$ piccoli; $\RH$ a due portatori, grande per $n$ piccolo, segno dalle mobilità.
\item \textbf{Reticolo e zona di Brillouin del grafene.} Nido d'ape non di Bravais: esagonale con base $A$, $B$; $\va_i$, $\vb_i$ con $\va_i\cdot\vb_j=2\pi\delta_{ij}$;
esagono con $\Gamma$, $M$, $K$, $K'$.
\item \textbf{Ricava gli orbitali ibridi sp$^2$.} Combinazione di $2s$, $2p_x$, $2p_y$; ortonormalità; $\alpha_l=1/\sqrt3$; primo lungo $\yh$; ortogonalità
e normalizzazione per gli altri; direzioni a \ang{120}; legami $\sigma$ e $\pi$.
\item \textbf{Bande del grafene: perché è un semimetallo?} 8 elettroni e 8 bande per cella, 4 piene; legami forti con $p_z$: $E=\pm t|f|$, $f(K)=0$;
cono $\pm\hbar\vF q$; $g(E)\propto|E|$, nulla in $\EF$ senza gap.
\item \textbf{Massa efficace nel grafene.} $m^*$ con la curvatura diverge; $\tilde m=\hbar^2k/(\dd E/\dd k)$ da $p=\tilde mv_g$; $m_0$ per la parabola,
$\hbar k/\vF\to0$ per il cono; analogia con $E^2=(\hbar kc)^2+(m_{\mathrm r}c^2)^2$, gap $2m_{\mathrm r}c^2$.
\item \textbf{Il pacchetto d'onde di Bloch gaussiano.} $\Psi=\int\Phi e^{i(kx-\omega t)}u_k\,\dd k$; $b\ll2\pi/a$, $u_k\simeq u_{k_0}$; Taylor al primo ordine;
integrale gaussiano; $|\Psi|^2$ trasla con $v_g$; allargamento al secondo ordine; $\dd E=-e\mathcal Ev_g\dd t\Rightarrow\hbar\dot k=-e\mathcal E$.
\item \textbf{Resistenza e quasi-impulso.} Onde di Bloch stazionarie: nessuna resistenza nel reticolo perfetto (contro Drude); cause degli urti e $\tau$;
$-i\hbar\nabla\psi_{\vk}\ne\hbar\vk\psi_{\vk}$; impulso cristallino conservato a meno di $\hbar\mathbf G$; urto con un fonone, \emph{umklapp}.
\end{domande}

\subsection*{Esercizi d'esame}
\addcontentsline{toc}{subsection}{Esercizi d'esame}
Nello stile dell'esame, con soluzioni in fondo; $\hbar^2/2m_0=\SI{3.81}{eV.\angstrom^2}$, $\hbar^2/m_0=\SI{7.62}{eV.\angstrom^2}$.

\textbf{Risposta multipla} (una sola risposta giusta).
\begin{enumerate}[label=\textbf{R\arabic*.}, leftmargin=2.3em, itemsep=2pt]
\item Cristallo con due atomi per cella, ciascuno con un orbitale $s$ e tre $p$: le bande a legami forti sono
(a) 4; (b) 2; (c) 8.
\item Nella catena biatomica la distanza minima tra le due bande è
(a) $E_B-E_A$, al bordo di zona; (b) $2t$, in $k=0$; (c) $E_B-E_A+2W$, al bordo di zona.
\item Catena biatomica con $E_A=E_B$:
(a) gap $2t$ al bordo; (b) gap $E_A$; (c) nessun gap: è la banda della catena di passo $a/2$ ripiegata.
\item Rispetto ai livelli atomici $E_A<E_B$, le bande della catena biatomica stanno
(a) fuori: valenza sotto $E_A$, conduzione sopra $E_B$; (b) tra $E_A$ ed $E_B$; (c) entrambe sopra $E_B$.
\item Nel metodo $\vk\cdot\mathbf p$, $H_0=\hp^2/2m_0+V$ è
(a) l'Hamiltoniana di un elettrone libero; (b) l'Hamiltoniana esatta in $\Gamma$; (c) l'Hamiltoniana atomica.
\item Una banda vicina che sta \emph{sopra} la banda $n$ rende la massa efficace di $n$
(a) più piccola; (b) uguale a $m_0$; (c) più grande, o negativa (contributo negativo a $1/m^*$).
\item Nelle formule di Kane la banda non accoppiata alla conduzione è quella delle
(a) lacune pesanti; (b) lacune leggere; (c) lacune split-off.
\item La banda split-off è fatta di
(a) stati $s$, con $L=0$; (b) stati $p$ con $J=\frac12$, più in basso di $\Delta=\frac32\lambda\hbar^2$; (c) stati $p$ con $J=\frac32$.
\item Un pacchetto gaussiano di onde di Bloch, al primo ordine in $k-k_0$,
(a) trasla con $v_g$ senza deformarsi; (b) si allarga linearmente nel tempo; (c) resta fermo.
\item Un elettrone di Bloch in un cristallo perfetto e statico
(a) è diffuso dagli ioni a ogni passo reticolare; (b) ha velocità media nulla; (c) non subisce urti: la resistenza viene da difetti e vibrazioni.
\item In un urto elettrone--fonone si conserva
(a) la quantità di moto dell'elettrone; (b) il quasi-impulso totale, a meno di $\hbar\mathbf G$; (c) solo l'energia.
\item Nella regione del core la funzione d'onda di valenza
(a) è quasi nulla; (b) oscilla più delle funzioni di core, per essere ortogonale a tutte; (c) è un'onda piana.
\item Il potenziale $V'=\sum_c(E-E_c)|\psi_{c\vk}\rangle\langle\psi_{c\vk}|$ è
(a) repulsivo e non locale; (b) attrattivo e locale; (c) nullo nel core.
\item Lo pseudopotenziale di un metallo per $q\to0$ vale
(a) $-\infty$, come lo ione nudo; (b) zero; (c) $-\frac23\EF$.
\item Nel Si i fattori di forma dello pseudopotenziale che contano sono
(a) $V^s_3$, $V^s_8$, $V^s_{11}$; (b) $V^s_3$, $V^s_4$, $V^s_8$; (c) tutti i $V_{\vG}$ fino a $|\vG|^2\sim100$.
\item Sostituendo in $V'$ il fattore $E-E_c$ con una funzione arbitraria $f(E,E_c)$
(a) cambiano le bande; (b) le energie non cambiano: lo pseudopotenziale non è unico; (c) la pseudo-funzione diventa lo stato vero.
\item In zona ridotta uno stato della banda $n$ con $\vk$ fuori dalla prima zona si disegna in
(a) $\vk'=\vk-\vG$ nella prima zona, con la stessa energia; (b) $-\vk$; (c) $\vk$, nella banda $n+1$.
\item Nella zona di Brillouin dell'fcc \emph{non} sta sul bordo il punto
(a) $X$; (b) $\Gamma$; (c) $L$.
\item Nella fotoemissione, attraversando la superficie, si conserva
(a) tutto il vettore d'onda; (b) la componente parallela alla superficie; (c) la componente perpendicolare.
\item Un cristallo con 3 elettroni di valenza per cella primitiva è
(a) un metallo; (b) un isolante, se c'è un gap assoluto; (c) un semiconduttore a gap indiretto.
\item Il berillio (hcp, 4 elettroni di valenza per cella) è un metallo perché
(a) il numero di elettroni per cella è dispari; (b) le bande $2s$ e $2p$ si sovrappongono; (c) il livello $1s$ forma una banda.
\item Il gap del Si è indiretto perché
(a) la cima della valenza è in $\Gamma$ e il minimo della conduzione vicino a $X$; (b) non è assoluto; (c) la valenza è piena solo in parte.
\item Il bismuto ha 10 elettroni di valenza per cella ed è un metallo perché
(a) il numero di elettroni è dispari; (b) valenza e conduzione si sovrappongono di circa \SI{40}{meV} in punti diversi della zona; (c) tutti i 10
elettroni contribuiscono all'effetto Hall.
\item Nel grafene le bande $\pi$ e $\pi^*$ si toccano
(a) in $\Gamma$; (b) in $M$; (c) nei vertici $K$ e $K'$ della zona.
\item Vicino a $\EF$ la densità degli stati del grafene è
(a) costante, come nel gas 2D; (b) proporzionale a $|E-\EF|$ e nulla in $\EF$; (c) proporzionale a $\sqrt{E}$.
\item Con la definizione $m^*=\hbar^2/(\dd^2E/\dd k^2)$ gli elettroni del grafene vicino a $K$ hanno massa
(a) nulla; (b) infinita; (c) uguale a $m_0$.
\end{enumerate}

\textbf{Esercizi numerici.}
\begin{enumerate}[label=\textbf{E\arabic*.}, leftmargin=2.3em, itemsep=2pt]
\item \emph{(Catena biatomica)} $E_A=0$, $E_B=\SI{2}{eV}$, $t=\SI{0.5}{eV}$, $a=\SI{4}{\angstrom}$. (a) $E_\pm$ in $k=0$ e $k=\pi/a$; gap e distanza in $k=0$.
(b) Larghezza esatta di ciascuna banda e stima $W=4t^2/(E_B-E_A)$. (c) Massa efficace al gap. (d) Peso $|b_A|^2$ dello stato di valenza in $k=0$ e
$k=\pi/a$.
\item \emph{($\vk\cdot\mathbf p$)} InSb: $\Eg=\SI{0.18}{eV}$, $\Delta=\SI{0.82}{eV}$, $E_P=\SI{20}{eV}$. (a) $m_c$, $m_{lh}$, $m_{so}$ dalle formule di Kane.
(b) $m_c$ con il modello a due bande. (c) Confronta con le misure ($m_c=0{,}0145\,m_0$, $m_{lh}=0{,}021\,m_0$).
\item \emph{(Spin--orbita)} Nel Ge $\Delta=\SI{0.29}{eV}$. (a) Trova $\lambda\hbar^2$. (b) Di quanto si spostano i livelli $J=\frac32$ e $J=\frac12$ rispetto al
livello $p$ senza spin--orbita? (c) A che energia sta quel livello rispetto alla cima della valenza?
\item \emph{(Pacchetto d'onde)} Banda $E=\bar E-2t\cos ka$, $t=\SI{1}{eV}$, $a=\SI{3}{\angstrom}$, pacchetto con $b=0{,}02\,(2\pi/a)$. (a) $v_g$ per
$k_0=\pi/2a$. (b) Larghezza $1/b$ in \AA{} e in celle. (c) Per $k_0=\pi/4a$: $m^*$ e tempo di allargamento $2m^*/\hbar b^2$, da confrontare con
$\tau\sim\SI{e-14}{s}$.
\item \emph{(Quasi-impulso)} Catena di passo $a=\SI{3}{\angstrom}$, banda $\bar E-2t\cos ka$ con $t=\SI{1}{eV}$. Un elettrone in $k=0{,}9\pi/a$ emette un
fonone con $q=-0{,}2\pi/a$ e $\hbar\omega\ll t$. (a) Trova $k_f$ nella prima zona e $G$. (b) Velocità prima e dopo l'urto. (c) Perché questi urti
contano per la resistenza?
\item \emph{(Scatola)} Elettrone in una scatola cubica di lato $L=\SI{10}{\angstrom}$. (a) Energie dei primi tre livelli distinti e degenerazioni.
(b) Energia e degenerazione dei due livelli successivi.
\item \emph{(Quasi liberi 2D)} Reticolo quadrato con $a=\SI{3}{\angstrom}$ e il potenziale delle slide 63--65. (a) $E_X$ ed $E_M$ dell'elettrone libero.
(b) Il valore massimo di $|V_0|$ per avere un metallo secondo la condizione della slide e secondo la stima accurata ($E_X/3$).
\item \emph{(Ione nudo e pseudopotenziale)} Si: $a=\SI{5.431}{\angstrom}$, $Z=14$, $\Omega=a^3/8$. (a) Stima $|V_{\vG}|$ del potenziale nudo per
$|\vG|^2=3$, $8$, $11$ (unità $(2\pi/a)^2$). (b) Confronta con l'energia cinetica libera in $\vG_{111}/2$ e con i fattori di forma della tabella.
\item \emph{(Ione schermato)} Na: $n=\SI{2.65e28}{m^{-3}}$, $Z_{\mathrm{ion}}=1$. (a) $\EF$, $g(\EF)$ e $k_{\mathrm{TF}}$. (b) $v(0)$ dalla formula
$-Z_{\mathrm{ion}}e^2/\varepsilon_0\Omega k_{\mathrm{TF}}^2$ e confronto con $-\frac23\EF$.
\item \emph{(Masse in $\Gamma$)} Con $E_P=\SI{20}{eV}$ calcola $m^*_c/m_0$ per GaAs ($\Eg(\Gamma)=\SI{1.55}{eV}$) e InSb (\SI{0.24}{eV}) con
$\Eg/E_P$ e con $1/(1+E_P/\Eg)$; confronta con la tabella del §\ref{sec:pseudoconseg}.
\item \emph{(Sfera di Fermi dell'alluminio)} Al: fcc, $a=\SI{4.05}{\angstrom}$, 3 elettroni per atomo. (a) $n$, $\kF$, $\EF$ del gas libero. (b) Confronta
$\kF$ con $|\Gamma L|$, $|\Gamma X|$, $|\Gamma W|$ e calcola il rapporto tra il volume della sfera e quello della zona. (c) Ripeti con un elettrone per
atomo: la sfera tocca il bordo?
\item \emph{(Fotoemissione)} Riga He I, $\hbar\omega=\SI{21.2}{eV}$, su Al con $\phi_{\mathrm e}=\SI{4.3}{eV}$. (a) $E_{\mathrm{kin}}$ degli elettroni
che vengono da $\EF$ e da uno stato \SI{2}{eV} sotto. (b) $k_\parallel$ a $\theta=\ang{30}$ per gli elettroni da $\EF$, e l'angolo per cui
$k_\parallel=2\pi/a$. (c) Confronta l'impulso del fotone con $2\pi/a$.
\item \emph{(Conteggio)} Prevedi con il conteggio degli elettroni, e poi spiega: Na (bcc, $3s^1$), Mg (hcp, $3s^2$), C diamante ($2s^22p^2$), NaCl
(fcc con base Na e Cl), Ca (fcc, $4s^2$). Quali sono metalli per forza? Che cosa serve agli altri per essere isolanti?
\item \emph{(Grafene)} $a=\SI{2.46}{\angstrom}$, $t=\SI{2.8}{eV}$. (a) $a_{\mathrm{CC}}$, $|\vb_i|$, $|\Gamma M|$, $|\Gamma K|$. (b) $\vF$. (c) Grafene drogato
con $\EF=\SI{0.1}{eV}$ sopra il punto di Dirac: densità di elettroni per unità di area. (d) $\tilde m/m_0$ a quell'energia.
\item \emph{(Bismuto)} $n_{\mathrm{at}}=\SI{2.82e28}{m^{-3}}$, $n=p=\SI{3e23}{m^{-3}}$. (a) $\RH$ del gas libero con 5 elettroni per atomo. (b) $\RH$ con due
portatori per $b_e=2b_h$ e per $b_e=b_h$. (c) Quanti atomi per portatore?
\end{enumerate}

\textbf{Soluzioni.}
\begingroup\small
\textbf{R1} (c): $n_{\mathrm{at}}M=2\times4$. \textbf{R2} (a): $\cos(ka/2)=0$ al bordo; $\Eg+2W$ è la distanza \emph{massima}, in $k=0$. \textbf{R3} (c):
senza differenza tra gli atomi il gap si chiude. \textbf{R4} (a): la radice vale almeno $\frac12(E_B-E_A)$, i livelli si respingono. \textbf{R5} (b):
$H'_{\vk}$ si annulla per $\vk=0$. \textbf{R6} (c): denominatore $E_n-E_{n'}<0$. \textbf{R7} (a): $|X\rangle$, $|Y\rangle$ non hanno elemento di
$\hat p_z$ con $|S\rangle$. \textbf{R8} (b): $J=L-S$ con $L=1$; l'etichetta $L=0$ della slide 44 va letta $J=\frac12$. \textbf{R9} (a): l'allargamento
viene dal termine $\omega''$. \textbf{R10} (c): stati stazionari, $\mathbf v_n(\vk)$ costante. \textbf{R11} (b): $\vk=\vk_f+\vq+\mathbf G$. \textbf{R12} (b): energia cinetica
maggiore e ortogonalità (11.5). \textbf{R13} (a): $\langle\phi|V'|\phi\rangle=\sum(E-E_c)|\langle\psi_c|\phi\rangle|^2\ge0$, agisce proiettando.
\textbf{R14} (c): la schermatura toglie la divergenza. \textbf{R15} (a): $V^s_4=0$ per il fattore di struttura, gli altri sono piccoli. \textbf{R16} (b):
il termine con $f$ sparisce nel prodotto scalare con $\psi_{\vk}$, ortogonale al core. \textbf{R17} (a): $E_n(\vk+\vG)=E_n(\vk)$. \textbf{R18} (b):
$\Gamma$ è il centro della zona. \textbf{R19} (b): simmetria di traslazione nel piano; $k_\perp$ cambia alla superficie. \textbf{R20} (a): con un
numero dispari una banda resta piena a metà. \textbf{R21} (b): il numero pari è necessario ma non sufficiente. \textbf{R22} (a): massimo e minimo a
$\vk$ diversi. \textbf{R23} (b): numero pari, ma bande sovrapposte. \textbf{R24} (c): $f(\mathbf K)=0$ per la simmetria tra $A$ e $B$.
\textbf{R25} (b): cono di Dirac in 2D. \textbf{R26} (b): $\dd^2E/\dd k^2=0$; con $\tilde m=\hbar k/\vF$ invece tende a zero.

\smallskip\noindent\textbf{E1.} (a) $\sqrt{1+4\cdot0{,}25}=\SI{1.414}{eV}$: $E_\pm(0)=1\pm1{,}414$, cioè \SI{2.414}{eV} e \SI{-0.414}{eV}; in $\pi/a$:
\SI{2}{eV} e \SI{0}{eV}. Gap $\SI{2}{eV}$, distanza in $k=0$ \SI{2.83}{eV}. (b) Larghezza $1{,}414-1=\SI{0.414}{eV}$; $W=4\cdot0{,}25/2=\SI{0.5}{eV}$
(lo sviluppo richiede $4t^2\ll\frac14(E_B-E_A)^2$, qui $1$ contro $1$: stima grossolana). (c) $m^*/m_0=7{,}62\cdot2/(2\cdot0{,}25\cdot16)=1{,}9$.
(d) In $k=0$: $|b_A|^2=0{,}85$; in $\pi/a$: $|b_A|^2=1$.

\smallskip\noindent\textbf{E2.} (a) $m_0/m_c=1+\frac{20}3\bigl(\frac2{0{,}18}+\frac1{1{,}00}\bigr)=81{,}7$, $m_c=0{,}012\,m_0$;
$m_0/m_{lh}=1+\frac{40}{0{,}54}=75{,}1$, $m_{lh}=0{,}013\,m_0$; $m_0/m_{so}=1+\frac{20}{3{,}00}=7{,}67$, $m_{so}=0{,}13\,m_0$. (b) $m_c=m_0/(1+20/0{,}18)
=0{,}009\,m_0$. (c) Stesso ordine di grandezza delle misure (il modello trascura le bande lontane); la formula completa è più vicina del modello a due
bande.

\smallskip\noindent\textbf{E3.} (a) $\lambda\hbar^2=\frac23\Delta=\SI{0.193}{eV}$. (b) $J=\frac32$: $+\frac12\lambda\hbar^2=\SI{+0.097}{eV}$; $J=\frac12$:
$-\lambda\hbar^2=\SI{-0.193}{eV}$ (baricentro invariato: $4\cdot0{,}097=2\cdot0{,}193$). (c) \SI{0.097}{eV} sotto la cima della valenza.

\smallskip\noindent\textbf{E4.} (a) $v_g=\frac{2ta}\hbar\sin k_0a=\SI{9.1e5}{m/s}$ (il massimo, nel flesso). (b) $b=\SI{0.0419}{\angstrom^{-1}}$,
$1/b=\SI{23.9}{\angstrom}\simeq8$ celle. (c) $m^*=\hbar^2/(2ta^2\cos k_0a)=0{,}60\,m_0$; $2m^*/\hbar b^2=\SI{5.9e-14}{s}$, circa $6\tau$: entro il tempo
tra due urti il pacchetto resta compatto.

\smallskip\noindent\textbf{E5.} (a) $k-q=1{,}1\pi/a$ è fuori dalla zona: $k_f=1{,}1\pi/a-2\pi/a=-0{,}9\pi/a$, con $G=2\pi/a$ (\emph{umklapp}); l'energia
$E(\pm0{,}9\pi/a)$ è la stessa, compatibile con $\hbar\omega\ll t$. (b) $v=\frac{2ta}\hbar\sin ka=\SI{9.1e5}{m/s}\times\sin(0{,}9\pi)=\SI{+2.8e5}{m/s}$
prima, \SI{-2.8e5}{m/s} dopo. (c) Un fonone di piccolo $q$ inverte la velocità: questi urti degradano la corrente in modo molto efficace.

\smallskip\noindent\textbf{E6.} (a) $\frac{\hbar^2\pi^2}{2mL^2}=3{,}81\cdot\pi^2/100=\SI{0.376}{eV}$: \SI{1.13}{eV} (1 stato), \SI{2.26}{eV} (3 stati),
\SI{3.38}{eV} (3 stati). (b) $n^2=11$: \SI{4.14}{eV} (3 stati, $(311)$ e permutazioni); $n^2=12$: \SI{4.51}{eV} (1 stato, $(222)$).

\smallskip\noindent\textbf{E7.} (a) $E_X=3{,}81\,(\pi/3)^2=\SI{4.18}{eV}$, $E_M=2E_X=\SI{8.36}{eV}$. (b) Slide: $|V_0|\le E_X/2=\SI{2.09}{eV}$; stima
accurata: $|V_0|<E_X/3=\SI{1.39}{eV}$ (il calcolo con molte onde piane dà \SI{1.46}{eV}).

\smallskip\noindent\textbf{E8.} (a) $\Omega=\SI{20.0}{\angstrom^3}$, $a_0^3/\Omega=0{,}0074$; $|\vG|=\num{2.00}$, $\num{3.27}$, \SI{3.84}{\angstrom^{-1}}:
$|V_{\vG}|\simeq\SI{31}{eV}$, \SI{12}{eV}, \SI{8.6}{eV}. (b) $E_{\vG_{111}/2}=3{,}81\cdot1{,}00^2=\SI{3.8}{eV}$: il potenziale nudo è molto più grande, non
è una perturbazione; i fattori di forma sono $-2{,}9$, $0{,}5$, \SI{1.1}{eV}, da circa 8 a 20 volte più piccoli.

\smallskip\noindent\textbf{E9.} (a) $\kF=(3\pi^2n)^{1/3}=\SI{0.922}{\angstrom^{-1}}$, $\EF=3{,}81\cdot0{,}922^2=\SI{3.24}{eV}$;
$g(\EF)=3n/2\EF=\SI{0.0123}{eV^{-1}.\angstrom^{-3}}$; $k_{\mathrm{TF}}^2=4\pi\cdot14{,}4\cdot0{,}0123=\SI{2.22}{\angstrom^{-2}}$,
$k_{\mathrm{TF}}=\SI{1.49}{\angstrom^{-1}}$ (con $e^2/4\pi\varepsilon_0=\SI{14.4}{eV.\angstrom}$). (b) $\Omega=1/n=\SI{37.7}{\angstrom^3}$:
$v(0)=-4\pi\cdot14{,}4/(37{,}7\cdot2{,}22)=\SI{-2.16}{eV}=-\frac23\EF$.

\smallskip\noindent\textbf{E10.} GaAs: $1{,}55/20=0{,}078$ (tabella $0{,}078$), $1/(1+20/1{,}55)=0{,}072$; InSb: $0{,}012$ e $0{,}012$ (tabella $0{,}012$).
Misurate $0{,}067$ e $0{,}015$: la riga della tabella è $\Eg/E_P$.

\smallskip\noindent\textbf{E11.} (a) $n=12/a^3=\SI{1.81e29}{m^{-3}}$, $\kF=(3\pi^2n)^{1/3}=\SI{1.75}{\angstrom^{-1}}$, $\EF=3{,}81\cdot1{,}749^2=
\SI{11.65}{eV}$. (b) $2\pi/a=\SI{1.551}{\angstrom^{-1}}$: $|\Gamma L|=\num{1.34}$, $|\Gamma X|=\num{1.55}$, $|\Gamma W|=\SI{1.73}{\angstrom^{-1}}<\kF$: la
sfera contiene tutta la prima zona, prima banda piena; $\frac43\pi\kF^3/4(2\pi/a)^3=\frac32$. (c) $\kF'=\SI{1.21}{\angstrom^{-1}}<|\Gamma L|$ (90\,\%):
non tocca il bordo; $\EF'=\SI{5.6}{eV}$.

\smallskip\noindent\textbf{E12.} (a) $E_{\mathrm{kin}}=21{,}2-4{,}3=\SI{16.9}{eV}$; dallo stato \SI{2}{eV} sotto, \SI{14.9}{eV}. (b) $K=\sqrt{E_{\mathrm{kin}}
/3{,}81}=\SI{2.11}{\angstrom^{-1}}$; a \ang{30} $k_\parallel=\SI{1.05}{\angstrom^{-1}}=0{,}68\,(2\pi/a)$; $\sin\theta=1{,}551/2{,}107$, $\theta=\ang{47}$.
(c) $\hbar\omega/\hbar c=21{,}2/1973=\SI{0.011}{\angstrom^{-1}}$, lo $0{,}7\,\%$ di $2\pi/a$: trascurabile.

\smallskip\noindent\textbf{E13.} $n_{\mathrm{val}}$ per cella: Na 1, dispari: metallo per forza. Mg $2\cdot2=4$ e Ca 2: pari, ma metalli per la
sovrapposizione delle bande, come il Be. Diamante $2\cdot4=8$ e NaCl $1+7=8$: pari e con un gap assoluto (\SI{5.5}{eV}, circa \SI{8.5}{eV}): isolanti.

\smallskip\noindent\textbf{E14.} (a) $a_{\mathrm{CC}}=a/\sqrt3=\SI{1.42}{\angstrom}$; $|\vb_i|=4\pi/(\sqrt3a)=\SI{2.95}{\angstrom^{-1}}$;
$|\Gamma M|=\SI{1.47}{\angstrom^{-1}}$, $|\Gamma K|=4\pi/3a=\SI{1.70}{\angstrom^{-1}}$. (b) $\vF=3ta_{\mathrm{CC}}/2\hbar=\SI{9.1e5}{m/s}$
($\hbar\vF=\SI{5.96}{eV.\angstrom}$). (c) $n=\int_0^{\EF}g\,\dd E=\EF^2/(\pi\hbar^2\vF^2)=\SI{8.9e15}{m^{-2}}=\SI{8.9e11}{cm^{-2}}$. (d) $\tilde m=\hbar k/\vF=\EF/\vF^2=\SI{1.95e-32}{kg}=0{,}021\,m_0$.

\smallskip\noindent\textbf{E15.} (a) $\RH=-1/(5n_{\mathrm{at}}e)=\SI{-4.4e-11}{m^3/C}$. (b) $\RH=(b_h-b_e)/[en(b_h+b_e)]$: per $b_e=2b_h$,
$-1/(3en)=\SI{-6.9e-6}{m^3/C}$, $1{,}6\cdot10^5$ volte (a); per $b_e=b_h$, $\RH=0$. (c) $n_{\mathrm{at}}/n\approx\num{9e4}$.
\endgroup

\end{document}
